\section{The GHZ--blocked-cluster tensor}
\label{sec:mpug_ghz_blocked_cluster}

\begin{definition}[GHZ--blocked-cluster tensor and its two blocks]
    \label{def:mpug_ghz_blocked_cluster_tensor}
    \lean{MPSTensor.clusterBlocked,
        MPSTensor.ghzSectorTensor,
        MPSTensor.ghzSectorTensor_apply_same,
        MPSTensor.z4z2GHZClusterTensor,
        MPSTensor.z4z2GHZClusterBlock,
        MPSTensor.z4z2GHZClusterBlock_apply,
        MPSTensor.z4z2GHZClusterDirectSum}
    \leanok
    \discussion{7699}
    \uses{def:ghz_tensor, def:cluster_tensor,
        def:mps_independent_tensor_product}
    Let $C=\{C^q\}_{q=0}^3$ be the once-blocked cluster-state tensor. The tensor
    obtained by taking its independent tensor product with the GHZ tensor has
    physical dimension $8$ and bond dimension $4$:
    \begin{align}
        A & =A^{\mathrm{GHZ}}\boxtimes C.
        \label{eq:mpug_ghz_blocked_cluster_tensor}
    \end{align}
    For $x\in\{0,1\}$, let $E_x$ be the one-dimensional GHZ sector supported
    at the physical value $x$. The two bond-dimension-$2$ tensors
    \begin{align}
        A_x & =E_x\boxtimes C
    \end{align}
    satisfy, for $p,x\in\{0,1\}$ and $q\in\{0,1,2,3\}$,
    \begin{align}
        A_x^{(p,q)} & =\delta_{p,x}C^q.
        \label{eq:mpug_ghz_blocked_cluster_blocks}
    \end{align}
    As in \cite[lines~4427--4479]{FrancoRubio2025MPUGauging}, their canonical
    direct sum is denoted $A_0\oplus A_1$.
    % Source anchor: arXiv:2502.20257, example:z4z2, eq:z4z2MPS and
    % lines 4427--4479.
\end{definition}

\begin{lemma}[Injective and locally orthogonal GHZ--cluster blocks]
    \label{thm:mpug_ghz_blocked_cluster_blocks}
    \lean{MPSTensor.ghzSectorTensor_isInjective,
        MPSTensor.z4z2GHZClusterBlock_isInjective,
        MPSTensor.z4z2GHZClusterBlocks_isBNTLocallyOrthogonal,
        MPSTensor.z4z2GHZClusterDirectSum_eq_tensor}
    \leanok
    \discussion{7699}
    \uses{def:mpug_ghz_blocked_cluster_tensor, def:injective,
        def:is_bnt_locally_orthogonal}
    Each tensor $A_x$ is injective. For $x\neq y$ and every
    $X\in\operatorname{Mat}_2(\C)$, their one-site mixed transfer map satisfies
    \begin{align}
        \mathcal E_{x,y}(X)
         & =\sum_{p=0}^1\sum_{q=0}^3
        A_x^{(p,q)}X\bigl(A_y^{(p,q)}\bigr)^\dagger
        =0.
        \label{eq:mpug_ghz_blocked_cluster_local_orthogonality}
    \end{align}
    Thus the two blocks are locally orthogonal as observed at
    \cite[line~4481]{FrancoRubio2025MPUGauging}, and
    \begin{align}
        A_0\oplus A_1 & =A
        \label{eq:mpug_ghz_blocked_cluster_direct_sum}
    \end{align}
    after identifying the direct-sum and product bond
    coordinates.
    % The Z4 x Z2 action, its action tensors and L-symbols are treated in the
    % lemmas below; the modified fusion operators and gauging are not.
\end{lemma}

\begin{proof}\leanok
    \uses{def:mpug_ghz_blocked_cluster_tensor, def:injective,
        def:is_bnt_locally_orthogonal, thm:mps_word_span_tensor_product,
        thm:cluster_normal}
    The one-dimensional sector $E_x$ is injective, and injectivity is preserved
    by independent tensor products, so $A_x=E_x\boxtimes C$ is injective.
    If $x\neq y$, formula~(\ref{eq:mpug_ghz_blocked_cluster_blocks}) makes
    every summand in
    (\ref{eq:mpug_ghz_blocked_cluster_local_orthogonality}) vanish: a physical
    value $p$ cannot equal both $x$ and $y$. For $p,x,y\in\{0,1\}$,
    $q\in\{0,1,2,3\}$, and $a,b\in\{0,1\}$, the direct-sum and product tensors
    have the same entry in their canonical bond coordinates:
    \begin{align}
        (A_0\oplus A_1)^{(p,q)}_{(x,a),(y,b)}
         & =A^{(p,q)}_{(x,a),(y,b)}
        =\delta_{x,y}\delta_{p,x}C^q_{a,b}.
    \end{align}
    Equality of all entries proves
    (\ref{eq:mpug_ghz_blocked_cluster_direct_sum}).
\end{proof}

\begin{lemma}[GHZ--cluster physical representation and virtual action table]
    \label{thm:mpug_ghz_cluster_action_table}
    % Principal declaration: MPSTensor.z4z2GHZClusterBlock_covariance.
    \lean{MPSTensor.z4z2GHZClusterBlock_covariance,
        MPSTensor.z4z2GHZClusterAction,
        MPSTensor.z4z2GHZClusterAction_unitary,
        MPSTensor.z4z2GHZClusterAction_generators,
        MPSTensor.z4z2GHZClusterIndex,
        MPSTensor.z4z2GHZClusterBlockAction,
        MPSTensor.z4z2GHZClusterVirtual,
        MPSTensor.z4z2GHZClusterVirtual_table,
        MPSTensor.z4z2GHZClusterVirtual_adjoint,
        MPSTensor.z4z2GHZClusterVirtual_sq,
        MPSTensor.z4z2GHZClusterVirtual_stabilizer}
    \leanok
    \discussion{7747}
    \uses{def:mpug_ghz_blocked_cluster_tensor}
    As in \cite[lines~4481--4489]{FrancoRubio2025MPUGauging}, write the
    physical basis as $|p,s,r\rangle$, with integer index $4p+2s+r$ for
    $p,s,r\in\{0,1\}$. Thus the blocked cluster index is $q=2s+r$. The
    operators
    $U_{(1,0)}=(X_1\otimes X_2)\mathrm{CNOT}_{1\to2}$ and
    $U_{(0,1)}=X_3$ generate a unitary representation of
    $G=\Z_4\times\Z_2$, with $U_{(2,0)}=X_2$.
    Put $g\cdot x=x+a\pmod 2$ for $g=(a,b)$ and $x\in\{0,1\}$.
    Let $V_{g,x}$ be given by the following table, where
    $Y=\left(\begin{smallmatrix}0&-i\\i&0\end{smallmatrix}\right)$:
    \begin{center}
        \begin{tabular}{ccccccccc}
            $g$   & $(0,0)$ & $(2,0)$ & $(0,1)$ & $(2,1)$
                  & $(1,0)$ & $(3,0)$ & $(1,1)$ & $(3,1)$                             \\
            $x=0$ & $\Id$   & $X$     & $Z$     & $Y$     & $X$   & $\Id$ & $Y$ & $Z$ \\
            $x=1$ & $\Id$   & $X$     & $Z$     & $Y$     & $\Id$ & $X$   & $Z$ & $Y$ \\
        \end{tabular}
    \end{center}
    Every $V_{g,x}$ is Hermitian and squares to the identity. For every
    physical index $i$, the blocks satisfy
    \begin{align}
        \sum_j (U_g)_{i,j} A_x^j
         & =V_{g,x}A_{g\cdot x}^iV_{g,x}^{\dagger}.
        \label{eq:mpug_ghz_cluster_block_covariance}
    \end{align}
    On the subgroup $H=\{(2a,b):a,b\in\{0,1\}\}$, the comparison with the
    cluster virtual representatives $W_{(b,a)}=Z^bX^a$ is
    $V_{(2a,b),x}=(-i)^{ab}W_{(b,a)}$.
    The $L$-symbols, the nonextension of the stabilizer class and the failure
    of block independence are treated in
    Lemma~\ref{thm:mpug_ghz_cluster_l_symbols} and
    Lemma~\ref{thm:mpug_ghz_cluster_not_block_independent}.
\end{lemma}

\begin{proof}\leanok
    \uses{def:mpug_ghz_blocked_cluster_tensor, thm:mpug_monomial_matrix_laws}
    On basis vectors the generators send $(p,s,r)$ to $(p+1,s+p+1,r)$
    and $(p,s,r+1)$, with coordinates added modulo two. Their permutations
    have orders four and two and commute. The corresponding permutation
    matrices are unitary and preserve multiplication.
    Substituting $A_x^{(p,q)}=\delta_{p,x}C^q$ into
    (\ref{eq:mpug_ghz_cluster_block_covariance}) and multiplying the matrices
    for each of the eight group elements and two block labels gives the
    displayed table. Hermiticity and $V_{g,x}^2=\Id$ turn the forward action
    on the physical index into the stated conjugation identity.
    Finally, $ZX=iY$, so the only phase needed to compare the restriction
    with $Z^bX^a$ is $-i$ when $a=b=1$.
\end{proof}

\begin{lemma}[L-symbols of matrix action tensors]
    \label{thm:mpug_l_symbols_action_matrices}
    \lean{TNLean.Algebra.LSymbol.isCompatible_one_of_mul_eq_smul,
        TNLean.Algebra.LSymbol.isNormalized_of_mul_eq_smul}
    \leanok
    \uses{def:mpug_l_symbols, def:mpug_l_normalized}
    Let $G$ act on $\mathsf X$, and let $S_{g,x}$ be nonzero square matrices
    and $L^x_{g,h}$ nonzero scalars with
    \begin{align}
        S_{g,hx}S_{h,x} & =L^x_{g,h}S_{gh,x}
        \label{eq:mpug_action_matrix_l_symbols}
    \end{align}
    for all $g,h\in G$ and $x\in\mathsf X$, as for an onsite symmetry moved to
    the virtual level of each block
    \cite[lines~4363--4368]{FrancoRubio2025MPUGauging}. Then $L$ is compatible
    with the trivial three-cochain. If in addition $S_{e,x}=\Id$ for every $x$,
    then $L$ is normalized.
\end{lemma}

\begin{proof}\leanok
    \uses{def:mpug_l_symbols}
    Both bracketings of $S_{g,hkx}S_{h,kx}S_{k,x}$ are evaluated with
    (\ref{eq:mpug_action_matrix_l_symbols}). One gives
    $L^{kx}_{g,h}L^x_{gh,k}S_{ghk,x}$ and the other gives
    $L^x_{g,hk}L^x_{h,k}S_{ghk,x}$. Since $S_{ghk,x}\neq0$, the scalars agree.
    With $S_{e,x}=\Id$, the relation at $(g,e)$ and at $(e,g)$ reads
    $S_{g,x}=L^x_{g,e}S_{g,x}$ and $S_{g,x}=L^x_{e,g}S_{g,x}$.
\end{proof}

\begin{lemma}[L-symbols of the GHZ--cluster state]
    \label{thm:mpug_ghz_cluster_l_symbols}
    \lean{MPSTensor.z4z2GHZClusterBlockMulAction,
        MPSTensor.z4z2_smul_def,
        MPSTensor.z4z2ActionMatrix,
        MPSTensor.twistedTensor_z4z2GHZClusterBlock,
        MPSTensor.z4z2ActionMatrix_one,
        MPSTensor.z4z2ActionMatrix_mul_conjTranspose,
        MPSTensor.z4z2ActionMatrix_ne_zero,
        MPSTensor.z4z2PauliWord,
        MPSTensor.z4z2XCode,
        MPSTensor.z4z2PauliPhase,
        MPSTensor.z4z2ActionMatrix_eq_pauliWord,
        MPSTensor.z4z2PauliWord_mul,
        MPSTensor.z4z2XCode_mul,
        MPSTensor.z4z2IUnit,
        MPSTensor.z4z2LSymbol,
        MPSTensor.coe_z4z2LSymbol,
        MPSTensor.z4z2ActionMatrix_mul,
        MPSTensor.z4z2LSymbol_isCompatible,
        MPSTensor.z4z2LSymbol_isNormalized,
        MPSTensor.z4z2LSymbol_pow_four,
        MPSTensor.z4z2U,
        MPSTensor.z4z2V,
        MPSTensor.z4z2LSymbol_zero_sq_v}
    \leanok
    \uses{thm:mpug_ghz_cluster_action_table, thm:mpug_l_symbols_action_matrices,
        def:mpug_l_symbols, def:mpug_l_normalized}
    In the setting of Lemma~\ref{thm:mpug_ghz_cluster_action_table}, put
    $S_{g,x}=V_{g,x}$, the action tensors $A^\Gamma_{g,x}$ of
    \cite[lines~4482--4487]{FrancoRubio2025MPUGauging}. Then
    \begin{align}
        \sum_j (U_g)_{i,j}A_x^j & =S_{g,x}^\dagger A_{g\cdot x}^iS_{g,x},
        \qquad S_{e,x}=\Id,
    \end{align}
    each $S_{g,x}$ is unitary, and each is a Pauli word
    $i^{\alpha\beta}X^\alpha Z^\beta$, with $\beta=b$ for $g=(a,b)$. There are
    scalars $L^x_{g,h}$ with
    \begin{align}
        S_{g,h\cdot x}S_{h,x} & =L^x_{g,h}S_{gh,x},
        \label{eq:mpug_ghz_cluster_l_symbols}
    \end{align}
    as at \cite[line~4366]{FrancoRubio2025MPUGauging}. They are compatible with
    the trivial three-cochain, normalized, and fourth roots of unity, so they
    take values in $\{\pm1,\pm i\}$ as stated at
    \cite[line~4495]{FrancoRubio2025MPUGauging}. With $u=(1,0)$ and $v=(0,1)$,
    \begin{align}
        L^0_{u^2,v} & =-i, & L^0_{v,u^2} & =i.
        \label{eq:mpug_ghz_cluster_l_values}
    \end{align}
\end{lemma}

\begin{proof}\leanok
    \uses{thm:mpug_ghz_cluster_action_table, thm:mpug_l_symbols_action_matrices}
    The conjugation identity is the covariance of
    Lemma~\ref{thm:mpug_ghz_cluster_action_table}, since every $V_{g,x}$ is
    Hermitian. In the table, $\Id$, $X$, $Z$ and $Y=iXZ$ are the Pauli words
    with exponents $(0,0)$, $(1,0)$, $(0,1)$ and $(1,1)$. Using $X^2=Z^2=\Id$
    and $ZX=-XZ$,
    \begin{align}
        \bigl(i^{\alpha\beta}X^\alpha Z^\beta\bigr)
        \bigl(i^{\alpha'\beta'}X^{\alpha'}Z^{\beta'}\bigr)
         & =i^{\alpha\beta+\alpha'\beta'+2\beta\alpha'
                -(\alpha+\alpha')(\beta+\beta')}\,
        i^{(\alpha+\alpha')(\beta+\beta')}X^{\alpha+\alpha'}Z^{\beta+\beta'},
        \notag
    \end{align}
    with exponents of $X$ and $Z$ added modulo two. The $X$ exponent of
    $S_{g,x}$ is additive along the action,
    $\alpha(gh,x)=\alpha(g,h\cdot x)+\alpha(h,x)$, as one checks on the
    sixteen pairs, and the $Z$ exponent is additive. This gives
    (\ref{eq:mpug_ghz_cluster_l_symbols}) with $L^x_{g,h}$ the displayed power
    of $i$. Lemma~\ref{thm:mpug_l_symbols_action_matrices} gives compatibility
    and normalization. Finally $S_{u^2,0}=X$ and $S_{v,0}=Z$, and
    $XZ=-iY$ while $ZX=iY$.
\end{proof}

\begin{lemma}[Block independence fails for the GHZ--cluster state]
    \label{thm:mpug_ghz_cluster_not_block_independent}
    \lean{MPSTensor.z4z2Unbroken,
        MPSTensor.z4z2StabilizerCocycle,
        MPSTensor.z4z2StabilizerCocycle_isMulCocycle₂,
        MPSTensor.z4z2StabilizerCocycle_H2π_ne_zero,
        MPSTensor.commPhase_z4z2U_sq_z4z2V,
        MPSTensor.z4z2StabilizerCocycle_not_extends,
        MPSTensor.z4z2StabilizerCocycle_H2π_not_mem_range,
        MPSTensor.z4z2LSymbol_not_isBlockIndependent}
    \leanok
    \uses{thm:mpug_ghz_cluster_l_symbols, lem:mpug_bi_basepoint_extension,
        thm:comm_phase_bicharacter, thm:nontrivial_of_comm_phase,
        thm:mpug_h2_exact_extension, def:mpug_block_independence_gauge_forms}
    Let $H=\{(2a,b)\}$ be the stabilizer of the block $0$, the unbroken
    subgroup of \cite[line~4489]{FrancoRubio2025MPUGauging}, and let
    $\psi(a,b)=L^0_{a,b}$ for $a,b\in H$. Then $\psi$ is a two-cocycle with
    nonzero class in $H^2(H,\C^\times)$. No two-cocycle of
    $G=\Z_4\times\Z_2$ restricts to $\psi$, so $[\psi]$ is not in the image of
    the restriction map $H^2(G,\C^\times)\to H^2(H,\C^\times)$
    \cite[line~4385]{FrancoRubio2025MPUGauging}. Consequently the $L$-symbols
    of Lemma~\ref{thm:mpug_ghz_cluster_l_symbols} are not block independent
    after any joint scalar gauge, as asserted at
    \cite[line~4490]{FrancoRubio2025MPUGauging}. Block independence is meant
    here as the scalar condition on these $L$-symbols, with coefficients in
    $\C^\times$; see \cite{gap:fbc25_mpu_action_stabilizer_source_audit}. The
    modified fusion operators and the state-level gauging are not treated.
\end{lemma}

\begin{proof}\leanok
    \uses{thm:mpug_ghz_cluster_l_symbols, lem:mpug_bi_basepoint_extension,
        thm:comm_phase_bicharacter, thm:nontrivial_of_comm_phase,
        thm:mpug_h2_exact_extension}
    Compatibility of $L$ with $\omega=1$ at the block $0$ and arguments in $H$
    is the cocycle equation of $\psi$, because $H$ fixes the block $0$. By
    (\ref{eq:mpug_ghz_cluster_l_values}), the commutator phase of $\psi$ at the
    commuting pair $(u^2,v)$ is $(-i)/i=-1$, so the class of $\psi$ is
    nonzero by Lemma~\ref{thm:nontrivial_of_comm_phase}. If a two-cocycle
    $\Psi$ of $G$ restricted to $\psi$, then, $G$ being commutative,
    Lemma~\ref{thm:comm_phase_bicharacter} would give
    \begin{align}
        \beta_\Psi(u^2,v) & =\beta_\Psi(u,v)^2=\beta_\Psi(u,v^2)
        =\beta_\Psi(u,e)=1,
        \notag
    \end{align}
    against $\beta_\Psi(u^2,v)=\beta_\psi(u^2,v)=-1$. The class statement
    follows by Lemma~\ref{thm:mpug_h2_exact_extension}. Finally, block
    independence would produce such a $\Psi$ by
    Lemma~\ref{lem:mpug_bi_basepoint_extension}.
\end{proof}

\section{The controlled-\texorpdfstring{$Z$}{Z} symmetry and its gauging}
\label{sec:mpug_controlled_z}

The simplest symmetry that is not onsite is the ring of controlled-$Z$ gates
$U=\prod_j\mathrm{CZ}_{j,j+1}$ on a chain of qubits
\cite[lines~3794--3968]{FrancoRubio2025MPUGauging}. Write
$s(a,b)=(-1)^{ab}$ for bits $a,b$; these are the entries of $\sqrt2H$, and
$s(a,b)$ is the phase of $\mathrm{CZ}\ket{a,b}$.

\begin{definition}[The controlled-\texorpdfstring{$Z$}{Z} matrix product unitary]
    \label{def:mpug_controlled_z_tensor}
    \lean{ControlledZGauging.czSign,
        ControlledZGauging.czMPUTensor,
        ControlledZGauging.czMPUTensor_apply,
        ControlledZGauging.czPhase,
        ControlledZGauging.czBondDim,
        ControlledZGauging.czRepTensor,
        ControlledZGauging.czFamily}
    \leanok
    \uses{def:mpo_tensor}
    Following \cite[eq.~(CZ) and eq.~(deltas)]{FrancoRubio2025MPUGauging}, the
    gate is a copy tensor $\delta_{ijr}$ on each qubit joined along the bond by
    $\sqrt2H$. Absorbing each $\sqrt2H$ into the site on its right gives the
    bond-two tensor
    \begin{align}
        M^{ij}_{lr} & =\delta_{ij}\,\delta_{rj}\,s(l,j),
        \label{eq:mpug_controlled_z_tensor}
    \end{align}
    with output index $i$ and input index $j$. Together with the identity
    tensor of bond dimension one, it gives the family $\{\Id,M\}$ over $\Z_2$.
\end{definition}

\begin{lemma}[Tensors with a forced outgoing bond]
    \label{thm:mps_forced_column_normality}
    \lean{MPSTensor.isNBlkInjective_two_of_forced_column,
        MPSTensor.isNormal_of_forced_column,
        MPSTensor.mpv_of_forced_right_bond}
    \leanok
    \uses{def:mps_tensor}
    Let $A$ be a tensor of bond dimension $D$. Suppose that for every bond
    label $j$ some letter $\ell_j$ has the single nonzero column $j$, equal to
    a vector $v_j$ with no vanishing entry, and that $v_1,\dots,v_D$ are the
    columns of an invertible matrix. Then the words of length two span all
    $D\times D$ matrices, so $A$ is normal. If every entry $A^i_{lr}$ vanishes
    unless $r=\beta(i)$, and then equals $\varphi(i,l)$, the periodic amplitude
    on a ring is $\prod_n\varphi(\sigma_{n+1},\beta(\sigma_n))$.
\end{lemma}

\begin{proof}\leanok
    \uses{def:mps_tensor}
    The product $A^{\ell_j}A^{\ell_k}$ is $v_k(j)\,v_je_k^{T}$, so the matrix
    unit $e_le_k^T$ is the combination
    $\sum_j P_{jl}\,v_k(j)^{-1}A^{\ell_j}A^{\ell_k}$, where $P$ is the inverse
    of the matrix with columns $v_j$. For the amplitude, only the closed bond
    configuration with bond $\beta(\sigma_n)$ after site $n$ contributes.
\end{proof}

\begin{lemma}[The controlled-\texorpdfstring{$Z$}{Z} symmetry is not anomalous]
    \label{thm:mpug_controlled_z_not_anomalous}
    \lean{ControlledZGauging.mpo_czMPUTensor,
        ControlledZGauging.mulVec_mpo_czMPUTensor,
        ControlledZGauging.mpo_czMPUTensor_mul_self,
        ControlledZGauging.czSign_mul_self,
        ControlledZGauging.czSign_ne_zero,
        ControlledZGauging.toMPSTensor_czMPUTensor,
        ControlledZGauging.czMPUTensor_isNormal,
        ControlledZGauging.czFamily_isNormalRepresentation,
        ControlledZGauging.zeroState,
        ControlledZGauging.zeroState_isNormal,
        ControlledZGauging.mpv_zeroState,
        ControlledZGauging.fixesMPV_zeroState,
        ControlledZGauging.isTrivialGaugeClass_omega_czFamily,
        ControlledZGauging.czFamily_anomalyClass_eq_zero}
    \leanok
    \uses{def:mpug_controlled_z_tensor, thm:mpug_anomaly_obstruction,
        def:mpug_anomaly_class, thm:mps_forced_column_normality}
    On every ring of $N\geq1$ qubits the periodic operator of $M$ is the
    diagonal matrix with entries $\prod_ks(t_k,t_{k+1})$, that is,
    $U=\prod_j\mathrm{CZ}_{j,j+1}$, and $U^2=\Id$. The tensor $M$ is normal, so
    $\{\Id,M\}$ represents $\Z_2$ exactly by normal tensors. The product state
    $\ket{0\cdots0}$ is normal and fixed by $U$. Hence, for every choice of
    fusion tensors, the anomaly three-cocycle has trivial gauge class and the
    anomaly class in $H^3(\Z_2,\C^\times)$ vanishes, as stated at
    \cite[line~3873]{FrancoRubio2025MPUGauging}.
\end{lemma}

\begin{proof}\leanok
    \uses{thm:mpug_anomaly_obstruction, thm:mps_forced_column_normality}
    The outgoing bond of $M$ copies the input, so the periodic operator is the
    product of the entries along the forced bonds, $s(t_k,t_{k+1})$ at each
    site, and squares to the identity since $s^2=1$. The letter $(j,j)$ of $M$
    has the single nonzero column $j$, equal to the column $j$ of $\sqrt2H$, an
    invertible matrix with no vanishing entry, so $M$ is normal by
    Lemma~\ref{thm:mps_forced_column_normality}. Every entry
    $\prod_ks(0,0)=1$ fixes $\ket{0\cdots0}$, and
    Theorem~\ref{thm:mpug_anomaly_obstruction} applies.
\end{proof}

\begin{definition}[The invariant state, its defect, and the gauged tensor]
    \label{def:mpug_controlled_z_defect}
    \lean{ControlledZGauging.stateInt,
        ControlledZGauging.defectInt,
        ControlledZGauging.stateTensor,
        ControlledZGauging.defectTensor,
        ControlledZGauging.gaugedLetter,
        ControlledZGauging.gaugedTensor,
        ControlledZGauging.gaugedTensor_finProdFinEquiv,
        ControlledZGauging.twoSite,
        ControlledZGauging.gaussPerm,
        ControlledZGauging.gaussOperator}
    \leanok
    \uses{def:mpug_controlled_z_tensor}
    As in \cite[lines~3874--3896]{FrancoRubio2025MPUGauging}, let
    \begin{align}
        A^0   & =\begin{pmatrix}1&0&0\\0&0&0\\0&1&1\end{pmatrix},    &
        A^1   & =\begin{pmatrix}0&1&-1\\-1&0&0\\0&0&0\end{pmatrix},
        \notag                                                         \\
        A_g^0 & =\begin{pmatrix}0&-1&-1\\0&0&0\\-1&0&0\end{pmatrix}, &
        A_g^1 & =\begin{pmatrix}1&0&0\\0&-1&1\\0&0&0\end{pmatrix}.
        \notag
    \end{align}
    The gauged tensor has letters $C^{ie}=A^i$ and $C^{ig}=A_g^i$, with the
    gauge label written as a bit \cite[line~3966]{FrancoRubio2025MPUGauging}.
    The Gauss operator $\mathcal G_g=\mathrm{CZ}\otimes X\otimes X$ acts on two
    gauged sites by $\mathrm{CZ}$ on the two matter labels and $X$ on each
    gauge label \cite[line~3968]{FrancoRubio2025MPUGauging}.
\end{definition}

\begin{lemma}[Defect identities and the Gauss law of the controlled-\texorpdfstring{$Z$}{Z} gauging]
    \label{thm:mpug_controlled_z_gauss_law}
    \lean{ControlledZGauging.sign_smul_state_mul_state,
        ControlledZGauging.sign_smul_defect_mul_state,
        ControlledZGauging.stateTensor_isNormal,
        ControlledZGauging.czSign_smul_gaugedLetter_mul,
        ControlledZGauging.gaussOperator_mul_self,
        ControlledZGauging.sum_gaussOperator_smul_twoSite}
    \leanok
    \uses{def:mpug_controlled_z_defect}
    The tensor $A$ is normal \cite[line~3885]{FrancoRubio2025MPUGauging}, and
    for all bits $i,j$ \cite[eq.~(Z2\_def)]{FrancoRubio2025MPUGauging},
    \begin{align}
        s(i,j)\,A^iA^j   & =A_g^iA_g^j, &
        s(i,j)\,A_g^iA^j & =A^iA_g^j.
        \label{eq:mpug_controlled_z_defect}
    \end{align}
    The Gauss operator satisfies $\mathcal G_g^2=\Id$ and fixes the two-site
    contraction of the gauged tensor:
    $\sum_x(\mathcal G_g)_{y,x}C^{x_1}C^{x_2}=C^{y_1}C^{y_2}$ for every pair
    $y=(y_1,y_2)$ of gauged labels.
\end{lemma}

\begin{proof}\leanok
    \uses{def:mpug_controlled_z_defect}
    The identities (\ref{eq:mpug_controlled_z_defect}) are identities of
    integer matrices, checked on the four pairs $(i,j)$. Twice every matrix
    unit is a signed sum of at most four words of $A$ of length four, so the
    words of length four span all matrices and $A$ is normal. By
    (\ref{eq:mpug_controlled_z_defect}) and $s^2=1$, flipping both gauge labels
    multiplies $C^{ia}C^{jb}$ by $s(i,j)$, which is the Gauss law in letters.
    The operator $\mathcal G_g$ is the monomial matrix of the double gauge
    flip with phases $s(i,j)$, so it is an involution, and its single nonzero
    entry in each row turns the Gauss law in letters into the stated
    invariance.
\end{proof}

\begin{lemma}[The split-index form of the gauged tensor]
    \label{thm:mpug_controlled_z_split_index}
    \lean{ControlledZGauging.simpsRInt,
        ControlledZGauging.simpsLInt,
        ControlledZGauging.xInt,
        ControlledZGauging.simpsR,
        ControlledZGauging.simpsL,
        ControlledZGauging.simpsGauge,
        ControlledZGauging.simpsTensor,
        ControlledZGauging.stateTensor_eq_simpsR_mul_simpsL,
        ControlledZGauging.defectTensor_eq_simpsR_mul_mul_simpsL,
        ControlledZGauging.gaugedLetter_eq_simps,
        ControlledZGauging.xInt_mul_simpsTensor_mul_xInt}
    \leanok
    \uses{def:mpug_controlled_z_defect}
    Let
    \begin{align}
        R^0 & =\begin{pmatrix}1&0\\0&0\\0&-1\end{pmatrix}, &
        R^1 & =\begin{pmatrix}1&0\\0&-1\\0&0\end{pmatrix}, &
        L^0 & =\begin{pmatrix}1&0&0\\0&-1&-1\end{pmatrix}, &
        L^1 & =\begin{pmatrix}0&1&-1\\1&0&0\end{pmatrix}.
        \notag
    \end{align}
    Then $A^i=R^iL^i$, $A_g^i=R^iXL^i$, and $C^{ia}=R^iX^aL^i$, the internal
    tensor of the virtual representation $S_e=\Id$, $S_g=X$ of
    \cite[lines~4073--4209]{FrancoRubio2025MPUGauging}. The split-index tensor
    $S^{ab}=L^aR^b$ satisfies $XS^{ab}X=s(a,b)\,S^{ab}$, the identity of
    \cite[lines~4002--4072]{FrancoRubio2025MPUGauging}. The source draws $L$
    and $R$ without listing their entries; the matrices above are one choice
    realizing the displayed identities.
\end{lemma}

\begin{proof}\leanok
    \uses{def:mpug_controlled_z_defect}
    Each identity is an identity of integer matrices, checked on the finitely
    many values of the bits. For instance $S^{00}=S^{01}=\Id$, $S^{10}=X$ and
    $S^{11}=\left(\begin{smallmatrix}0&-1\\1&0\end{smallmatrix}\right)$, which
    anticommutes with $X$.
\end{proof}

\begin{lemma}[Split-index presentation of a periodic amplitude]
    \label{lem:mps_split_index_trace}
    \lean{MPSTensor.prod_ofFn_mul_eq_mul_prod_mul,
        MPSTensor.mpv_eq_trace_prod_ofFn_of_eq_mul,
        MPSTensor.mul_prod_ofFn_mul_eq_smul,
        MPSTensor.prod_mul_trace_prod_ofFn_eq}
    \leanok
    \uses{def:mps_tensor}
    Let $A$ be a tensor of bond dimension $D$ whose letters factor as
    $A^i=R^iL^i$ through an intermediate bond of dimension $m$, with $R^i$ a
    $D\times m$ and $L^i$ an $m\times D$ matrix, and let
    $S^{ab}=L^aR^b$. On a ring of $N\geq1$ sites,
    \begin{align}
        \tr\bigl(A^{\sigma_0}A^{\sigma_1}\cdots A^{\sigma_{N-1}}\bigr)
         & =\tr\bigl(S^{\sigma_0\sigma_1}S^{\sigma_1\sigma_2}\cdots
        S^{\sigma_{N-1}\sigma_0}\bigr).
        \label{eq:mps_split_index_trace}
    \end{align}
    If $XY=YX=\Id$ and $XS_kY=c_kS_k$ for matrices $S_1,\dots,S_N$ and scalars
    $c_1,\dots,c_N$, then $X(S_1\cdots S_N)Y=(c_1\cdots c_N)\,S_1\cdots S_N$,
    and hence $c_1\cdots c_N\tr(S_1\cdots S_N)=\tr(S_1\cdots S_N)$.
\end{lemma}

\begin{proof}\leanok
    The product $A^{\sigma_0}\cdots A^{\sigma_{N-1}}$ regroups as
    $R^{\sigma_0}\bigl(L^{\sigma_0}R^{\sigma_1}\bigr)\cdots
        \bigl(L^{\sigma_{N-2}}R^{\sigma_{N-1}}\bigr)L^{\sigma_{N-1}}$, and the
    trace is invariant under moving the rectangular factor $R^{\sigma_0}$ to
    the end. For the second statement insert $YX=\Id$ between consecutive
    factors, and then use $YX=\Id$ once more after moving $X$ to the end
    inside the trace.
\end{proof}

\begin{lemma}[The controlled-\texorpdfstring{$Z$}{Z} state as a split-index state]
    \label{lem:mpug_controlled_z_simps_trace}
    \lean{ControlledZGauging.zInt,
        ControlledZGauging.simpsTensor_eq,
        ControlledZGauging.mpv_stateTensor_eq_trace_simpsTensor,
        ControlledZGauging.czPhase_mul_mpv_stateTensor,
        ControlledZGauging.fixesMPV_czMPUTensor_stateTensor,
        ControlledZGauging.fixesMPV_stateTensor}
    \leanok
    \uses{thm:mpug_controlled_z_split_index, thm:mpug_controlled_z_not_anomalous,
        def:mpug_single_block_l_symbol}
    The split-index tensor of Lemma~\ref{thm:mpug_controlled_z_split_index} is
    $S^{ab}=X^aZ^{ab}$, and on a ring of $N\geq1$ qubits the amplitude of the
    state generated by $A$ is
    $\tr\bigl(S^{\sigma_0\sigma_1}S^{\sigma_1\sigma_2}\cdots
        S^{\sigma_{N-1}\sigma_0}\bigr)$, the split-index presentation drawn in
    \cite[lines~3970--4001]{FrancoRubio2025MPUGauging}. Every amplitude is
    fixed by the controlled-$Z$ phase,
    $\prod_ks(\sigma_k,\sigma_{k+1})\,\psi(\sigma)=\psi(\sigma)$, so the state
    is fixed by $U$ and by every operator of the family $\{\Id,M\}$ on every
    ring of $N\geq1$ qubits.
\end{lemma}

\begin{proof}\leanok
    \uses{lem:mps_split_index_trace, thm:mpug_controlled_z_split_index,
        thm:mpug_controlled_z_not_anomalous}
    The formula for $S^{ab}$ is an identity of integer matrices. The
    amplitude is the case $A^i=R^iL^i$ of
    Lemma~\ref{lem:mps_split_index_trace}. Since $XS^{ab}X=s(a,b)\,S^{ab}$
    and $X^2=\Id$, the second part of that lemma gives the invariance of each
    amplitude under the phase, and the periodic operator of $M$ is the
    diagonal matrix of these phases.
\end{proof}

\begin{lemma}[The onsite symmetry \texorpdfstring{$\prod_jZ_j$}{prod Z}]
    \label{lem:mpug_controlled_z_onsite_z}
    \lean{ControlledZGauging.zSign,
        ControlledZGauging.zAction,
        ControlledZGauging.zAction_ofAdd_one,
        ControlledZGauging.stateZGaugeInt,
        ControlledZGauging.stateZGauge,
        ControlledZGauging.twistedTensor_stateTensor_zAction,
        ControlledZGauging.stateZGauge_mul_stateTensor_mul_inv,
        ControlledZGauging.stateTensor_isOnSiteSymmetric_zAction}
    \leanok
    \uses{def:mpug_controlled_z_defect, def:twisted_tensor, def:on_site_symmetric}
    Let $\Z_2$ act on a qubit by $g\mapsto Z^g$, with
    $Z=\operatorname{diag}(1,-1)$. Then
    $\sum_jZ_{ij}A^j=(-1)^iA^i=VA^iV^{-1}$ for
    $V=\operatorname{diag}(1,-1,-1)$, so the state generated by $A$ has the
    onsite symmetry $\prod_jZ_j$ stated in
    \cite[line~3970]{FrancoRubio2025MPUGauging}.
\end{lemma}

\begin{proof}\leanok
    \uses{def:gauge_equiv, thm:gauge_invariance}
    The twisted letters are $Z_{ii}A^i=(-1)^iA^i$. Conjugation by $V$
    multiplies the entry $(r,c)$ of a matrix by $V_{rr}V_{cc}$; the nonzero
    entries of $A^0$ sit at $(0,0)$, $(2,1)$ and $(2,2)$, and those of $A^1$
    at $(0,1)$, $(0,2)$ and $(1,0)$, which gives the signs $+1$ and $-1$.
    Gauge-equivalent tensors generate the same states.
\end{proof}

\begin{lemma}[The combined \texorpdfstring{$\Z_2\times\Z_2$}{Z2 x Z2} symmetry on the split-index bond]
    \label{lem:mpug_controlled_z_combined_class}
    \lean{ControlledZGauging.splitLegPhase,
        ControlledZGauging.clusterProjRep_mul_simpsTensor,
        ControlledZGauging.simpsTensor_isNontrivialClass}
    \leanok
    \uses{lem:mpug_controlled_z_simps_trace, def:cluster_omega}
    Let $g=(g_1,g_2)\in\Z_2\times\Z_2$ act on the two legs $(a,b)$ of a
    split-index tensor by the phase
    $\chi_g(a,b)=(-1)^{g_1a}\,s(a,b)^{g_2}$: the first generator is the onsite
    symmetry $\prod_jZ_j$, whose factor at each site is assigned to the left
    leg of the split-index tensor at that site, and the second is $U$. Let
    $V_g$ be the projective representation of $\Z_2\times\Z_2$ on the
    split-index bond with $V_{(1,0)}=Z$ and $V_{(0,1)}=X$, whose factor
    system is that of Definition~\ref{def:cluster_omega}. Then
    \begin{align}
        V_gS^{ab} & =\chi_g(a,b)\,S^{ab}V_g
        \label{eq:mpug_controlled_z_combined_class}
    \end{align}
    for all $g$, $a$, $b$, and the class of this factor system in
    $H^2(\Z_2\times\Z_2,\C^\times)$ is nontrivial. This is the invariant
    behind the remark of \cite[line~3970]{FrancoRubio2025MPUGauging} that the
    state lies in the nontrivial symmetry-protected phase of the combined
    $\Z_2\times\Z_2$ symmetry. That this class labels the phase of a state
    protected by a symmetry that is not onsite is the classification of
    \cite{Stephen2024SIMPS}, which is not proved here; see
    \cite{gap:fbc25_controlled_z_split_index_class_scope}.
\end{lemma}

\begin{proof}\leanok
    \uses{thm:cluster_nontrivial_spt}
    By Lemma~\ref{lem:mpug_controlled_z_simps_trace}, $S^{ab}=X^aZ^{ab}$, so
    (\ref{eq:mpug_controlled_z_combined_class}) is an identity of integer
    matrices, checked on the generators, on their product and on the bits
    $a,b$. The matrices $Z$ and $X$ anticommute, so the commutator phase of
    the factor system on the two generators is $-1$, and the class is
    nontrivial by Lemma~\ref{thm:cluster_nontrivial_spt}.
\end{proof}

\section{The gauge-invariant subspace of the CZX circuit tuple}
\label{sec:mpug_czx_gauge_invariant}

The CZX model of \cite[lines~4503--5183]{FrancoRubio2025MPUGauging} is the
anomalous $\Z_2=\{e,g\}$ matrix product unitary acting on the GHZ state.
After blocking, one site carries two matter qubits, and the fusion operators on
two neighboring blocked sites are $\lambda^R_{e,e}=\lambda^R_{e,g}=\Id$,
$\lambda^R_{g,e}=w_R$, together with the modified fusion operator
$\tilde\lambda^R_{g,g}=\lambda^R_{g,g}Z_1$ introduced at the end of that
passage.
For the displayed circuit tuple $(\Id,\Id,w,\tilde\lambda)$, the common $+1$
eigenspace of the resulting Gauss projectors has dimension $2^N$ on
$N\geq3$ blocked sites, answering the scaling question of
\cite[lines~5198--5204]{FrancoRubio2025MPUGauging} for this tuple only. The
displayed tuple is not identified with a member of the full physical
completion class of the CZX defect data; that identification requires the
three defect domains and maps, which the displayed matrices do not specify.

\subsection{Monomial operators and their common fixed space}

\begin{definition}[Monomial matrix]
    \label{def:mpug_monomial_matrix}
    \lean{Matrix.monomial,
        Matrix.monomial_apply,
        Matrix.monomial_apply_perm,
        Matrix.monomial_apply_of_ne}
    \leanok
    Let $\sigma$ be a permutation of a finite index set $S$ and let
    $\varphi\colon S\to\C$. The monomial matrix $M(\sigma,\varphi)$ is the
    operator on $\C^S$ with
    \begin{align}
        M(\sigma,\varphi)\ket{s} & =\varphi(s)\ket{\sigma(s)}.
        \label{eq:mpug_monomial_matrix}
    \end{align}
    Its entry in row $\sigma(s)$ and column $s$ is $\varphi(s)$, and every
    other entry vanishes.
\end{definition}

\begin{lemma}[Calculus of monomial matrices]
    \label{thm:mpug_monomial_matrix_laws}
    \lean{Matrix.monomial_mulVec,
        Matrix.monomial_mulVec_eq_self_iff,
        Matrix.monomial_mul_monomial,
        Matrix.conjTranspose_monomial,
        Matrix.smul_monomial,
        Matrix.neg_monomial,
        Matrix.monomial_one,
        Matrix.reindex_monomial,
        Matrix.monomial_mem_unitaryGroup,
        Matrix.trace_monomial,
        Matrix.trace_monomial_eq_zero}
    \leanok
    \uses{def:mpug_monomial_matrix}
    A vector $v\in\C^S$ is fixed by $M(\sigma,\varphi)$ if and only if
    $v_{\sigma(s)}=\varphi(s)v_s$ for every $s\in S$. Products, adjoints,
    scalar multiples, negatives, and simultaneous relabellings of monomial
    matrices are monomial:
    \begin{align}
        M(\sigma,\varphi)M(\tau,\psi)
         & =M\bigl(\sigma\tau,\ s\mapsto\varphi(\tau(s))\psi(s)\bigr),
        \label{eq:mpug_monomial_mul}                                               \\
        M(\sigma,\varphi)^\dagger
         & =M\bigl(\sigma^{-1},\ t\mapsto\overline{\varphi(\sigma^{-1}(t))}\bigr).
        \label{eq:mpug_monomial_adjoint}
    \end{align}
    The identity is $M(\id,1)$, and $M(\sigma,\varphi)$ is unitary when
    $|\varphi(s)|=1$ for every $s$. The trace is
    $\tr M(\sigma,\varphi)=\sum_{s:\,\sigma(s)=s}\varphi(s)$, which vanishes
    when $\sigma$ has no fixed point.
\end{lemma}

\begin{proof}\leanok
    \uses{def:mpug_monomial_matrix}
    The action of $M(\sigma,\varphi)$ on a vector is
    $(M(\sigma,\varphi)v)_t=\varphi(\sigma^{-1}(t))v_{\sigma^{-1}(t)}$, which
    gives the fixed-vector criterion after the substitution $t=\sigma(s)$.
    The remaining identities are entrywise: in a product only the column
    $\tau(s)$ of the first factor meets the row $\tau(s)$ of the second, the
    adjoint transposes the single nonzero entry of each column, and a
    relabelling conjugates the permutation. For unitarity,
    (\ref{eq:mpug_monomial_mul}) and (\ref{eq:mpug_monomial_adjoint}) give
    $M(\sigma,\varphi)^\dagger M(\sigma,\varphi)=M(\id,|\varphi|^2)=\Id$.
    The diagonal entry at $s$ is $\varphi(s)$ if $\sigma(s)=s$ and zero
    otherwise, which gives the trace.
\end{proof}

\begin{lemma}[Flat phases on the hypercube have a potential]
    \label{thm:mpug_hypercube_potential}
    \lean{TNLean.Algebra.hypercube_induction,
        TNLean.Algebra.exists_potential_of_flat}
    \leanok
    Let $M$ be a commutative group and let
    $\varphi_j\colon\mathbb F_2^r\to M$ for $0\leq j<r$ be phases attached
    to the edges $(\gamma,\gamma+e_j)$ of the hypercube. Assume that
    \begin{align}
        \varphi_j(\gamma)\varphi_j(\gamma+e_j)
         & =1,
        \label{eq:mpug_hypercube_involutive}       \\
        \varphi_j(\gamma)\varphi_k(\gamma+e_j)
         & =\varphi_k(\gamma)\varphi_j(\gamma+e_k)
        \label{eq:mpug_hypercube_flat}
    \end{align}
    for all $j,k$ and $\gamma$. Then there is $\Phi\colon\mathbb F_2^r\to M$
    with $\Phi(0)=1$ and $\Phi(\gamma+e_j)=\varphi_j(\gamma)\Phi(\gamma)$ for
    every $j$ and $\gamma$. In addition, a property of vectors of
    $\mathbb F_2^r$ that holds at $0$ and is preserved by every coordinate
    flip holds everywhere.
\end{lemma}

\begin{proof}\leanok
    For the connectivity claim, every $\gamma$ is the sum of the $e_j$ over
    its support, so induction on the support carries the property from $0$ to
    $\gamma$.

    The potential is constructed by induction on $r$. Writing
    $\gamma=(a,\gamma')$ with $a\in\mathbb F_2$ and
    $\gamma'\in\mathbb F_2^{r-1}$, the phases
    $\gamma'\mapsto\varphi_{j+1}(0,\gamma')$ for $0\leq j<r-1$ satisfy
    (\ref{eq:mpug_hypercube_involutive}) and (\ref{eq:mpug_hypercube_flat})
    on $\mathbb F_2^{r-1}$, so the induction hypothesis gives a potential
    $\Phi'$ for them. Setting
    \begin{align}
        \Phi(0,\gamma') & =\Phi'(\gamma'),
        \qquad
        \Phi(1,\gamma')=\varphi_0(0,\gamma')\Phi'(\gamma'),
        \label{eq:mpug_hypercube_potential_step}
    \end{align}
    the edge relation in direction $0$ at $(0,\gamma')$ holds by definition,
    and at $(1,\gamma')$ it follows from
    (\ref{eq:mpug_hypercube_involutive}). The edge relation in a direction
    $j+1$ at $(0,\gamma')$ is the relation for $\Phi'$, and at $(1,\gamma')$
    it is the relation for $\Phi'$ combined with
    (\ref{eq:mpug_hypercube_flat}) for the square spanned by the directions
    $0$ and $j+1$ at $(0,\gamma')$.
\end{proof}

\begin{definition}[Trivial holonomy on a fiber]
    \label{def:mpug_trivial_holonomy}
    \lean{TNLean.Algebra.fiberFlip,
        TNLean.Algebra.fiberFlip_apply,
        TNLean.Algebra.IsTrivialHolonomy}
    \leanok
    Let $X$ be a finite set and consider the basis $X\times\mathbb F_2^r$.
    For $0\leq j<r$ let $\sigma_j$ be the coordinate flip
    $(x,\gamma)\mapsto(x,\gamma+e_j)$ and let
    $\varphi_j\colon X\times\mathbb F_2^r\to\C$. The phases have trivial
    holonomy on the fiber over $x\in X$ when
    \begin{align}
        \varphi_j(x,\gamma)\varphi_k(x,\gamma+e_j)
         & =\varphi_k(x,\gamma)\varphi_j(x,\gamma+e_k)
        \label{eq:mpug_trivial_holonomy}
    \end{align}
    for all $j,k$ and $\gamma\in\mathbb F_2^r$.
\end{definition}

\begin{lemma}[Monomial fixed-space criterion]
    \label{thm:mpug_monomial_fixed_space_criterion}
    \lean{TNLean.Algebra.reindex_mulVec,
        TNLean.Algebra.commonFixedSubmodule_reindex,
        TNLean.Algebra.finrank_commonFixedSubmodule_reindex,
        TNLean.Algebra.add_self_eq_zero_pi_zmod_two,
        TNLean.Algebra.eq_zero_of_mem_commonFixedSubmodule_monomial_fiberFlip,
        TNLean.Algebra.finrank_commonFixedSubmodule_monomial_fiberFlip}
    \leanok
    \uses{def:mpug_monomial_matrix, def:mpug_trivial_holonomy,
        def:mpug_common_fixed_subspace}
    In the setting of Definition~\ref{def:mpug_trivial_holonomy}, let
    $T_j=M(\sigma_j,\varphi_j)$ and assume that every $T_j$ is an involution,
    that is, $\varphi_j(s)\varphi_j(\sigma_j(s))=1$ for every basis element
    $s$. Then
    \begin{align}
        \dim\mathcal V(T)
         & =\#\{x\in X:\text{the phases have trivial holonomy on the fiber over }x\}.
        \label{eq:mpug_monomial_fixed_space_count}
    \end{align}
    The dimension of a common fixed subspace is also unchanged by a
    simultaneous relabelling of the basis.
\end{lemma}

\begin{proof}\leanok
    \uses{thm:mpug_monomial_matrix_laws, thm:mpug_hypercube_potential,
        thm:mpug_mem_common_fixed_subspace}
    By Lemma~\ref{thm:mpug_monomial_matrix_laws}, a vector $v$ lies in
    $\mathcal V(T)$ exactly when
    $v_{(x,\gamma+e_j)}=\varphi_j(x,\gamma)v_{(x,\gamma)}$ for all $j$, $x$,
    and $\gamma$. Let $X_0$ be the set of fibers with trivial holonomy and
    consider the linear map $E\colon\mathcal V(T)\to\C^{X_0}$,
    $v\mapsto(v_{(x,0)})_{x\in X_0}$.

    The map $E$ is injective: if $v\in\mathcal V(T)$ vanishes at one point
    of a fiber, the edge relation and the connectivity statement of
    Lemma~\ref{thm:mpug_hypercube_potential} make it vanish on the whole
    fiber. On a fiber with trivial holonomy, $v$ therefore vanishes as soon
    as $v_{(x,0)}=0$. On a fiber where (\ref{eq:mpug_trivial_holonomy})
    fails at some $(j,k,\gamma_0)$, the two paths from $(x,\gamma_0)$ to
    $(x,\gamma_0+e_j+e_k)$ give
    \begin{align}
        v_{(x,\gamma_0+e_j+e_k)} & =\varphi_j(x,\gamma_0)\varphi_k(x,\gamma_0+e_j)\,v_{(x,\gamma_0)}, \\
        v_{(x,\gamma_0+e_j+e_k)} & =\varphi_k(x,\gamma_0)\varphi_j(x,\gamma_0+e_k)\,v_{(x,\gamma_0)}.
    \end{align}
    The two phase products differ, so $v_{(x,\gamma_0)}=0$ and again $v$
    vanishes on the fiber.

    The map $E$ is also surjective: given $w\in\C^{X_0}$,
    Lemma~\ref{thm:mpug_hypercube_potential} provides a potential $\Phi_x$
    on each fiber with trivial holonomy for the invertible phases
    $\gamma\mapsto\varphi_j(x,\gamma)$, and setting
    $v_{(x,\gamma)}=w_x\Phi_x(\gamma)$ there and $v=0$ on the other fibers
    satisfies the potential edge relation, so $v\in\mathcal V(T)$,
    and $\Phi_x(0)=1$ gives $Ev=w$.

    Hence $\dim\mathcal V(T)=\dim\C^{X_0}=\#X_0$. For the relabelling
    statement, a relabelled matrix acts by relabelling the vector, acting,
    and relabelling back, so the common fixed subspace is the image of the
    original one under the coordinate relabelling.
\end{proof}

\begin{lemma}[Local Gauss operator of monomial matter unitaries]
    \label{thm:mpug_gauss_operator_monomial}
    \lean{TNLean.Algebra.gaussMonomialPerm,
        TNLean.Algebra.gaussMonomialPerm_apply,
        TNLean.Algebra.gaussMonomialPhase,
        TNLean.Algebra.gaussOperator_eq_monomial}
    \leanok
    \uses{def:mpug_gauss_operator, def:mpug_monomial_matrix}
    Suppose that every matter unitary of the tuple $R$ is monomial,
    $R_{a,b}=M(\sigma_{a,b},\varphi_{a,b})$. Then the local Gauss operator
    of Definition~\ref{def:mpug_gauss_operator} is the monomial operator
    \begin{align}
        \mathcal G_g\ket{i,(a,b)}
         & =\overline{\varphi_{T_g(a,b)}\bigl(\sigma_{T_g(a,b)}^{-1}\sigma_{a,b}(i)\bigr)}\,
        \varphi_{a,b}(i)\,
        \ket{\sigma_{T_g(a,b)}^{-1}\sigma_{a,b}(i),\ T_g(a,b)}.
        \label{eq:mpug_gauss_operator_monomial}
    \end{align}
\end{lemma}

\begin{proof}\leanok
    \uses{thm:mpug_monomial_matrix_laws, def:mpug_gauss_operator}
    The block of $\mathcal G_g$ from $(a,b)$ to $T_g(a,b)$ is
    $R_{T_g(a,b)}^\dagger R_{a,b}$, which is monomial by
    (\ref{eq:mpug_monomial_mul}) and (\ref{eq:mpug_monomial_adjoint}); all
    other blocks vanish.
\end{proof}

\begin{lemma}[Periodic placement of monomial operators]
    \label{thm:mpug_embed_monomial}
    \lean{MPOTensor.windowPerm,
        MPOTensor.windowPerm_apply,
        MPOTensor.embedLocalOperator_monomial,
        MPOTensor.extractWindow_two,
        MPOTensor.replaceWindow_two_apply}
    \leanok
    \uses{def:mpug_monomial_matrix}
    Let $M(\sigma,\varphi)$ be a monomial operator on $L$ consecutive sites.
    Its placement on the periodic window beginning at site $i$ of an $N$-site
    chain is the monomial operator whose permutation replaces the window of a
    configuration $s$ by $\sigma$ applied to that window and whose phase is
    $\varphi$ evaluated on the window of $s$. For $L=2$ and $N\geq2$ the
    window at $j$ is $(s_j,s_{j+1})$, and replacing it changes exactly the
    sites $j$ and $j+1$.
\end{lemma}

\begin{proof}\leanok
    \uses{def:mpug_monomial_matrix}
    The placed operator has entry $(t,s)$ equal to the window entry
    $M(\sigma,\varphi)_{t_W,s_W}$ when $t$ and $s$ agree outside the window
    and $0$ otherwise. Both conditions together say that $t$ is $s$ with its
    window replaced by $\sigma(s_W)$, and the entry is then
    $\varphi(s_W)$. The two-site statements are the cyclic offset
    computations $(j+N-j)\bmod N=0$ and $(j+1+N-j)\bmod N=1$.
\end{proof}

\subsection{The exact CZX matrix product unitary}

\begin{definition}[The once-blocked CZX tensor]
    \label{def:mpug_czx_exact_tensor}
    % Principal declaration: MPOTensor.CZX.tensor.
    \lean{MPOTensor.CZX.decoratedSiteTensor, MPOTensor.CZX.tensor,
        MPOTensor.CZX.complementSite, MPOTensor.CZX.complement,
        MPOTensor.CZX.edgeExponent, MPOTensor.CZX.cyclicExponent}
    \leanok
    \uses{def:mpdo_two_site_blocking}
    Write a physical index as $2a+b$, with $a,b\in\{0,1\}$, and let
    $\overline{(a,b)}=(1-a,1-b)$. The tensor in
    \cite[lines~4503--4547]{FrancoRubio2025MPUGauging}
    is the two-site block $B^{ab,cd}=C^{a,c}C^{b,d}$ of the bond-two tensor
    \begin{align}
        C^{u,d}_{\ell r}
         & =\delta_{u,1-d}\delta_{\ell,u}(-1)^{u+\ell r}.
    \end{align}
    Here the constituent indices $u,d,\ell,r$ are binary, the weight is
    $W_{\ell r}=(-1)^{\ell r}=\sqrt{2}H_{\ell r}$, where $H$ is the
    normalized Hadamard matrix, and the factor $(-1)^u$ is the output
    Pauli $Z$. Thus, if $P$ denotes the normalized tensor printed in
    \cite[lines~1906--1943]{Cirac2017MPU}, then
    $B=(Z\otimes Z)\operatorname{block}_2(\sqrt{2}P)$, not
    $\operatorname{block}_2(P)$.
    Define the local integer exponent $e((a,b),r)=a+b+ab+br$.
    For a periodic configuration $s_j=(a_j,b_j)$ of positive length $N$,
    let $\kappa_N(s)_j=\overline{s_j}$ and set
    \begin{align}
        E(s) & =\sum_{j=0}^{N-1}e(s_j,a_{j+1})
        =\sum_{j=0}^{N-1}(a_j+b_j+a_jb_j+b_ja_{j+1}),
    \end{align}
    where $j+1$ is interpreted modulo $N$ and the bit values in the sum
    are integers.
\end{definition}

\begin{lemma}[Coordinates and unitarity of the exact CZX tensor]
    \label{thm:mpug_czx_exact_tensor_unitary}
    % Principal declaration: MPOTensor.CZX.tensor_isMPU.
    \lean{MPOTensor.CZX.tensor_apply, MPOTensor.CZX.tensor_apply_bits,
        MPOTensor.CZX.mpo_tensor_apply, MPOTensor.CZX.mpo_tensor,
        MPOTensor.CZX.mpo_tensor_mem_unitaryGroup, MPOTensor.CZX.tensor_isMPU}
    \leanok
    \uses{def:mpug_czx_exact_tensor, def:mpug_monomial_matrix, def:mpu}
    The tensor $B$ has bond dimension two and coordinates
    \begin{align}
        B^{ab,cd}_{\ell r}
         & =\delta_{a,1-c}\delta_{b,1-d}\delta_{\ell,a}
        (-1)^{a+b+ab+br}.
    \end{align}
    Its periodic operator is unitary for every $N>0$, with entries and
    monomial form
    \begin{align}
        O_N(B)_{s,t} & =\delta_{s,\kappa_N(t)}(-1)^{E(s)},                     \\
        O_N(B)       & =M\bigl(\kappa_N,\ t\mapsto(-1)^{E(\kappa_N(t))}\bigr).
    \end{align}
    In particular, $B$ is a matrix product unitary. This assertion does
    not include its phase classification, GHZ action, or defect data.
\end{lemma}
\begin{proof}
    \leanok
    \uses{thm:trace_eval_word_eq_sum_cyclic, thm:mpug_monomial_matrix_laws}
    The two-site block is the matrix product contracted over a binary
    intermediate bond $m$. The second factor vanishes unless $m=b$, so
    \begin{align}
        B^{ab,cd}_{\ell r}
         & =\sum_{m\in\{0,1\}}C^{a,c}_{\ell m}C^{b,d}_{mr} \\
         & =C^{a,c}_{\ell b}C^{b,d}_{br}                   \\
         & =\delta_{a,1-c}\delta_{b,1-d}\delta_{\ell,a}
        (-1)^{a+\ell b+b+br}                               \\
         & =\delta_{a,1-c}\delta_{b,1-d}\delta_{\ell,a}
        (-1)^{a+b+ab+br}.
    \end{align}
    For the periodic contraction with $s_j=(a_j,b_j)$ and
    $t_j=(c_j,d_j)$, the cyclic trace expansion is
    \begin{align}
        O_N(B)_{s,t}
         & =\sum_{g:\Z/N\Z\to\{0,1\}}
        \prod_{j=0}^{N-1}B^{s_j,t_j}_{g_j,g_{j+1}}.
    \end{align}
    The factor at $j$ vanishes unless $g_j=a_j$, so only the path
    $g_j=a_j$ can contribute. Substituting it gives
    \begin{align}
        O_N(B)_{s,t}
         & =\prod_{j=0}^{N-1}B^{s_j,t_j}_{a_j,a_{j+1}} \\
         & =\prod_{j=0}^{N-1}
        \delta_{a_j,1-c_j}\delta_{b_j,1-d_j}
        (-1)^{a_j+b_j+a_jb_j+b_ja_{j+1}}               \\
         & =\delta_{s,\kappa_N(t)}(-1)^{E(s)}.
    \end{align}
    Thus the operator is the monomial matrix for complementation, with
    column phase $(-1)^{E(\kappa_N(t))}$. Each phase has absolute value one,
    so Lemma~\ref{thm:mpug_monomial_matrix_laws} proves unitarity.
\end{proof}

\begin{lemma}[Injectivity of the once-blocked CZX tensor]
    \label{thm:mpug_czx_exact_tensor_injective}
    \lean{MPOTensor.CZX.tensor_isInjective}
    \leanok
    \discussion{7354}
    \uses{def:mpug_czx_exact_tensor, def:injective}
    As in \cite[lines~4505--4547]{FrancoRubio2025MPUGauging}, the
    once-blocked CZX tensor $B$ of Definition~\ref{def:mpug_czx_exact_tensor},
    with the unnormalized Hadamard weights, both output $Z$ gates, and
    physical indexing $B^{2a+b,2c+d}=B^{ab,cd}$ for binary $a,b,c,d$, is
    injective:
    \begin{align}
        \operatorname{span}_{\C}\{B^{i,j}:0\leq i,j<4\}
         & =\operatorname{Mat}_2(\C).
    \end{align}
    % Source: eq:MPU_CZX. This proves only injectivity of the shared tensor;
    % it does not complete the GHZ, fusion, or action-data scope of #7354.
\end{lemma}

\begin{proof}\leanok
    \uses{thm:mpug_czx_exact_tensor_unitary}
    Let $E_{ab}$ be the bond matrix units, with $a,b\in\{0,1\}$.
    The coordinate formula in
    Lemma~\ref{thm:mpug_czx_exact_tensor_unitary} gives
    \begin{align}
        E_{00}
         & =\tfrac12(B^{0,3}-B^{1,2}),  \\
        E_{01}
         & =\tfrac12(B^{0,3}+B^{1,2}),  \\
        E_{10}
         & =-\tfrac12(B^{2,1}+B^{3,0}), \\
        E_{11}
         & =\tfrac12(B^{3,0}-B^{2,1}).
    \end{align}
    Thus every bond matrix unit lies in the span of the physical letters.
    Since every $2\times2$ matrix is a linear combination of these units,
    the letters span the full bond matrix algebra.
\end{proof}

\subsection{The explicit CZX adjoint gauge}

\begin{definition}[The Pauli adjoint gauge]
    \label{def:mpug_czx_dagger_gauge}
    \lean{MPOTensor.CZX.daggerGauge}
    \leanok
    \discussion{7354}
    On the bond space $\C^2$, with basis indexed by $0,1$, the matrix on
    both virtual legs in the adjoint diagram of
    \cite[lines~4547--4559]{FrancoRubio2025MPUGauging} is the Pauli matrix
    \begin{align}
        Y & =\begin{pmatrix}0&-\mathrm{i}\\ \mathrm{i}&0\end{pmatrix}.
    \end{align}
\end{definition}

\begin{lemma}[Elementary identities for the Pauli adjoint gauge]
    \label{thm:mpug_czx_dagger_gauge_laws}
    \lean{MPOTensor.CZX.daggerGauge_conjTranspose,
        MPOTensor.CZX.daggerGauge_mul_self,
        MPOTensor.CZX.daggerGauge_mem_unitaryGroup}
    \leanok
    \uses{def:mpug_czx_dagger_gauge}
    As used for the unitary gauge in
    \cite[lines~4547--4559]{FrancoRubio2025MPUGauging}, the matrix $Y$ of
    Definition~\ref{def:mpug_czx_dagger_gauge} satisfies
    \begin{align}
        Y^\dagger & =Y,   \\
        Y^2       & =I_2,
    \end{align}
    and in particular $Y\in\mathrm{U}(2)$.
\end{lemma}

\begin{proof}\leanok
    \uses{def:mpug_czx_dagger_gauge}
    Conjugate-transposing the displayed matrix gives $Y^\dagger=Y$, and
    the Pauli multiplication rule $Y^2=I_2$ then yields $YY^\dagger=I_2$.
\end{proof}

\begin{lemma}[The physical adjoint of the exact CZX tensor]
    \label{thm:mpug_czx_physical_adjoint}
    \lean{MPOTensor.CZX.physicalAdjointTensor_tensor}
    \leanok
    \discussion{7354}
    \uses{def:mpug_czx_exact_tensor, def:mpug_czx_dagger_gauge,
        def:mpdo_physical_adjoint}
    Let $B$ be the exact once-blocked CZX tensor of
    Definition~\ref{def:mpug_czx_exact_tensor}, with physical adjoint
    $(B^\dagger)^{u,v}_{\ell r}=\overline{B^{v,u}_{\ell r}}$ in which the
    virtual indices are not transposed. As in the adjoint diagram of
    \cite[lines~4547--4559]{FrancoRubio2025MPUGauging} with no scalar
    prefactor, for $0\leq u,v<4$,
    \begin{align}
        (B^\dagger)^{u,v} & =YB^{u,v}Y.
    \end{align}
\end{lemma}

\begin{proof}\leanok
    \uses{thm:mpug_czx_exact_tensor_unitary,
        def:mpug_czx_dagger_gauge, def:mpdo_physical_adjoint}
    Inserting the coordinate formula from
    Lemma~\ref{thm:mpug_czx_exact_tensor_unitary} on both sides, expanding
    the two matrix products over the binary bond indices, and substituting
    the four physical values and the two bond values verifies each entry.
\end{proof}

\begin{lemma}[An explicit unitary adjoint gauge for CZX]
    \label{thm:mpug_czx_unitary_dagger_gauge}
    \lean{MPOTensor.CZX.daggerGaugeUnitary,
        MPOTensor.CZX.physicalAdjointTensor_tensor_eq_unitaryGauge}
    \leanok
    \uses{def:mpug_czx_exact_tensor, def:mpug_czx_dagger_gauge,
        def:mpdo_physical_adjoint}
    In the unitary gauge convention of
    \cite[lines~1552--1557 and 4547--4559]{FrancoRubio2025MPUGauging}, for
    the exact once-blocked CZX tensor $B$, the matrix $T=Y$, regarded as an
    element of $\mathrm{U}(2)$, satisfies, for $0\leq u,v<4$,
    \begin{align}
        (B^\dagger)^{u,v} & =T^\dagger B^{u,v}T.
    \end{align}
    % This exhibits the explicit witness, not an identification with the
    % choice-selected generic dagger-inverse gauge.
\end{lemma}

\begin{proof}\leanok
    \uses{thm:mpug_czx_dagger_gauge_laws, thm:mpug_czx_physical_adjoint}
    Lemma~\ref{thm:mpug_czx_dagger_gauge_laws} gives $Y\in\mathrm{U}(2)$
    and $Y^\dagger=Y$, and substituting these identities into
    Lemma~\ref{thm:mpug_czx_physical_adjoint} yields the claim.
\end{proof}

\begin{lemma}[The CZX adjoint-gauge sign]
    \label{thm:mpug_czx_dagger_sign}
    \lean{MPOTensor.CZX.daggerGauge_map_star,
        MPOTensor.CZX.daggerGauge_mul_map_star}
    \leanok
    \discussion{7354}
    \uses{def:mpug_czx_dagger_gauge}
    For the Pauli adjoint gauge $Y$, entrywise complex conjugation gives
    \begin{align}
        \overline{Y}  & =-Y,   \\
        Y\overline{Y} & =-I_2.
    \end{align}
    Hence, as stated in \cite[lines~4547--4559]{FrancoRubio2025MPUGauging},
    the sign in $Y\overline{Y}=\sigma_g I_2$ is $\sigma_g=-1$.
    % These entries do not determine fusion tensors, local action tensors,
    % L-symbols, or a gauge selected by the generic existence theorem.
\end{lemma}

\begin{proof}\leanok
    \uses{def:mpug_czx_dagger_gauge, thm:mpug_czx_dagger_gauge_laws}
    Conjugating each entry of the displayed Pauli matrix gives
    $\overline{Y}=-Y$. Then $Y\overline{Y}=-Y^2=-I_2$ by
    Lemma~\ref{thm:mpug_czx_dagger_gauge_laws}.
\end{proof}

\subsection{The displayed CZX source factors}

\begin{definition}[Four explicit CZX source factors]
    \label{def:mpug_czx_displayed_factors}
    \lean{MPOTensor.CZX.displayedX₂, MPOTensor.CZX.displayedY₂,
        MPOTensor.CZX.displayedX₁, MPOTensor.CZX.displayedY₁}
    \leanok
    \discussion{7354}
    \uses{def:mpug_czx_dagger_gauge}
    For $0\leq i<4$, write $i=2a(i)+b(i)$ with $a(i),b(i)\in\{0,1\}$,
    and put $\overline i=3-i$. Let $Y$ be the Pauli matrix of
    Definition~\ref{def:mpug_czx_dagger_gauge}. Define matrices
    $X_2\in\operatorname{Mat}_{(2\times4)\times4}(\C)$ and
    $Y_2\in\operatorname{Mat}_{4\times(4\times2)}(\C)$ by
    \begin{align}
        (X_2)_{(\alpha,i),k}
         & =\delta_{i,k}\delta_{\alpha,a(i)}(-1)^{a(i)+b(i)}, \\
        (Y_2)_{k,(j,\beta)}
         & =\delta_{k,\overline j}(-1)^{a(k)b(k)+b(k)\beta},
    \end{align}
    where $0\leq i,j,k<4$ and $\alpha,\beta\in\{0,1\}$. Define
    $X_1\in\operatorname{Mat}_{(4\times2)\times4}(\C)$ and
    $Y_1\in\operatorname{Mat}_{4\times(2\times4)}(\C)$ by
    \begin{align}
        (X_1)_{(i,\beta),k}
         & =\sum_{t=0}^1Y_{\beta t}\overline{(Y_2)_{k,(i,t)}}, \\
        (Y_1)_{k,(\alpha,j)}
         & =\sum_{t=0}^1\overline{(X_2)_{(t,j),k}}Y_{t\alpha}.
    \end{align}
    Product indices in these matrix spaces retain their displayed order.
    As in the filled and empty factors of
    \cite[lines~4559--4659]{FrancoRubio2025MPUGauging}, the weights
    $(-1)^{ab}$ are unnormalized Hadamard weights, $(-1)^{a(i)+b(i)}$
    includes both output $Z$ gates, and the Pauli entry in the formula for
    $X_1$ is $Y_{\beta t}$ rather than $Y_{t\beta}$.
\end{definition}

\begin{lemma}[The displayed factors recover both CZX cuts]
    \label{thm:mpug_czx_displayed_cuts}
    \lean{MPOTensor.CZX.displayed_cuts}
    \leanok
    \uses{def:mpug_czx_displayed_factors, def:mpug_czx_exact_tensor,
        def:mpu_source_matrix_one, def:mpu_source_matrix_two}
    As in the two decompositions of
    \cite[lines~4559--4659]{FrancoRubio2025MPUGauging}, for the exact
    once-blocked CZX tensor $B$ with source matrices
    $(M_1(B))_{(i,\beta),(\alpha,j)}=B^{i,j}_{\alpha\beta}$ and
    $(M_2(B))_{(\alpha,i),(j,\beta)}=B^{i,j}_{\alpha\beta}$, the factors of
    Definition~\ref{def:mpug_czx_displayed_factors} satisfy
    \begin{align}
        X_1Y_1 & =M_1(B), \\
        X_2Y_2 & =M_2(B).
    \end{align}
\end{lemma}

\begin{proof}\leanok
    \uses{def:mpug_czx_displayed_factors, def:mpug_czx_dagger_gauge,
        thm:mpug_czx_exact_tensor_unitary,
        def:mpu_source_matrix_one, def:mpu_source_matrix_two}
    Expanding each matrix product over its four-dimensional factor index,
    substituting the displayed formulas and the coordinate formula from
    Lemma~\ref{thm:mpug_czx_exact_tensor_unitary}, and evaluating the
    binary bond indices and the four physical values proves both entry
    identities.
\end{proof}

\begin{lemma}[Four Gram matrices of the displayed CZX factors]
    \label{thm:mpug_czx_displayed_grams}
    \lean{MPOTensor.CZX.displayed_grams}
    \leanok
    \uses{def:mpug_czx_displayed_factors}
    With the factors of two retaining the unnormalized Hadamard weights of
    \cite[lines~4559--4659]{FrancoRubio2025MPUGauging}, the four explicit
    factors of Definition~\ref{def:mpug_czx_displayed_factors} satisfy
    \begin{align}
        X_1^\dagger X_1 & =2I_4,
        \label{eq:mpug_czx_displayed_gram_x1} \\
        X_2^\dagger X_2 & =I_4,
        \label{eq:mpug_czx_displayed_gram_x2} \\
        Y_1Y_1^\dagger  & =I_4,               \\
        Y_2Y_2^\dagger  & =2I_4.
    \end{align}
\end{lemma}

\begin{proof}\leanok
    \uses{def:mpug_czx_displayed_factors, thm:mpug_czx_dagger_gauge_laws}
    The empty factors are Pauli-gauge contractions of the filled factors,
    $X_1=(I_4\otimes Y)Y_2^\dagger$ and $Y_1=X_2^\dagger(Y\otimes I_4)$.
    By Lemma~\ref{thm:mpug_czx_dagger_gauge_laws}, $Y$ is unitary, hence so
    are $I_4\otimes Y$ and $Y\otimes I_4$; composing a unitary matrix with
    $Y_2^\dagger$, respectively with $X_2^\dagger$, leaves the Gram matrix
    unchanged, so $X_1^\dagger X_1=Y_2Y_2^\dagger$ and
    $Y_1Y_1^\dagger=X_2^\dagger X_2$. Expanding the two filled
    Gram matrices $X_2^\dagger X_2$ and $Y_2Y_2^\dagger$ over their
    product-index spaces and substituting the factor coordinates at each
    pair of factor indices gives the two displayed scalar identities.
\end{proof}

\begin{definition}[Right inverses of the displayed CZX factors]
    \label{def:mpug_czx_displayed_right_inverses}
    \lean{MPOTensor.CZX.displayedZ₁, MPOTensor.CZX.displayedZ₂}
    \leanok
    \uses{def:mpug_czx_displayed_factors}
    For the factors of Definition~\ref{def:mpug_czx_displayed_factors}, the
    right inverses associated with the explicit decompositions in
    \cite[lines~4559--4659]{FrancoRubio2025MPUGauging} are
    \begin{align}
        Z_1 & =Y_1^\dagger,         \\
        Z_2 & =\tfrac12Y_2^\dagger.
    \end{align}
\end{definition}

\begin{lemma}[Right inverses and weighted normalization]
    \label{thm:mpug_czx_displayed_normalizations}
    \lean{MPOTensor.CZX.displayed_normalizations}
    \leanok
    \uses{def:mpug_czx_displayed_factors,
        def:mpug_czx_displayed_right_inverses, def:mpu_source_weight}
    Set $\rho=I_2/2$. The displayed factors and right inverses satisfy
    \begin{align}
        Y_1Z_1                         & =I_4,
        \label{eq:mpug_czx_displayed_right_inverse_y1} \\
        Y_2Z_2                         & =I_4,
        \label{eq:mpug_czx_displayed_right_inverse_y2} \\
        X_1^\dagger(I_4\otimes\rho)X_1 & =I_4.
    \end{align}
\end{lemma}

\begin{proof}\leanok
    \uses{thm:mpug_czx_displayed_grams,
        def:mpug_czx_displayed_right_inverses, def:mpu_source_weight}
    The two inverse identities follow directly from
    Lemma~\ref{thm:mpug_czx_displayed_grams}. Since
    $I_4\otimes\rho=I_8/2$, the weighted identity is
    $\tfrac12X_1^\dagger X_1=I_4$.
\end{proof}

\begin{lemma}[Both CZX source-cut ranks are four]
    \label{thm:mpug_czx_displayed_ranks}
    \lean{MPOTensor.CZX.displayed_ranks}
    \leanok
    \uses{def:mpug_czx_exact_tensor, def:mpu_source_matrix_one,
        def:mpu_source_matrix_two}
    For the exact once-blocked CZX tensor $B$,
    \begin{align}
        d_r=\operatorname{rank}M_1(B) & =4, \\
        d_l=\operatorname{rank}M_2(B) & =4.
    \end{align}
    Thus the four-dimensional legs in
    \cite[lines~4559--4659]{FrancoRubio2025MPUGauging} have the
    source-cut dimensions.
\end{lemma}

\begin{proof}\leanok
    \uses{thm:mpug_czx_displayed_cuts, thm:mpug_czx_displayed_grams,
        thm:mpug_czx_displayed_normalizations}
    Each cut factors through a four-dimensional space, so its rank is at
    most four. Setting $L_1=X_1^\dagger/2$ for the first cut and
    $L_2=X_2^\dagger$ for the second, (\ref{eq:mpug_czx_displayed_gram_x1})
    and (\ref{eq:mpug_czx_displayed_gram_x2}) give $L_iX_i=I_4$, while
    (\ref{eq:mpug_czx_displayed_right_inverse_y1}) and
    (\ref{eq:mpug_czx_displayed_right_inverse_y2}) give $Y_iZ_i=I_4$, hence
    $L_i(X_iY_i)Z_i=I_4$. Multiplication cannot increase rank, so each cut
    has rank at least four.
\end{proof}

\begin{definition}[The displayed factors in source-rank coordinates]
    \label{def:mpug_czx_displayed_source_factors}
    \lean{MPOTensor.CZX.displayedRightEquiv, MPOTensor.CZX.displayedLeftEquiv,
        MPOTensor.CZX.displayedSourceFactors,
        MPOTensor.CZX.displayedSourceFactors_coordinates,
        MPOTensor.CZX.displayedSourceFactors_inverse_coordinates}
    \leanok
    \uses{def:mpug_czx_displayed_factors,
        def:mpug_czx_displayed_right_inverses, thm:mpug_czx_displayed_ranks,
        def:mpug_czx_exact_tensor, def:mpu_weighted_source_factors,
        def:mpug_czx_dagger_gauge}
    Let $d_r,d_l$ be the two source-cut ranks of the exact CZX tensor $B$.
    By Lemma~\ref{thm:mpug_czx_displayed_ranks}, both are four. Let
    \begin{align}
        e_r & :\{0,\ldots,d_r-1\}\longrightarrow\{0,1,2,3\}, \\
        e_l & :\{0,\ldots,d_l-1\}\longrightarrow\{0,1,2,3\}
    \end{align}
    be the bijections that preserve each integer index. Pulling back the
    four-dimensional factor index of $X_1,Y_1,Z_1$ along $e_r$ and that
    of $X_2,Y_2,Z_2$ along $e_l$ yields matrices
    $\widehat X_i,\widehat Y_i,\widehat Z_i$. Together with $\rho=I_2/2$,
    these give source factors for $B$ in the sense of
    Definition~\ref{def:mpu_weighted_source_factors}, satisfying both cut
    factorizations, the two left-factor normalizations, and both
    right-inverse identities.

    Explicitly, for $0\leq k,i<4$ and $\alpha,\beta\in\{0,1\}$, evaluating
    the pullbacks gives
    \begin{align}
        (\widehat X_1)_{(i,\alpha),e_r^{-1}(k)} & =(X_1)_{(i,\alpha),k}, \\
        (\widehat Y_1)_{e_r^{-1}(k),(\alpha,i)} & =(Y_1)_{k,(\alpha,i)}, \\
        (\widehat Z_1)_{(\alpha,i),e_r^{-1}(k)} & =(Z_1)_{(\alpha,i),k}, \\
        (\widehat X_2)_{(\alpha,i),e_l^{-1}(k)} & =(X_2)_{(\alpha,i),k}, \\
        (\widehat Y_2)_{e_l^{-1}(k),(i,\alpha)} & =(Y_2)_{k,(i,\alpha)}, \\
        (\widehat Z_2)_{(i,\alpha),e_l^{-1}(k)} & =(Z_2)_{(i,\alpha),k},
    \end{align}
    and substituting the definitions of $X_1,Y_1,Z_1,Z_2$ gives the
    Pauli and inverse coordinates
    \begin{align}
        (\widehat X_1)_{(i,\beta),e_r^{-1}(k)}
         & =\sum_{t=0}^1Y_{\beta t}
        \overline{(\widehat Y_2)_{e_l^{-1}(k),(i,t)}},                           \\
        (\widehat Y_1)_{e_r^{-1}(k),(\beta,i)}
         & =\sum_{t=0}^1\overline{(\widehat X_2)_{(t,i),e_l^{-1}(k)}}Y_{t\beta}, \\
        (\widehat Z_1)_{(\beta,i),e_r^{-1}(k)}
         & =\overline{(\widehat Y_1)_{e_r^{-1}(k),(\beta,i)}},                   \\
        (\widehat Z_2)_{(i,\beta),e_l^{-1}(k)}
         & =\tfrac12\overline{(\widehat Y_2)_{e_l^{-1}(k),(i,\beta)}}.
    \end{align}
    In particular, as in the first factor of
    \cite[lines~4611--4659]{FrancoRubio2025MPUGauging}, pulling back
    preserves the Pauli orientation $Y_{\beta t}$.
    % This is an explicit SourceFactors witness, not an equality with the
    % choice-selected compact SVD factors. It supplies no canonical-form-II
    % presentation, fusion tensors, local action tensors, or anomaly data.
\end{definition}

\begin{proof}\leanok
    \uses{thm:mpug_czx_displayed_cuts, thm:mpug_czx_displayed_grams,
        thm:mpug_czx_displayed_normalizations, def:mpug_czx_displayed_factors,
        def:mpug_czx_displayed_right_inverses}
    Pulling back the cut factorizations and inverse identities along the
    two index bijections preserves products, conjugate transposes, and
    identities, so the weighted first normalization and the ordinary
    second normalization hold in the source-rank coordinates.
    Evaluating the pullbacks, using $e_r(e_r^{-1}(k))=k$ and
    $e_l(e_l^{-1}(k))=k$, gives the coordinate identities; substituting
    the definitions of $X_1,Y_1,Z_1,Z_2$ and cancelling the inverse
    bijections gives the Pauli and inverse coordinates.
\end{proof}

\subsection{Normalization of the displayed CZX comparison}

\begin{definition}[Reflected candidates and the two displayed comparisons]
    \label{def:mpug_czx_displayed_comparison}
    \lean{MPOTensor.CZX.displayedReflectedX₁, MPOTensor.CZX.displayedReflectedY₁,
        MPOTensor.CZX.displayedRawComparison,
        MPOTensor.CZX.displayedWeightedComparison}
    \leanok
    \uses{def:mpug_czx_displayed_factors, def:mpug_czx_dagger_gauge,
        def:mpu_source_weight}
    Let $Y$ be the Pauli adjoint gauge of
    Definition~\ref{def:mpug_czx_dagger_gauge} and let
    $X_1,Y_1,X_2,Y_2$ be the displayed factors of
    Definition~\ref{def:mpug_czx_displayed_factors}, with $d=4$ and $D=2$.
    Put $A=I_4\otimes\overline Y$ and $B=Y^{\mathsf T}\otimes I_4$, and
    define the reflected candidates
    \begin{align}
        \widehat X_1 & =AY_2^\dagger,  \\
        \widehat Y_1 & =X_2^\dagger B.
    \end{align}
    Evaluating the candidates and comparison of
    \cite[lines~5432--5441]{FrancoRubio2025MPUGauging} on the explicit
    factors of lines~4559--4659 with $\rho=I_2/2$ defines the unweighted
    and the density-weighted comparisons
    \begin{align}
        K_{\mathrm{unw}} & =X_1^\dagger\widehat X_1,                 \\
        K_\rho           & =X_1^\dagger(I_4\otimes\rho)\widehat X_1.
    \end{align}
\end{definition}

\begin{lemma}[The reflected candidates are the negatives of the displayed factors]
    \label{thm:mpug_czx_reflected_neg}
    \lean{MPOTensor.CZX.displayedReflected_eq_neg}
    \leanok
    \uses{def:mpug_czx_displayed_comparison, def:mpug_czx_displayed_factors}
    For the maintained Pauli gauge,
    \begin{align}
        \widehat X_1 & =-X_1, \\
        \widehat Y_1 & =-Y_1.
    \end{align}
\end{lemma}
\begin{proof}\leanok
    \uses{thm:mpug_czx_dagger_gauge_laws, thm:mpug_czx_dagger_sign,
        def:mpug_czx_displayed_factors, def:mpug_czx_displayed_comparison}
    Lemma~\ref{thm:mpug_czx_dagger_sign} gives $\overline Y=-Y$, and
    $Y^{\mathsf T}=\overline{Y^\dagger}=\overline Y=-Y$ by
    Lemma~\ref{thm:mpug_czx_dagger_gauge_laws}. Expanding both entries
    of the displayed factors against these two gauges gives the sign.
\end{proof}

\begin{lemma}[Coordinates of the two comparisons and the right overlap]
    \label{thm:mpug_czx_comparison_coordinates}
    \lean{MPOTensor.CZX.displayedComparison_coordinates,
        MPOTensor.CZX.displayedRawComparison_gram,
        MPOTensor.CZX.displayedRawComparison_not_isUnitaryBetween,
        MPOTensor.CZX.displayedRawComparison_ne_rightOverlap}
    \leanok
    \uses{def:mpug_czx_displayed_comparison}
    On the displayed factors,
    \begin{align}
        K_{\mathrm{unw}}        & =-2I_4,
        \label{eq:mpug_czx_raw_comparison}      \\
        K_\rho                  & =-I_4,
        \label{eq:mpug_czx_weighted_comparison} \\
        Y_1\widehat Y_1^\dagger & =-I_4.
        \label{eq:mpug_czx_right_overlap}
    \end{align}
    Consequently $K_{\mathrm{unw}}^\dagger K_{\mathrm{unw}}=4I_4$, so
    $K_{\mathrm{unw}}$ is not unitary, and
    $K_{\mathrm{unw}}\neq Y_1\widehat Y_1^\dagger$: the unweighted left
    overlap and the right overlap of
    \cite[lines~5438--5441]{FrancoRubio2025MPUGauging} disagree on these
    factors.
\end{lemma}
\begin{proof}\leanok
    \uses{thm:mpug_czx_reflected_neg, thm:mpug_czx_displayed_grams,
        thm:mpug_czx_displayed_normalizations,
        def:mpug_czx_displayed_comparison}
    Substituting Lemma~\ref{thm:mpug_czx_reflected_neg} turns each
    overlap into the negative of a Gram matrix of the displayed factors.
    The unweighted Gram $X_1^\dagger X_1=2I_4$ of
    Lemma~\ref{thm:mpug_czx_displayed_grams} gives
    (\ref{eq:mpug_czx_raw_comparison}); the weighted normalization
    $X_1^\dagger(I_4\otimes\rho)X_1=I_4$ of
    Lemma~\ref{thm:mpug_czx_displayed_normalizations} gives
    (\ref{eq:mpug_czx_weighted_comparison}); and $Y_1Y_1^\dagger=I_4$ gives
    (\ref{eq:mpug_czx_right_overlap}). Squaring
    (\ref{eq:mpug_czx_raw_comparison}) gives the Gram matrix $4I_4$, whose
    diagonal entry $4$ differs from $1$, and comparing the diagonal entries
    of (\ref{eq:mpug_czx_raw_comparison}) and
    (\ref{eq:mpug_czx_right_overlap}) separates the two overlaps.
\end{proof}

\begin{lemma}[The weighted comparison is unitary and recovers both factors]
    \label{thm:mpug_czx_weighted_comparison_spec}
    \lean{MPOTensor.CZX.displayedWeightedComparison_spec}
    \leanok
    \uses{def:mpug_czx_displayed_comparison}
    The weighted comparison $K_\rho$ is unitary, and
    \begin{align}
        X_1 & =\widehat X_1K_\rho^\dagger, \\
        Y_1 & =K_\rho\widehat Y_1.
    \end{align}
    The weight is the one carried by the normalization of the first left
    factor. This calculation concerns the explicit displayed factors of
    Definition~\ref{def:mpug_czx_displayed_factors}, not choice-selected
    factors, so the unequal overlaps of
    Lemma~\ref{thm:mpug_czx_comparison_coordinates} diagnose the
    intermediate formula of
    \cite[lines~5438--5441]{FrancoRubio2025MPUGauging} on these factors
    rather than contradicting the proposition being proved there. The
    weighted comparison is the one retained in
    \cite{gap:fbc25_inverse_compatible_tilde_omission}; no gauge phase is
    changed.
\end{lemma}
\begin{proof}\leanok
    \uses{thm:mpug_czx_comparison_coordinates, thm:mpug_czx_reflected_neg}
    By (\ref{eq:mpug_czx_weighted_comparison}), $K_\rho=-I_4$, which is
    unitary. Lemma~\ref{thm:mpug_czx_reflected_neg} then gives
    $\widehat X_1K_\rho^\dagger=(-X_1)(-I_4)=X_1$ and
    $K_\rho\widehat Y_1=(-I_4)(-Y_1)=Y_1$.
\end{proof}

\subsection{The blocked GHZ state and its CZX invariance}

\begin{definition}[The GHZ tensor in the blocked CZX alphabet]
    \label{def:mpug_czx_blocked_ghz}
    \lean{MPOTensor.CZX.blockedGHZTensor}
    \leanok
    \discussion{7354}
    \uses{def:ghz_tensor, def:blocked_tensor}
    Let $A^0=\ketbra{0}{0}$ and $A^1=\ketbra{1}{1}$ be the GHZ tensor,
    and identify the blocked physical alphabet with $\{0,1,2,3\}$ by
    $\varphi(2a+b)=(a,b)$ for $a,b\in\{0,1\}$. As in
    \cite[lines~4660--4670]{FrancoRubio2025MPUGauging}, with the same
    physical-index order as the CZX tensor $B$, the once-blocked GHZ
    tensor is
    \begin{align}
        \widehat A^{2a+b} & =A^aA^b=(A^{[2]})^{\varphi(2a+b)}.
    \end{align}
\end{definition}

\begin{lemma}[The two blocked GHZ sectors]
    \label{thm:mpug_czx_blocked_ghz_sectors}
    \lean{MPOTensor.CZX.blockTensor_ghzSectorTensor_apply}
    \leanok
    \uses{def:mpug_ghz_blocked_cluster_tensor, def:blocked_tensor}
    For $x\in\{0,1\}$, let $A_x^a=\delta_{a,x}$ be the
    bond-one GHZ sector. Under $\varphi(2a+b)=(a,b)$ its two-site block is
    \begin{align}
        (A_x^{[2]})^{\varphi(i)} & =\delta_{i,3x},
        \qquad i\in\{0,1,2,3\}.
    \end{align}
    Thus, as in \cite[lines~4660--4670]{FrancoRubio2025MPUGauging}, the
    blocked sector vectors are $\ket{00}$ and $\ket{11}$.
\end{lemma}

\begin{proof}\leanok
    For $i=2a+b$, the block coefficient is
    $A_x^aA_x^b=\delta_{a,x}\delta_{b,x}$, which equals one exactly when
    $a=b=x$, or equivalently $i=3x$.
\end{proof}

\begin{lemma}[The matrix product vector of the blocked GHZ tensor]
    \label{thm:mpug_czx_blocked_ghz_mpv}
    \lean{MPOTensor.CZX.mpv_blockedGHZTensor}
    \leanok
    \uses{def:mpug_czx_blocked_ghz, def:mpv}
    For every $N\geq0$, the periodic matrix product vector of
    $\widehat A$ is
    \begin{align}
        \ket{\psi_N}
         & =\ket{0}^{\otimes N}+\ket{3}^{\otimes N}.
        \label{eq:mpug_czx_blocked_ghz_mpv}
    \end{align}
    In the original two-qubit alphabet the two summands are
    $\ket{00}^{\otimes N}$ and $\ket{11}^{\otimes N}$.
    At $N=0$ they are the same empty configuration, with total coefficient
    two.
\end{lemma}

\begin{proof}\leanok
    The blocked letters are diagonal, with
    $\widehat A^i=\operatorname{diag}(\delta_{i,0},\delta_{i,3})$.
    Hence, for $s\in\{0,1,2,3\}^N$,
    \begin{align}
        \operatorname{tr}\bigl(\widehat A^{s_1}\cdots
        \widehat A^{s_N}\bigr)
         & =\prod_{j=1}^N\delta_{s_j,0}
        +\prod_{j=1}^N\delta_{s_j,3}.
    \end{align}
    These are the coefficients of the two constant configurations.
\end{proof}

\begin{lemma}[CZX invariance of the blocked GHZ state]
    \label{thm:mpug_czx_blocked_ghz_invariant}
    \lean{MPOTensor.CZX.mpo_tensor_mulVec_blockedGHZ}
    \leanok
    \discussion{7354}
    \uses{def:mpug_czx_exact_tensor, def:mpug_czx_blocked_ghz, def:mpv}
    Let $B$ be the once-blocked CZX tensor, with both output $Z$ gates and
    the unnormalized Hadamard weights, and let $N>0$.
    Reindex the matrix product vector of the two-site-blocked GHZ tensor
    $A^{[2]}$ by the alphabet bijection $\varphi(2a+b)=(a,b)$:
    \begin{align}
        \psi_N(s) & =V^{(N)}(A^{[2]})_{(\varphi(s_1),\ldots,\varphi(s_N))},
        \qquad s\in\{0,1,2,3\}^N.
    \end{align}
    Then, as in the GHZ-invariance assertion of
    \cite[lines~4660--4670]{FrancoRubio2025MPUGauging}, the periodic CZX
    operator satisfies
    \begin{align}
        O_N(B)\ket{\psi_N} & =\ket{\psi_N}.
    \end{align}
    % Only the actual blocked GHZ MPV and its positive-length invariance
    % are asserted. Fusion, action tensors, and classification remain open.
\end{lemma}

\begin{proof}\leanok
    \uses{thm:mpug_czx_blocked_ghz_mpv, thm:mpug_czx_exact_tensor_unitary}
    Reindexing the alphabet gives the vector in
    (\ref{eq:mpug_czx_blocked_ghz_mpv}). The CZX operator complements
    every matter bit. Its phase on either constant configuration is one:
    the cyclic exponent is zero on $0^N$ and $4N$ on $3^N$.
    Consequently
    \begin{align}
        O_N(B)\ket{0}^{\otimes N} & =\ket{3}^{\otimes N}, &
        O_N(B)\ket{3}^{\otimes N} & =\ket{0}^{\otimes N}.
    \end{align}
    Their sum is invariant.
\end{proof}

\subsection{The printed CZX action pair and its empty-word boundary}

The action tensors displayed in
\cite[lines~4690--4694]{FrancoRubio2025MPUGauging} satisfy the interior
identity for nonempty words and the exterior identity with nonempty buffers.
Their empty contraction is one half, whereas the general interior identity
in \cite[lines~1841--1870]{FrancoRubio2025MPUGauging} includes length zero.
The distinction and a local replacement are recorded in
\cite{gap:fbc25_czx_action_empty_word}. The replacement below is a different
pair, not a scalar gauge of the printed pair or a substitution in the
unbuffered defect formulas.

\begin{definition}[The acted GHZ sectors]
    \label{def:mpug_czx_acted_sectors}
    \lean{MPOTensor.CZX.blockedSector, MPOTensor.CZX.actedSector,
        MPOTensor.CZX.actionTarget, MPOTensor.CZX.actionLetter,
        MPOTensor.CZX.actionMatrix}
    \leanok
    \uses{def:mpug_czx_exact_tensor, def:mpug_ghz_blocked_cluster_tensor,
        def:blocked_tensor}
    For $x\in\{0,1\}$, let $A_x$ be the two-site block of the
    bond-one GHZ sector, with physical indices identified by
    $\varphi(2a+b)=(a,b)$. Thus $A_x^i=\delta_{i,3x}$.
    Write $B^{i,j}$ for the CZX tensor and define
    \begin{align}
        B_x^i & =\sum_{j=0}^3 B^{i,j}A_x^j, \\
        k_x   & =3(1-x),                    \\
        C_x   & =B^{k_x,3x}.
        \label{eq:mpug_czx_acted_sectors}
    \end{align}
    As in the physical contraction and exchanged GHZ blocks of
    \cite[lines~4660--4694]{FrancoRubio2025MPUGauging}, the target sector
    is $A_{1-x}$.
\end{definition}

\begin{lemma}[Coordinates of the acted sectors]
    \label{thm:mpug_czx_acted_sector_coordinates}
    \lean{MPOTensor.CZX.actedSector_apply, MPOTensor.CZX.actionTarget_apply,
        MPOTensor.CZX.actionMatrix_coordinates}
    \leanok
    \uses{def:mpug_czx_acted_sectors}
    For $x\in\{0,1\}$ and $i\in\{0,1,2,3\}$,
    \begin{align}
        B_x^i     & =\delta_{i,k_x}C_x,                     \\
        A_{1-x}^i & =\delta_{i,k_x},                        \\
        C_0       & =\begin{pmatrix}0&0\\-1&1\end{pmatrix}, \\
        C_1       & =\begin{pmatrix}1&1\\0&0\end{pmatrix}.
        \label{eq:mpug_czx_action_matrices}
    \end{align}
\end{lemma}
\begin{proof}\leanok
    \uses{def:mpug_czx_acted_sectors, thm:mpug_czx_blocked_ghz_sectors,
        def:mpug_czx_exact_tensor}
    Since $A_x^j=\delta_{j,3x}$, the defining sum of
    Definition~\ref{def:mpug_czx_acted_sectors} collapses to
    $B_x^i=\sum_{j=0}^3 B^{i,j}A_x^j=B^{i,3x}$. The exact CZX tensor of
    Definition~\ref{def:mpug_czx_exact_tensor} vanishes off its support,
    giving
    \begin{align}
        B^{i,3x} & =\delta_{i,3(1-x)}B^{3(1-x),3x}=\delta_{i,k_x}C_x,
    \end{align}
    so $B_x^i=\delta_{i,k_x}C_x$; likewise $A_x^j=\delta_{j,3x}$ at
    $x\mapsto1-x$ gives $A_{1-x}^i=\delta_{i,k_x}$. Expanding the defining
    CZX matrices at $B^{k_x,3x}$ evaluates the two displayed matrices
    $C_0,C_1$.
\end{proof}

\begin{definition}[The printed action pair]
    \label{def:mpug_czx_printed_action_pair}
    \lean{MPOTensor.CZX.printedActionBra, MPOTensor.CZX.printedActionKet}
    \leanok
    For $x\in\{0,1\}$, denote the printed left and right action tensors
    by $p_x$ and $q_x$. With $|\pm\rangle=(|0\rangle\pm|1\rangle)/\sqrt2$,
    the formulas of \cite[lines~4690--4694]{FrancoRubio2025MPUGauging} read
    \begin{align}
        p_0 & =(0,-i),                                    \\
        p_1 & =(1,0),                                     \\
        q_0 & =\tfrac12\begin{pmatrix}-i\\i\end{pmatrix}, \\
        q_1 & =\tfrac12\begin{pmatrix}1\\1\end{pmatrix}.
        \label{eq:mpug_czx_printed_action_pair}
    \end{align}
    The normalization of the kets is not specified explicitly in that
    passage. These coordinates agree with the later defect evaluations;
    see \cite{gap:fbc25_czx_action_empty_word}.
\end{definition}

\begin{lemma}[Finite contractions of the printed action pair]
    \label{thm:mpug_czx_printed_action_contractions}
    \lean{MPOTensor.CZX.printedAction_finite_identities}
    \leanok
    \uses{def:mpug_czx_acted_sectors, def:mpug_czx_printed_action_pair}
    For each $x\in\{0,1\}$,
    \begin{align}
        C_x^2        & =C_x,      \\
        p_xq_x       & =\tfrac12, \\
        p_xC_xq_x    & =1,        \\
        C_xq_xp_xC_x & =C_x.
        \label{eq:mpug_czx_printed_action_contractions}
    \end{align}
\end{lemma}
\begin{proof}\leanok
    \uses{thm:mpug_czx_acted_sector_coordinates, def:mpug_czx_printed_action_pair}
    Multiplying the two-by-two matrices and the printed rows and columns
    yields the four identities.
\end{proof}

\begin{lemma}[Nonempty acted and target words]
    \label{thm:mpug_czx_action_words}
    \lean{MPOTensor.CZX.evalWord_actedSector, MPOTensor.CZX.evalWord_actionTarget}
    \leanok
    \uses{def:mpug_czx_acted_sectors}
    For a word $\mathbf a$ in $\{0,1,2,3\}$, write $B_x^{\mathbf a}$
    and $A_{1-x}^{\mathbf a}$ for ordered products of letters. If
    $|\mathbf a|>0$, then
    \begin{align}
        B_x^{\mathbf a}     & =
        \begin{cases}C_x&\text{if every letter is }k_x,\\0&\text{otherwise},\end{cases} \\
        A_{1-x}^{\mathbf a} & =
        \begin{cases}1&\text{if every letter is }k_x,\\0&\text{otherwise}.\end{cases}
        \label{eq:mpug_czx_action_words}
    \end{align}
\end{lemma}
\begin{proof}\leanok
    \uses{thm:mpug_czx_acted_sector_coordinates, thm:mpug_czx_printed_action_contractions}
    An unsupported letter makes the product zero. Otherwise every acted
    letter is $C_x$, whose positive powers equal $C_x$ by idempotency.
\end{proof}

\begin{lemma}[Printed interior and buffered exterior identities]
    \label{thm:mpug_czx_printed_action_words}
    \lean{MPOTensor.CZX.printedAction_interior, MPOTensor.CZX.printedAction_exterior}
    \leanok
    \uses{def:mpug_czx_acted_sectors, def:mpug_czx_printed_action_pair}
    Let $x\in\{0,1\}$. For every nonempty word $\mathbf c$,
    \begin{align}
        p_xB_x^{\mathbf c}q_x & =A_{1-x}^{\mathbf c}.
        \label{eq:mpug_czx_printed_interior}
    \end{align}
    If $\mathbf a,\mathbf b$ are also nonempty, then
    \begin{align}
        B_x^{\mathbf a}q_xA_{1-x}^{\mathbf c}p_xB_x^{\mathbf b}
         & =B_x^{\mathbf a\mathbf c\mathbf b}.
        \label{eq:mpug_czx_printed_exterior}
    \end{align}
    No equality of buffer lengths is needed: while the equal-length case
    matches the exterior diagram of
    \cite[lines~1788--1838]{FrancoRubio2025MPUGauging}, the interior
    equation covers only the nonempty part of that source assertion; see
    \cite{gap:fbc25_czx_action_empty_word}.
\end{lemma}
\begin{proof}\leanok
    \uses{thm:mpug_czx_action_words, thm:mpug_czx_printed_action_contractions}
    If an unsupported letter occurs, both sides vanish. Otherwise the
    interior is $p_xC_xq_x=1$ and the exterior is $C_xq_xp_xC_x=C_x$.
\end{proof}

\begin{lemma}[The printed pair is not an all-word reduction]
    \label{thm:mpug_czx_printed_action_empty}
    \lean{MPOTensor.CZX.printedAction_empty_ne, MPOTensor.CZX.printedAction_not_isReduction}
    \leanok
    \uses{def:mpug_czx_acted_sectors, def:mpug_czx_printed_action_pair,
        def:mps_rectangular_reduction}
    As recorded in the empty-word obstruction of
    \cite{gap:fbc25_czx_action_empty_word}, for each $x\in\{0,1\}$, the
    printed pair $(p_x,q_x)$ is not a reduction from $B_x$ to $A_{1-x}$.
    Indeed, if $\varepsilon$ is the empty word, then
    \begin{align}
        p_xB_x^{\varepsilon}q_x & =\tfrac12\ne1=A_{1-x}^{\varepsilon}.
        \label{eq:mpug_czx_printed_action_empty}
    \end{align}
\end{lemma}
\begin{proof}\leanok
    \uses{thm:mpug_czx_printed_action_contractions}
    Empty products are identities, whereas $p_xq_x=1/2$.
\end{proof}

\begin{definition}[A separately dressed left action tensor]
    \label{def:mpug_czx_dressed_action}
    \lean{MPOTensor.CZX.dressedActionBra}
    \leanok
    \uses{def:mpug_czx_acted_sectors, def:mpug_czx_printed_action_pair}
    As the local replacement of \cite{gap:fbc25_czx_action_empty_word}
    rather than the printed pair, for $x\in\{0,1\}$, define $V_x=p_xC_x$
    and retain $W_x=q_x$. In coordinates $V_0=(i,-i)$ and
    $V_1=(1,1)$, so neither is a scalar multiple of the corresponding
    printed row. No replacement in an unbuffered defect formula is asserted.
\end{definition}

\begin{lemma}[Dressing is absorbed by a nonempty buffer]
    \label{thm:mpug_czx_dressed_action_absorption}
    \lean{MPOTensor.CZX.dressedActionBra_absorption}
    \leanok
    \uses{def:mpug_czx_dressed_action}
    For $x\in\{0,1\}$ and every nonempty word $\mathbf b$,
    \begin{align}
        V_xB_x^{\mathbf b} & =p_xB_x^{\mathbf b}.
        \label{eq:mpug_czx_dressed_action_absorption}
    \end{align}
\end{lemma}
\begin{proof}\leanok
    \uses{thm:mpug_czx_action_words, thm:mpug_czx_printed_action_contractions,
        def:mpug_czx_dressed_action}
    The word evaluates to either zero or $C_x$, and in the latter case
    $C_x^2=C_x$ gives the equality.
\end{proof}

\begin{lemma}[An all-word reduction with exterior buffer length one]
    \label{thm:mpug_czx_dressed_action_reduction}
    \lean{MPOTensor.CZX.dressedAction_isReduction,
        MPOTensor.CZX.dressedAction_isReductionExteriorBufferLength}
    \leanok
    \uses{def:mpug_czx_dressed_action, def:mps_rectangular_reduction,
        def:mps_reduction_exterior_buffer_length}
    For each $x\in\{0,1\}$, $(V_x,W_x)$ is a reduction from $B_x$
    to $A_{1-x}$ with exterior buffer length one. The reduction identity
    holds for every word $\mathbf c$, including the empty word; the exterior
    identity holds for nonempty words $\mathbf a$, $\mathbf b$, and
    $\mathbf c$. Thus
    \begin{align}
        V_xW_x                & =1,                                 \\
        V_xB_x^{\mathbf c}W_x & =A_{1-x}^{\mathbf c},               \\
        B_x^{\mathbf a}W_xA_{1-x}^{\mathbf c}V_xB_x^{\mathbf b}
                              & =B_x^{\mathbf a\mathbf c\mathbf b}.
        \label{eq:mpug_czx_dressed_action_reduction}
    \end{align}
    This proves the all-word relation for the separately dressed pair,
    not the empty-word claim for the printed pair.
\end{lemma}
\begin{proof}\leanok
    \uses{thm:mpug_czx_printed_action_contractions,
        thm:mpug_czx_printed_action_words, thm:mpug_czx_dressed_action_absorption,
        def:mpug_czx_dressed_action}
    The empty contraction is $V_xW_x=p_xC_xq_x=1$. For a nonempty
    interior word, absorbing $C_x$ and applying the printed interior
    identity yields the second equation, while in the exterior identity,
    absorbing $C_x$ into the right buffer and applying
    (\ref{eq:mpug_czx_printed_exterior}) completes the proof.
\end{proof}

\subsection{The CZX \texorpdfstring{$g,g$}{g,g} fusion caps}

\begin{definition}[The CZX $g,g$ fusion caps]
    \label{def:mpug_czx_fusion_caps}
    \lean{MPOTensor.CZX.fusionV, MPOTensor.CZX.fusionW}
    \leanok
    As in the concrete caps of
    \cite[lines~4671--4689]{FrancoRubio2025MPUGauging} (the $g=h=g^{-1}$
    specialization of the general fusion-tensor display in
    \cite[lines~2000--2037]{FrancoRubio2025MPUGauging}), in the ordered bond
    coordinates $0,1,2,3$, define the row vector
    $F^<_{g,g}\colon\C^4\to\C$ and the column vector
    $F^>_{g,g}\colon\C\to\C^4$ by
    \begin{align}
        F^<_{g,g} & =(0,-i,i,0),                     \\
        F^>_{g,g} & =\tfrac12(0,i,-i,0)^{\mathsf T}.
    \end{align}
    % Only the concrete g,g slice of #7354 is covered here.
\end{definition}

\begin{lemma}[Product coordinates of the CZX fusion caps]
    \label{thm:mpug_czx_fusion_caps_apply}
    \lean{MPOTensor.CZX.fusion_caps_apply}
    \leanok
    \uses{def:mpug_czx_fusion_caps, def:mpug_czx_dagger_gauge}
    Let $F^<_{g,g},F^>_{g,g}$ be the caps of
    Definition~\ref{def:mpug_czx_fusion_caps}, and let $Y$ be the Pauli
    matrix of Definition~\ref{def:mpug_czx_dagger_gauge}. Identify the
    product-bond coordinate $(a,b)$ with $2a+b$. For all bond bits
    $a,b\in\{0,1\}$,
    \begin{align}
        F^<_{g,g}(a,b) & =Y_{a,b},
        \label{eq:mpug_czx_fusion_caps_apply} \\
        F^>_{g,g}(a,b) & =\tfrac12Y_{b,a}.
        \label{eq:mpug_czx_fusion_caps_apply_right}
    \end{align}
    Thus, as in the printed diagram of
    \cite[lines~4671--4689]{FrancoRubio2025MPUGauging}, both caps vanish at
    the diagonal bond coordinates and carry $Y$ in opposite orientations.
\end{lemma}
\begin{proof}\leanok
    \uses{def:mpug_czx_fusion_caps, def:mpug_czx_dagger_gauge}
    Evaluating both caps at $2a+b$ for the four pairs $(a,b)$ and comparing
    with the entries of $Y$ yields
    (\ref{eq:mpug_czx_fusion_caps_apply}) and
    (\ref{eq:mpug_czx_fusion_caps_apply_right}).
\end{proof}

\begin{lemma}[The CZX fusion reduction]
    \label{thm:mpug_czx_fusion_reduction}
    \lean{MPOTensor.CZX.fusion_isReduction,
        MPOTensor.CZX.fusion_isReductionExteriorBufferLength_one}
    \leanok
    \uses{def:mpug_czx_fusion_caps, def:mpdo_operator_product,
        def:mpug_czx_exact_tensor, threeMPU, def:mpo_to_mps,
        def:mps_rectangular_reduction,
        def:mps_reduction_exterior_buffer_length}
    Let $B$ be the doubled-index tensor of the product
    $\mathcal U_g\cdot\mathcal U_g$
    (Definition~\ref{def:mpdo_operator_product} applied to
    Definition~\ref{def:mpug_czx_exact_tensor} with itself, viewed
    through Definition~\ref{def:mpo_to_mps}), and let $I$ be the
    doubled-index bond-one identity tensor of
    Definition~\ref{threeMPU}. Then $(F^<_{g,g},F^>_{g,g})$ is a
    reduction from $B$ to $I$ with exterior buffer length one. For every
    word $\mathbf c$, including the empty word,
    \begin{align}
        F^<_{g,g}F^>_{g,g}              & =1,             \\
        F^<_{g,g}B^{\mathbf c}F^>_{g,g} & =I^{\mathbf c}.
    \end{align}
    For nonempty words $\mathbf p,\mathbf q,\mathbf c$, this gives the
    concrete $g,g$ instance of the fusion reduction
    \cite[lines~1405--1498]{FrancoRubio2025MPUGauging} with the fusion
    tensors of \cite[lines~2000--2037]{FrancoRubio2025MPUGauging}
    specialized to the CZX display
    \cite[lines~4671--4689]{FrancoRubio2025MPUGauging}:
    \begin{align}
        B^{\mathbf p}F^>_{g,g}I^{\mathbf c}F^<_{g,g}B^{\mathbf q}
         & =B^{\mathbf p\mathbf c\mathbf q}.
        \label{eq:mpug_czx_fusion_reduction}
    \end{align}
\end{lemma}
\begin{proof}\leanok
    \uses{def:mpug_czx_fusion_caps, def:mpug_czx_exact_tensor,
        def:mpdo_operator_product, def:mpo_to_mps, threeMPU}
    With $V=F^<_{g,g}$ and $W=F^>_{g,g}$, substituting the cap and
    tensor coordinates gives the three local identities
    \begin{align}
        VW       & =1,
        \label{eq:mpug_czx_fusion_cap_product} \\
        VB^iW    & =I^i,
        \label{eq:mpug_czx_fusion_letter}      \\
        B^iWVB^j & =B^iB^j.
        \label{eq:mpug_czx_fusion_letter_insertion}
    \end{align}
    The all-word identity $VB^{\mathbf c}W=I^{\mathbf c}$ follows by
    induction on $\mathbf c$, removing its first letter. The empty word
    is (\ref{eq:mpug_czx_fusion_cap_product}) and a single letter is
    (\ref{eq:mpug_czx_fusion_letter}). For $\mathbf c=ij\mathbf c'$,
    inserting $WV$ after the first letter by
    (\ref{eq:mpug_czx_fusion_letter_insertion}), splitting the caps off both
    factors, and applying (\ref{eq:mpug_czx_fusion_letter}) together with
    the induction hypothesis for $j\mathbf c'$ gives
    \begin{align}
        VB^{ij\mathbf c'}W
         & =(VB^iW)(VB^{j\mathbf c'}W)
        =I^iI^{j\mathbf c'}=I^{ij\mathbf c'}.
    \end{align}
    For nonempty words $\mathbf p,\mathbf q$, a second induction on the left
    word $\mathbf p$, again removing its first letter, gives
    \begin{align}
        B^{\mathbf p}WVB^{\mathbf q} & =B^{\mathbf p\mathbf q}.
        \label{eq:mpug_czx_fusion_word_insertion}
    \end{align}
    For the singleton case $\mathbf p=i$ and
    $\mathbf q=j\mathbf q'$,
    (\ref{eq:mpug_czx_fusion_letter_insertion}) gives
    \begin{align}
        B^iWVB^{j\mathbf q'}
         & =(B^iWVB^j)B^{\mathbf q'}
        =B^iB^jB^{\mathbf q'}=B^{i j\mathbf q'}.
    \end{align}
    If $\mathbf p=i\mathbf p'$ with $\mathbf p'$ nonempty, the induction
    hypothesis applied to the tail gives
    \begin{align}
        B^{i\mathbf p'}WVB^{\mathbf q}
         & =B^i(B^{\mathbf p'}WVB^{\mathbf q})
        =B^iB^{\mathbf p'\mathbf q}=B^{i\mathbf p'\mathbf q}.
    \end{align}
    For nonempty buffers and middle word, substituting
    $I^{\mathbf c}=VB^{\mathbf c}W$ and applying
    (\ref{eq:mpug_czx_fusion_word_insertion}) twice yields
    \begin{align}
        B^{\mathbf p}WI^{\mathbf c}VB^{\mathbf q}
         & =B^{\mathbf p}WVB^{\mathbf c}WVB^{\mathbf q}
        =B^{\mathbf p\mathbf c\mathbf q}.
    \end{align}
\end{proof}

\subsection{Printed fusion--action coordinate contractions}

The contractions below evaluate the printed action pair and right fusion
cap of \cite[lines~4671--4694]{FrancoRubio2025MPUGauging} in fixed product-bond
order. They are compared with the fusion--action relation of
\cite[lines~1875--1913]{FrancoRubio2025MPUGauging}; the next subsection
identifies these coefficients with the $L$-symbols of two generators, up to the
orientation of the caps \cite{gap:fbc25_czx_fusion_coordinate_compatibility}.
No cap, action tensor, block label, or physical gate is changed.

\begin{definition}[Sequential action covector]
    \label{def:mpug_czx_sequential_action_bra}
    \lean{MPOTensor.CZX.sequentialActionBra}
    \leanok
    \uses{def:mpug_czx_printed_action_pair}
    For $x\in\{0,1\}$, let $p_x$ be the printed action row of
    Definition~\ref{def:mpug_czx_printed_action_pair}. Define the row
    $P_x\colon\C^4\to\C$ by the following coordinates, for
    $\alpha,\gamma\in\{0,1\}$:
    \begin{align}
        P_x                    & =p_{1-x}\otimes p_x,         &
        (P_x)_{2\alpha+\gamma} & =p_{1-x}(\alpha)p_x(\gamma).
    \end{align}
    The first action factor occupies the first product-bond coordinate.
\end{definition}

\begin{definition}[Sequential action vector]
    \label{def:mpug_czx_sequential_action_ket}
    \lean{MPOTensor.CZX.sequentialActionKet}
    \leanok
    \uses{def:mpug_czx_printed_action_pair}
    For $x\in\{0,1\}$, let $q_x$ be the printed action column of
    Definition~\ref{def:mpug_czx_printed_action_pair}. Define the column
    $Q_x\colon\C\to\C^4$ by the following coordinates, for
    $\alpha,\gamma\in\{0,1\}$:
    \begin{align}
        Q_x                    & =q_{1-x}\otimes q_x,         &
        (Q_x)_{2\alpha+\gamma} & =q_{1-x}(\alpha)q_x(\gamma).
    \end{align}
    Rows contract with columns by ordinary matrix multiplication, with no
    additional complex conjugation.
\end{definition}

\begin{definition}[Supported product letter]
    \label{def:mpug_czx_fusion_action_letter}
    \lean{MPOTensor.CZX.fusionActionLetter}
    \leanok
    \uses{def:mpug_czx_exact_tensor, def:mpdo_operator_product}
    Let $T$ be the exact blocked CZX tensor of
    Definition~\ref{def:mpug_czx_exact_tensor}. For $x\in\{0,1\}$, define
    $B_x\in\operatorname{Mat}_4(\C)$ to be the physical $(3x,3x)$ letter
    of $T\mathbin{\cdot}T$. Explicitly, for
    $\alpha,\gamma,\beta,\delta\in\{0,1\}$,
    \begin{align}
        (B_x)_{2\alpha+\gamma,2\beta+\delta}
         & =\sum_{j=0}^3 T^{3x,j}_{\alpha\beta}T^{j,3x}_{\gamma\delta}.
    \end{align}
    The physical index $3x$ encodes the two bits $(x,x)$.
\end{definition}

\begin{lemma}[Coordinates of the supported product letters]
    \label{thm:mpug_czx_supported_letter_coordinates}
    \lean{MPOTensor.CZX.mulTensor_diagonal_letter_coordinates}
    \leanok
    \uses{def:mpug_czx_fusion_action_letter}
    The two supported product letters have coordinates
    \begin{align}
        B_0 & =\begin{pmatrix}
                   0 & 0 & 0 & 0 \\-1&1&-1&1\\0&0&0&0\\0&0&0&0
               \end{pmatrix}, \\
        B_1 & =\begin{pmatrix}
                   0 & 0 & 0 & 0 \\0&0&0&0\\-1&-1&1&1\\0&0&0&0
               \end{pmatrix}.
        \label{eq:mpug_czx_supported_letters}
    \end{align}
\end{lemma}
\begin{proof}\leanok
    \uses{def:mpug_czx_exact_tensor, def:mpdo_operator_product}
    The letter of $T\mathbin{\cdot}T$ at a physical pair $(i,i')$ vanishes
    unless $i=i'$, because the blocked tensor complements the physical index
    it carries. At $i=i'=3x$ the summation index of
    Definition~\ref{def:mpug_czx_fusion_action_letter} retains the single
    term $j=3(1-x)$, whose two factors force the left virtual pair
    $(\alpha,\gamma)=(x,1-x)$, the row of index $x+1$. On that row the
    surviving sign is $(-1)^{1+(1-x)\delta+x\beta}$, which is the row
    displayed above.
\end{proof}

\begin{definition}[Printed fusion--action coefficient]
    \label{def:mpug_czx_fusion_action_coefficient}
    \lean{MPOTensor.CZX.fusionActionCoefficient}
    \leanok
    \uses{def:mpug_czx_sequential_action_bra,
        def:mpug_czx_fusion_action_letter, def:mpug_czx_fusion_caps}
    For $x\in\{0,1\}$, let $P_x$ and $B_x$ be the sequential action row
    and supported product letter, and put
    $W=F^>_{g,g}=\tfrac12(0,i,-i,0)^{\mathsf T}$, in the product-bond
    order of Definition~\ref{def:mpug_czx_fusion_caps}. Define
    \begin{align}
        c_x & =P_xB_xW\in\C.
    \end{align}
    The one-by-one matrix on the right is identified with its sole entry.
\end{definition}

\begin{lemma}[Values of the printed coefficient]
    \label{thm:mpug_czx_fusion_action_coefficient_coordinates}
    \lean{MPOTensor.CZX.fusionActionCoefficient_coordinates}
    \leanok
    \uses{def:mpug_czx_fusion_action_coefficient}
    The printed coefficients satisfy $c_0=1$ and $c_1=-1$.
\end{lemma}
\begin{proof}\leanok
    \uses{thm:mpug_czx_supported_letter_coordinates,
        def:mpug_czx_sequential_action_bra, def:mpug_czx_printed_action_pair,
        def:mpug_czx_fusion_caps}
    The printed rows are $P_0=(0,-i,0,0)$ and $P_1=(0,0,-i,0)$, and the
    column is $W=\tfrac12(0,i,-i,0)^{\mathsf T}$. With $e_0,\dots,e_3$ the
    standard basis of the product bond, the single nonzero row of each
    supported letter of (\ref{eq:mpug_czx_supported_letters}) contracts with
    $W$ as
    \begin{align}
        B_0W & =\tfrac12(i+i)e_1=ie_1,   \\
        B_1W & =\tfrac12(-i-i)e_2=-ie_2,
    \end{align}
    so the two finite matrix products are
    \begin{align}
        c_0 & =P_0B_0W=(0,-i,0,0)\,ie_1=1,     \\
        c_1 & =P_1B_1W=(0,0,-i,0)\,(-ie_2)=-1.
    \end{align}
\end{proof}

\begin{lemma}[Normalization, proportionality, and idempotence]
    \label{thm:mpug_czx_fusion_action_finite_identities}
    \lean{MPOTensor.CZX.fusionAction_finite_identities}
    \leanok
    \uses{def:mpug_czx_sequential_action_bra,
        def:mpug_czx_sequential_action_ket, def:mpug_czx_fusion_action_letter,
        def:mpug_czx_fusion_action_coefficient, def:mpug_czx_fusion_caps}
    For each $x\in\{0,1\}$, the printed rows, columns, and supported
    product letter satisfy
    \begin{align}
        P_xB_xQ_x & =1, & B_xW & =c_xB_xQ_x, & B_x^2 & =B_x.
    \end{align}
\end{lemma}
\begin{proof}\leanok
    \uses{thm:mpug_czx_supported_letter_coordinates,
        def:mpug_czx_fusion_action_letter,
        def:mpug_czx_sequential_action_bra, def:mpug_czx_sequential_action_ket,
        def:mpug_czx_printed_action_pair, def:mpug_czx_fusion_caps,
        def:mpug_czx_fusion_action_coefficient}
    Substituting the printed action columns gives
    $Q_0=\tfrac14(-i,i,-i,i)^{\mathsf T}$ and
    $Q_1=\tfrac14(-i,-i,i,i)^{\mathsf T}$. With $e_0,\dots,e_3$ the standard
    basis of the product bond, the two supported letters of
    (\ref{eq:mpug_czx_supported_letters}) contract with these columns and
    with the cap $W$ as
    \begin{align}
        B_0Q_0 & =ie_1,  \\
        B_0W   & =ie_1,  \\
        B_1Q_1 & =ie_2,  \\
        B_1W   & =-ie_2.
        \label{eq:mpug_czx_fusion_action_contractions}
    \end{align}
    The printed rows are $P_0=(0,-i,0,0)$ and $P_1=(0,0,-i,0)$, so applying
    $P_x$ to (\ref{eq:mpug_czx_fusion_action_contractions}) gives
    $P_xB_xQ_x=1$ and $P_xB_xW=c_x$, with $c_0=1$ and $c_1=-1$. Comparing
    the two right-hand sides for each $x$ gives $B_xW=c_xB_xQ_x$. Each
    matrix of (\ref{eq:mpug_czx_supported_letters}) has a single nonzero row,
    whose entry in the supporting column is $1$, so $B_x^2=B_x$.
\end{proof}

\begin{lemma}[Positive repetitions of the supported letter]
    \label{thm:mpug_czx_fusion_action_positive_power}
    \lean{MPOTensor.CZX.fusionAction_positive_power}
    \leanok
    \uses{def:mpug_czx_sequential_action_bra,
        def:mpug_czx_sequential_action_ket, def:mpug_czx_fusion_action_letter,
        def:mpug_czx_fusion_action_coefficient, def:mpug_czx_fusion_caps}
    For every $x\in\{0,1\}$ and every integer $n\geq0$,
    \begin{align}
        B_x^{n+1}       & =B_x, & P_xB_x^{n+1}W & =c_x,             \\
        P_xB_x^{n+1}Q_x & =1,   & B_x^{n+1}W    & =c_xB_x^{n+1}Q_x.
    \end{align}
    These statements concern repetitions of this single supported letter,
    not arbitrary words or an identification with the source $L$-symbol.
\end{lemma}
\begin{proof}\leanok
    \uses{thm:mpug_czx_fusion_action_finite_identities,
        def:mpug_czx_fusion_action_coefficient}
    By induction on $n$, the first power is $B_x$, and multiplication by
    $B_x$ together with idempotence proves the successor step.
    Substituting $B_x^{n+1}=B_x$ into the remaining three identities and
    applying the definition of $c_x$ along with the normalization and
    proportionality of
    Lemma~\ref{thm:mpug_czx_fusion_action_finite_identities} completes
    the proof.
\end{proof}
\subsection{The CZX anomaly and its \texorpdfstring{$L$}{L}-symbols}

The once-blocked CZX operators, the fusion caps of
Definition~\ref{def:mpug_czx_fusion_caps}, and the dressed action tensors of
Definition~\ref{def:mpug_czx_dressed_action} supply all data of the anomaly and
of the $L$-symbols on the two blocked GHZ sectors. The computation below
identifies the coefficients of
Lemma~\ref{thm:mpug_czx_fusion_action_coefficient_coordinates} with the
$L$-symbols of two generators. With the caps read in the product-bond order of
the stacked tensor, $L^0_{g,g}=1$ and $L^1_{g,g}=-1$; the source states
$L^0_{g,g}=-1$ \cite[line~4699]{FrancoRubio2025MPUGauging}. The two values are
related by the fusion gauge $\beta_{g,g}=-1$, which transposes both caps, since
$Y^{\mathsf T}=-Y$, and leaves the anomaly unchanged
\cite{gap:fbc25_czx_fusion_coordinate_compatibility}. The gauge-invariant ratio
$L^0_{g,g}/L^1_{g,g}=-1$ agrees with the source.

\begin{lemma}[The CZX representation of \texorpdfstring{$\Z_2$}{Z2}]
    \label{thm:mpug_czx_z2_representation}
    \lean{MPOTensor.CZX.mpo_tensor_mul_self,
        MPOTensor.CZX.family,
        MPOTensor.CZX.family_isNormalRepresentation}
    \leanok
    \discussion{7354}
    \uses{def:mpug_czx_exact_tensor, thm:mpug_czx_exact_tensor_injective,
        def:mpug_anomaly_three_cochain, threeMPU}
    Let $\mathcal U_g$ be the periodic operator of the once-blocked CZX tensor
    of Definition~\ref{def:mpug_czx_exact_tensor} on $N\geq1$ blocked sites.
    Then $\mathcal U_g^2=\Id$. Consequently the bond-one identity tensor for
    $e$ and the CZX tensor for $g$ represent $\Z_2=\{e,g\}$ by normal
    tensors, with $\mathcal U_e\mathcal U_h=\mathcal U_h\mathcal U_e=\mathcal U_h$
    and $\mathcal U_g\mathcal U_g=\mathcal U_e$ on every nonempty chain
    \cite[lines~4503--4547]{FrancoRubio2025MPUGauging}.
\end{lemma}

\begin{proof}\leanok
    \uses{thm:mpug_czx_exact_tensor_unitary, thm:mpug_czx_exact_tensor_injective}
    The periodic operator is the monomial operator of the complement of every
    matter bit, with a sign given by a cyclic sum of edge exponents. The
    complement is an involution, and the edge exponents of a configuration and
    of its complement differ by the sum of the leading bits of the two sites
    of the edge, modulo two. Summed around the ring these differences cancel,
    so the two signs multiply to one and $\mathcal U_g^2=\Id$. The identity
    tensor is normal, and the CZX tensor is injective, hence normal.
\end{proof}

\begin{lemma}[The displayed CZX anomaly]
    \label{thm:mpug_czx_anomaly}
    \lean{MPOTensor.CZX.fusionData,
        MPOTensor.CZX.tensorInt,
        MPOTensor.CZX.tensor_eq_complexOfInt,
        MPOTensor.CZX.fusionVInt,
        MPOTensor.CZX.fusionV_eq_smul,
        MPOTensor.CZX.mulTensor_tensor_eq_complexOfInt,
        MPOTensor.CZX.cubeInt,
        MPOTensor.CZX.tensorIntMPS,
        MPOTensor.CZX.fusionData_isAssociator,
        MPOTensor.CZX.fusionData_omega_gen_gen_gen,
        MPOTensor.CZX.kronId_fusionV_one,
        MPOTensor.CZX.fusionData_omega}
    \leanok
    \discussion{7354}
    \uses{thm:mpug_czx_z2_representation, thm:mpug_czx_fusion_reduction,
        def:mpug_anomaly_three_cochain, thm:mpug_anomaly_three_cochain_spec}
    Take the fusion tensors to be trivial whenever a factor is the identity
    and to be the caps $(F^<_{g,g},F^>_{g,g})$ of
    Definition~\ref{def:mpug_czx_fusion_caps} for two generators. Then, as
    displayed in \cite[lines~4695--4698]{FrancoRubio2025MPUGauging},
    \begin{align}
        \omega(h_1,h_2,h_3)
         & =(-1)^{\delta_{h_1,g}\delta_{h_2,g}\delta_{h_3,g}}.
        \label{eq:mpug_czx_anomaly}
    \end{align}
\end{lemma}

\begin{proof}\leanok
    \uses{thm:mpug_czx_fusion_reduction, thm:mpug_anomaly_three_cochain_spec}
    When one of $h_1,h_2,h_3$ is $e$, both fusion trees are built from the
    trivial fusion tensors and the cap of the remaining pair, and they agree,
    so $\omega(h_1,h_2,h_3)=1$. For three generators the two trees are
    $F^<_{g,g}\otimes\Id$ and $\Id\otimes F^<_{g,g}$. Both caps are $-i$ times
    an integer row, and the letters of the threefold product are images of
    integer matrices. On every letter $R^a$ of the threefold product the
    integer identity $(F\otimes\Id)R^a=-(\Id\otimes F)R^a$ holds, and both
    trees intertwine two letters with the target. Induction on the word gives
    the dressed identity with scalar $-1$, so $\omega(g,g,g)=-1$.
\end{proof}

\begin{lemma}[The CZX anomaly class is nonzero]
    \label{thm:mpug_czx_anomaly_class}
    \lean{MPOTensor.CZX.cyclicInvariant_omega,
        MPOTensor.CZX.not_isTrivialGaugeClass_omega,
        MPOTensor.CZX.family_anomalyClass_ne_zero}
    \leanok
    \discussion{7354}
    \uses{thm:mpug_czx_anomaly, def:mpug_anomaly_class,
        cor:mpug_anomaly_class}
    For every choice of fusion tensors of the CZX representation,
    $\omega(g,e,g)\,\omega(g,g,g)=-1$. Hence the anomaly three-cocycle is not
    cohomologous to the constant cochain one, and the anomaly class in
    $H^3(\Z_2,\C^\times)$ is nonzero.
\end{lemma}

\begin{proof}\leanok
    \uses{thm:mpug_czx_anomaly, cor:mpug_anomaly_class}
    The product $\omega(g,e,g)\,\omega(g,g,g)$ is invariant under the scalar
    gauge action of the fusion tensors because $g^2=e$, and every two choices
    of fusion tensors differ by that action. For the fusion tensors of
    Lemma~\ref{thm:mpug_czx_anomaly} it equals $1\cdot(-1)$. A trivial gauge
    class would make it one.
\end{proof}

\begin{lemma}[The two blocked GHZ sectors are permuted]
    \label{thm:mpug_czx_ghz_blocks}
    \lean{MPOTensor.CZX.blockedGHZTensor_apply,
        MPOTensor.CZX.cyclicPhase_ghz_constant,
        MPOTensor.CZX.block,
        MPOTensor.CZX.block_apply,
        MPOTensor.CZX.block_isNormal,
        MPOTensor.CZX.carriesMPV}
    \leanok
    \discussion{7354}
    \uses{thm:mpug_czx_blocked_ghz_sectors, thm:mpug_czx_blocked_ghz_invariant,
        thm:mpug_czx_z2_representation, def:mpug_permuted_block_l_symbol}
    Index the two bond-one blocked GHZ sectors $\ket{00}$ and $\ket{11}$ of
    Lemma~\ref{thm:mpug_czx_blocked_ghz_sectors} by $\Z_2$ acting on itself.
    Both are normal, and for every $h,x\in\Z_2$ and $N\geq1$ the CZX
    representation carries the periodic vector of the sector $x$ to that of
    the sector $hx$, with phase one.
\end{lemma}

\begin{proof}\leanok
    \uses{thm:mpug_czx_blocked_ghz_invariant}
    One letter of each sector is the scalar one, so the sector is normal. The
    identity tensor acts trivially. The CZX operator is monomial; it carries
    the constant configuration $x$ to the constant configuration $1-x$, and
    the cyclic sign of a constant configuration is one.
\end{proof}

\begin{definition}[Action tensors on the two blocked GHZ sectors]
    \label{def:mpug_czx_block_action_data}
    \lean{MPOTensor.CZX.actTensor_tensor_blockedSector,
        MPOTensor.CZX.actionV,
        MPOTensor.CZX.actionW,
        MPOTensor.CZX.blockActionData}
    \leanok
    \discussion{7354}
    \uses{def:mpug_czx_acted_sectors, def:mpug_czx_dressed_action,
        thm:mpug_czx_dressed_action_reduction, thm:mpug_czx_ghz_blocks,
        def:mpug_permuted_block_l_symbol, thm:mpug_permuted_block_l_symbol_spec}
    For the identity the action tensors are the identification of the bond
    spaces. For the generator, the CZX tensor acting on the sector $x$ is the
    acted sector of Definition~\ref{def:mpug_czx_acted_sectors}, and the action
    tensors are the dressed pair $(V_x,W_x)$ of
    Definition~\ref{def:mpug_czx_dressed_action}, which reduces it onto the
    sector $1-x$ and agrees with the printed action covector of
    \cite[lines~4690--4694]{FrancoRubio2025MPUGauging} against every nonempty
    word.
\end{definition}

\begin{lemma}[The CZX \texorpdfstring{$L$}{L}-symbols]
    \label{thm:mpug_czx_l_symbols}
    \lean{MPOTensor.CZX.blockInt,
        MPOTensor.CZX.pairActionInt,
        MPOTensor.CZX.actTensor_mulTensor_blockedSector,
        MPOTensor.CZX.successiveActionInt,
        MPOTensor.CZX.actV_gen_gen,
        MPOTensor.CZX.fuseV_gen_gen,
        MPOTensor.CZX.isDressedProportional_actV_fuseV_gen_gen,
        MPOTensor.CZX.lSymbol_one_one,
        MPOTensor.CZX.blockActionData_isCompatible,
        MPOTensor.CZX.blockActionData_lSymbol}
    \leanok
    \discussion{7354}
    \uses{def:mpug_czx_block_action_data, thm:mpug_czx_anomaly,
        thm:mpug_permuted_block_l_symbol_spec,
        thm:mpug_permuted_block_l_compatibility,
        thm:mpug_czx_fusion_action_coefficient_coordinates}
    With the fusion tensors of Lemma~\ref{thm:mpug_czx_anomaly} and the action
    tensors of Definition~\ref{def:mpug_czx_block_action_data}, the
    $L$-symbols on the two blocked GHZ sectors are
    \begin{align}
        L^x_{h,k} & =
        \begin{cases}
            (-1)^x & \text{if } h=k=g, \\
            1      & \text{otherwise},
        \end{cases}
        \label{eq:mpug_czx_l_symbols}
    \end{align}
    for $x\in\{0,1\}$. The source states $L^0_{g,g}=-1$
    \cite[line~4699]{FrancoRubio2025MPUGauging}; with both caps transposed,
    which is the fusion gauge $\beta_{g,g}=-1$, the two values $L^x_{g,g}$
    change sign and the printed value is obtained
    \cite{gap:fbc25_czx_fusion_coordinate_compatibility}.
\end{lemma}

\begin{proof}\leanok
    \uses{thm:mpug_permuted_block_l_symbol_spec,
        thm:mpug_permuted_block_l_compatibility, thm:mpug_czx_anomaly}
    An $L$-symbol is the unique scalar relating the two left boundaries against
    long words, because the target sector is normal. For two generators the
    boundary of two successive actions is $i$ times the integer row
    $(1,-1,1,-1)$ on $\ket{00}$ and $(1,1,-1,-1)$ on $\ket{11}$, and the
    boundary of the fusion followed by the trivial action is the cap
    $F^<_{g,g}=-i\,(0,1,-1,0)$. On every letter of the stacked action, which
    is the image of an integer matrix, the first row is $\mp1$ times the
    second; since a nonempty word begins with a letter, the dressed identity
    holds on every nonempty word with scalar $(-1)^x$. At two identity
    elements the two boundaries coincide, so $L^x_{e,e}=1$. The compatibility
    relation (\ref{eq:mpug_l_compatibility}) at $(e,e,g)$ and $(g,e,e)$, with
    $\omega=1$ there by (\ref{eq:mpug_czx_anomaly}), gives
    $L^x_{e,g}=L^{gx}_{e,e}$ and $L^x_{g,e}=L^x_{e,e}$.
\end{proof}

\begin{lemma}[The anomaly from the \texorpdfstring{$L$}{L}-symbols]
    \label{thm:mpug_czx_l_symbol_ratio}
    \lean{MPOTensor.CZX.omega_eq_lSymbol_ratio,
        MPOTensor.CZX.lSymbol_ratio_mul_omega_eq_neg_one}
    \leanok
    \discussion{7354}
    \uses{thm:mpug_czx_ghz_blocks, thm:mpug_permuted_block_l_compatibility,
        thm:mpug_involution_anomaly_ratio, thm:mpug_czx_anomaly_class}
    For every choice of fusion tensors and every choice of action tensors on
    the two blocked GHZ sectors,
    \begin{align}
        \omega(g,g,g)
         & =\frac{L^x_{g,g}L^x_{g,e}}{L^{gx}_{g,g}L^x_{e,g}},
         &
        \frac{L^0_{g,g}L^0_{g,e}}{L^1_{g,g}L^0_{e,g}}\,\omega(g,e,g)
         & =-1,
    \end{align}
    as in \cite[line~5121]{FrancoRubio2025MPUGauging}. When the fusion and
    action tensors are trivial at the identity this is
    $L^0_{g,g}/L^1_{g,g}=-1$, so no gauge makes every $L$-symbol one.
\end{lemma}

\begin{proof}\leanok
    \uses{thm:mpug_involution_anomaly_ratio, thm:mpug_czx_anomaly_class}
    The compatibility of
    Lemma~\ref{thm:mpug_permuted_block_l_compatibility} holds for every choice
    of action tensors, so Lemma~\ref{thm:mpug_involution_anomaly_ratio} gives
    the first identity. Multiplying by $\omega(g,e,g)$ and using
    Lemma~\ref{thm:mpug_czx_anomaly_class} gives the second.
\end{proof}

\subsection{The displayed CZX circuit tuple}

\begin{definition}[The displayed CZX circuit tuple]
    \label{def:mpug_czx_circuit_tuple}
    \lean{MPOTensor.CZX.bitFlip,
        MPOTensor.CZX.pauliX,
        MPOTensor.CZX.pauliZ,
        MPOTensor.CZX.controlledZ,
        MPOTensor.CZX.w,
        MPOTensor.CZX.lambda,
        MPOTensor.CZX.tildeLambda,
        MPOTensor.CZX.siteBits,
        MPOTensor.CZX.localBits,
        MPOTensor.CZX.localBits_apply_zero,
        MPOTensor.CZX.localBits_apply_one,
        MPOTensor.CZX.localBits_apply_two,
        MPOTensor.CZX.localBits_apply_three,
        MPOTensor.CZX.matterMatrix,
        MPOTensor.CZX.gen,
        MPOTensor.CZX.circuitTuple,
        MPOTensor.CZX.circuitTuple_one,
        MPOTensor.CZX.circuitTuple_gen_one,
        MPOTensor.CZX.circuitTuple_gen_gen}
    \leanok
    \uses{def:mpug_monomial_matrix}
    Number the four matter qubits of two neighboring blocked sites by
    $0,1,2,3$, so that $(x_0,x_1)$ are the qubits of the first site and
    $(x_2,x_3)$ those of the second, and write
    $\ket{x}=\ket{x_0}\otimes\ket{x_1}\otimes\ket{x_2}\otimes\ket{x_3}$ for
    $x\in\mathbb F_2^4$. The elementary gates act by
    \begin{align}
        X_r\ket{x}              & =\ket{x+e_r},          \\
        Z_r\ket{x}              & =(-1)^{x_r}\ket{x},    \\
        \mathrm{CZ}_{rs}\ket{x} & =(-1)^{x_rx_s}\ket{x},
    \end{align}
    and operator products act from right to left. In this qubit
    numbering, the fusion operators $\lambda^R_{g,e}=w_R$ and
    $\lambda^R_{g,g}$ and the modified fusion operator
    $\tilde\lambda^R_{g,g}=\lambda^R_{g,g}Z_1$ of
    \cite[lines~4503--5183]{FrancoRubio2025MPUGauging} are
    \begin{align}
        w             & =(X_0X_1)\,\mathrm{CZ}_{01}\mathrm{CZ}_{12}\,(Z_2Z_3),
        \label{eq:mpug_czx_w}                                                       \\
        \lambda       & =-i\,(X_0X_1X_2Z_3)\,\mathrm{CZ}_{01}\mathrm{CZ}_{12}\,X_2,
        \label{eq:mpug_czx_lambda}                                                  \\
        \tilde\lambda & =\lambda Z_0.
        \label{eq:mpug_czx_tilde_lambda}
    \end{align}
    The displayed circuit tuple, indexed by the ordered pair of
    gauge labels in $\Z_2$ with $e\leftrightarrow0$ and $g\leftrightarrow1$,
    is
    \begin{align}
        R^{\mathrm{CZX}}
         & =\bigl(U_{0,0},U_{0,1},U_{1,0},U_{1,1}\bigr)
        =\bigl(\Id,\Id,w,\tilde\lambda\bigr).
        \label{eq:mpug_czx_tuple}
    \end{align}
    Each site index in $\{0,1,2,3\}$ is read as the pair of matter bits
    $(m^{(1)},m^{(2)})$ with $m^{(1)}$ the leading binary digit, and the
    four-qubit operators are reindexed to the two-site matter space by
    this identification.
    % Source anchor: arXiv:2502.20257, lines 4503--5183 (CZX model and modified
    % fusion operator).
\end{definition}

\begin{lemma}[Phase tables of the CZX fusion operators]
    \label{thm:mpug_czx_phase_tables}
    \lean{MPOTensor.CZX.barFlip,
        MPOTensor.CZX.barFlip_apply,
        MPOTensor.CZX.barFlip_mul_barFlip,
        MPOTensor.CZX.barFlip_symm,
        MPOTensor.CZX.eExponent,
        MPOTensor.CZX.hExponent,
        MPOTensor.CZX.fExponent,
        ZMod.neg_one_pow_val_add,
        MPOTensor.CZX.w_eq,
        MPOTensor.CZX.lambda_eq,
        MPOTensor.CZX.tildeLambda_eq,
        MPOTensor.CZX.fExponent_barFlip,
        MPOTensor.CZX.eExponent_barFlip,
        MPOTensor.CZX.matterBarFlip,
        MPOTensor.CZX.matterBarFlip_symm,
        MPOTensor.CZX.localBits_matterBarFlip,
        MPOTensor.CZX.matterMatrix_monomial,
        MPOTensor.CZX.w_mem_unitaryGroup,
        MPOTensor.CZX.tildeLambda_mem_unitaryGroup}
    \leanok
    \uses{def:mpug_czx_circuit_tuple}
    Write $\overline x=(x_0+1,x_1+1,x_2,x_3)$ and
    \begin{align}
        e(x) & =x_0x_1+x_1x_2+x_2+x_3,     \\
        h(x) & =x_0x_1+x_1x_2+x_1+x_3,     \\
        f(x) & =x_0x_1+x_1x_2+x_0+x_1+x_3.
    \end{align}
    Then
    \begin{align}
        w\ket{x}             & =(-1)^{e(x)}\ket{\overline x},
        \label{eq:mpug_czx_w_action}                            \\
        \lambda\ket{x}       & =-i(-1)^{h(x)}\ket{\overline x},
        \label{eq:mpug_czx_lambda_action}                       \\
        \tilde\lambda\ket{x} & =-i(-1)^{f(x)}\ket{\overline x},
        \label{eq:mpug_czx_tilde_lambda_action}
    \end{align}
    the exponents satisfy $f(\overline x)=e(x)+1$ and
    $e(\overline x)=f(x)+1$. After reindexing to the two-site matter coordinates,
    $w$ and $\tilde\lambda$ are unitary.
\end{lemma}

\begin{proof}\leanok
    \uses{thm:mpug_monomial_matrix_laws, def:mpug_czx_circuit_tuple}
    Each gate is monomial, so (\ref{eq:mpug_monomial_mul}) expresses each
    product as a monomial operator. Reading (\ref{eq:mpug_czx_w}) from right
    to left, the diagonal gates contribute the phase
    $(-1)^{x_0x_1+x_1x_2+x_2+x_3}$ and the two bit flips send $x$ to
    $\overline x$. In (\ref{eq:mpug_czx_lambda}) the first $X_2$ replaces
    $x_2$ by $x_2+1$ in the phase of $\mathrm{CZ}_{12}$, giving
    $x_1(x_2+1)=x_1x_2+x_1$, and the second $X_2$ restores $x_2$, so the
    permutation is again $\overline{\,\cdot\,}$. The factor $Z_0$ in
    (\ref{eq:mpug_czx_tilde_lambda}) adds $x_0$ to the exponent. The
    identities between $e$ and $f$ are checked on the sixteen bit strings.
    After relabelling the computational basis as two four-level matter sites,
    unitarity follows from Lemma~\ref{thm:mpug_monomial_matrix_laws} because
    all phases are unimodular.
\end{proof}

\begin{lemma}[Right movement coordinates of the displayed CZX factors]
    \label{thm:mpug_czx_displayed_right_movement_coordinates}
    \lean{MPOTensor.CZX.displayedSourceFactors_sourceWR_coordinates}
    \leanok
    \discussion{7355}
    \uses{def:mpug_czx_displayed_source_factors, def:mpug_staircase_gates,
        def:mpug_czx_circuit_tuple}
    Let $w_R$ be the crossed contraction of the displayed source factors
    of the exact CZX tensor $B$, and let $e_r$ be their right-rank
    identification with $\{0,1,2,3\}$. Let $W$ denote the matrix of
    $w=(X_0X_1)\mathrm{CZ}_{01}\mathrm{CZ}_{12}(Z_2Z_3)$ on two
    four-level matter sites, with each site read in leading-bit-first
    binary order. Identifying the right movement contraction in
    \cite[lines~811--869, 4559--4659, 4860--4888]{FrancoRubio2025MPUGauging}
    without additional hypotheses, for every $i,r,a,j\in\{0,1,2,3\}$,
    \begin{align}
        (w_R)_{(i,e_r^{-1}(r)),(e_r^{-1}(a),j)}
         & =W_{(i,r),(a,j)}.
    \end{align}
\end{lemma}

\begin{proof}\leanok
    \uses{def:mpug_czx_displayed_source_factors, def:mpug_staircase_gates,
        def:mpug_czx_displayed_factors, def:mpug_czx_circuit_tuple,
        thm:mpug_czx_phase_tables}
    Pulling back the rank indices expands the crossed contraction as
    $\sum_{\beta=0}^1(X_1)_{(i,\beta),a}(Y_1)_{r,(\beta,j)}$.
    The sparse action (\ref{eq:mpug_czx_w_action}) gives
    \begin{align}
        W_{(i,r),(a,j)}
         & =\begin{cases}
                (-1)^{e(x)}, & i=3-a\text{ and }r=j, \\
                0,           & \text{otherwise},
            \end{cases}
    \end{align}
    where $x=(x_0,x_1,x_2,x_3)\in\{0,1\}^4$ satisfies
    $a=2x_0+x_1$ and $j=2x_2+x_3$.
    Substituting the displayed factors and first discarding the zero
    entries of the empty factor $X_1$ by comparing $i$ and $a$ leaves
    only the comparison of $r$ and $j$, and on the surviving support the
    phases agree.
\end{proof}

\begin{lemma}[Left movement coordinates of the displayed CZX factors]
    \label{thm:mpug_czx_displayed_left_movement_coordinates}
    \lean{MPOTensor.CZX.displayedSourceFactors_sourceWL_coordinates}
    \leanok
    \discussion{7355}
    \uses{def:mpug_czx_displayed_source_factors, def:mpug_staircase_gates,
        def:mpug_czx_circuit_tuple}
    Let $w_L$ be the crossed contraction of the displayed source factors
    of the exact CZX tensor $B$, and let $e_l$ be their left-rank
    identification with $\{0,1,2,3\}$. Let $W$ denote the matrix of
    $w=(X_0X_1)\mathrm{CZ}_{01}\mathrm{CZ}_{12}(Z_2Z_3)$ on two
    four-level matter sites, with each site read in leading-bit-first
    binary order. In the same matter coordinates and without additional
    hypotheses, the left movement identification in
    \cite[lines~811--869, 4559--4659, 4860--4888]{FrancoRubio2025MPUGauging}
    holds for every $l,i,j,k\in\{0,1,2,3\}$:
    \begin{align}
        (w_L)_{(e_l^{-1}(l),i),(j,e_l^{-1}(k))}
         & =(W^\dagger)_{(l,i),(j,k)}.
    \end{align}
\end{lemma}

\begin{proof}\leanok
    \uses{def:mpug_czx_displayed_source_factors, def:mpug_staircase_gates,
        def:mpug_czx_displayed_factors, def:mpug_czx_circuit_tuple,
        thm:mpug_czx_phase_tables}
    Pulling back the rank indices expands the crossed contraction as
    $\sum_{\beta=0}^1(Y_2)_{l,(j,\beta)}(X_2)_{(\beta,i),k}$, while the
    conjugated reversed entry is
    $(W^\dagger)_{(l,i),(j,k)}=\overline{W_{(j,k),(l,i)}}$.
    By (\ref{eq:mpug_czx_w_action}), it vanishes unless $j=3-l$ and
    $k=i$; on this support it is $\overline{(-1)^{e(x)}}$, where
    $x=(x_0,x_1,x_2,x_3)\in\{0,1\}^4$ satisfies
    $l=2x_0+x_1$ and $i=2x_2+x_3$.
    Substituting the displayed factors and first discarding the zero
    entries of the filled factor $Y_2$ by comparing $l$ and $j$ leaves
    the comparison of $i$ and $k$, on which the remaining phases agree.
\end{proof}

% Scope of #7355 here: only the two displayed movement-coordinate identities.
% No full truncated-symmetry formula, associator, sign, or parent issue
% completion is asserted.

\begin{lemma}[Local monomial action of the CZX Gauss operator]
    \label{thm:mpug_czx_local_monomial_action}
    \lean{MPOTensor.CZX.iPow,
        MPOTensor.CZX.iPow_add,
        MPOTensor.CZX.iPow_natCast,
        MPOTensor.CZX.iPow_zero,
        MPOTensor.CZX.iPow_two,
        MPOTensor.CZX.iPow_three,
        MPOTensor.CZX.iPow_two_mul_val,
        MPOTensor.CZX.iPow_two_mul_cast,
        MPOTensor.CZX.uExponent,
        MPOTensor.CZX.phaseExponent,
        MPOTensor.CZX.phaseExponent_zero_zero,
        MPOTensor.CZX.phaseExponent_zero_one,
        MPOTensor.CZX.phaseExponent_one_zero,
        MPOTensor.CZX.phaseExponent_one_one,
        MPOTensor.CZX.localPerm,
        MPOTensor.CZX.localPhase,
        MPOTensor.CZX.gen_inv,
        MPOTensor.CZX.gen_mul_gen,
        MPOTensor.CZX.toAdd_gen_ne_zero,
        MPOTensor.CZX.gaussOperator_circuitTuple_gen}
    \leanok
    \uses{def:mpug_czx_circuit_tuple, def:mpug_gauss_operator}
    Let $a,a'\in\mathbb F_2$ be the two gauge bits and write
    $s=(x_0,x_1,a,x_2,x_3,a')$ and
    $\sigma(s)=(x_0+1,x_1+1,a+1,x_2,x_3,a'+1)$. The nontrivial local Gauss
    operator obtained from (\ref{eq:mpug_czx_tuple}) and
    (\ref{eq:mpug_gauss_operator}) is the monomial operator
    \begin{align}
        G\ket{s} & =i^{q(s)}\ket{\sigma(s)}
        \label{eq:mpug_czx_monomial_action}
    \end{align}
    with, for $u(x)=2(x_0x_1+x_1x_2+x_3)\bmod4$, the phase table
    \begin{align}
        q(x_0,x_1,0,x_2,x_3,0) & =3+u(x)+2x_2,      \\
        q(x_0,x_1,0,x_2,x_3,1) & =2+u(x)+2x_0+2x_1, \\
        q(x_0,x_1,1,x_2,x_3,0) & =u(x)+2x_2,        \\
        q(x_0,x_1,1,x_2,x_3,1) & =3+u(x)+2x_0+2x_1
        \label{eq:mpug_czx_phase_table}
    \end{align}
    modulo $4$.
\end{lemma}

\begin{proof}\leanok
    \uses{thm:mpug_gauss_operator_monomial, thm:mpug_czx_phase_tables}
    By Lemma~\ref{thm:mpug_gauss_operator_monomial} applied to the monomial
    unitaries of Lemma~\ref{thm:mpug_czx_phase_tables}, the operator
    $\mathcal G_g$ is monomial. The permutation part is
    $\overline{\,\cdot\,}$ on the matter bits in all four gauge sectors,
    since $\overline{\,\cdot\,}$ is an involution, and $T_g$ flips both
    gauge bits. The phases are read off
    (\ref{eq:mpug_czx_w_action})--(\ref{eq:mpug_czx_tilde_lambda_action})
    and the adjoint formula (\ref{eq:mpug_monomial_adjoint}) together with
    $f(\overline x)=e(x)+1$ and $e(\overline x)=f(x)+1$: for example the
    first row is $-i(-1)^{e(x)}=i^{3+2e(x)}=i^{3+u(x)+2x_2}$, and the other
    three rows follow in the same way. The four identities between the phase
    table and $e$, $f$ are checked on the sixteen bit strings.
\end{proof}

\subsection{Orbits, holonomy, and the exact dimension}

\begin{definition}[Chain configurations, bond flips, and labels]
    \label{def:mpug_czx_chain_labels}
    \lean{MPOTensor.CZX.Site,
        MPOTensor.CZX.siteDecode,
        MPOTensor.CZX.siteDecode_equivFin,
        MPOTensor.CZX.chainDecode,
        MPOTensor.CZX.chainDecode_apply,
        MPOTensor.CZX.flipPattern,
        MPOTensor.CZX.chainFlip,
        MPOTensor.CZX.chainFlip_apply,
        MPOTensor.CZX.succ_ne_self,
        MPOTensor.CZX.succ_succ_ne_self,
        MPOTensor.CZX.flipPattern_apply_self,
        MPOTensor.CZX.flipPattern_apply_succ,
        MPOTensor.CZX.flipPattern_apply_of_ne,
        MPOTensor.CZX.bondExponent,
        MPOTensor.CZX.bLabel,
        MPOTensor.CZX.pLabel}
    \leanok
    \uses{thm:mpug_czx_local_monomial_action}
    A computational-basis configuration of the periodic chain of $N$ blocked
    sites is
    \begin{align}
        s & =\bigl(m^{(1)}_k,m^{(2)}_k,g_k\bigr)_{k\in\Z/N\Z}\in S_N:=(\mathbb F_2^3)^N,
        \label{eq:mpug_czx_configuration}
    \end{align}
    read from the enumerated site indices of the placed Gauss projector. The
    bond-$j$ flip $\sigma_j$ adds $1$ to $m^{(1)}_j$, $m^{(2)}_j$, $g_j$, and
    $g_{j+1}$. The bond-$j$ phase exponent $q_j(s)$ is the table
    (\ref{eq:mpug_czx_phase_table}) evaluated at
    $x=(m^{(1)}_j,m^{(2)}_j,m^{(1)}_{j+1},m^{(2)}_{j+1})$ and
    $(a,a')=(g_j,g_{j+1})$. The affine labels are
    \begin{align}
        p_k(s) & =m^{(1)}_k+m^{(2)}_k,
        \label{eq:mpug_czx_p_label}            \\
        b_k(s) & =g_k+m^{(1)}_k+m^{(1)}_{k-1}.
        \label{eq:mpug_czx_b_label}
    \end{align}
\end{definition}

\begin{lemma}[Permutation-orbit classification]
    \label{thm:mpug_czx_orbit_classification}
    \lean{MPOTensor.CZX.fiberEquiv,
        MPOTensor.CZX.fiberEquiv_apply,
        MPOTensor.CZX.pLabel_fiberEquiv,
        MPOTensor.CZX.bLabel_fiberEquiv,
        MPOTensor.CZX.chainFlip_fiberEquiv}
    \leanok
    \uses{def:mpug_czx_chain_labels}
    For $N\geq2$, the map
    \begin{align}
        \bigl((p,b),\gamma\bigr)
         & \longmapsto
        \bigl(m^{(1)}_k,m^{(2)}_k,g_k\bigr)_k
        =\bigl(\gamma_k,\ p_k+\gamma_k,\ b_k+\gamma_k+\gamma_{k-1}\bigr)_k
        \label{eq:mpug_czx_fiber_coordinates}
    \end{align}
    is a bijection $\mathbb F_2^N\times\mathbb F_2^N\times\mathbb F_2^N\to S_N$
    whose inverse reads off $p(s)$, $b(s)$, and $\gamma=m^{(1)}$. In these
    coordinates $\sigma_j$ is $\gamma\mapsto\gamma+e_j$ and leaves $(p,b)$
    unchanged.
\end{lemma}

\begin{proof}\leanok
    \uses{def:mpug_czx_chain_labels}
    Both compositions of (\ref{eq:mpug_czx_fiber_coordinates}) with the
    label map are the identity, by the identity $a+a=0$ in $\mathbb F_2$.
    The flip $\sigma_j$ adds $((1,1),1)$ at site $j$ and $((0,0),1)$ at
    site $j+1$; in fiber coordinates this is $\gamma_j\mapsto\gamma_j+1$
    because $k-1=j$ exactly when $k=j+1$.
\end{proof}

\begin{lemma}[Involution and neighboring holonomy]
    \label{thm:mpug_czx_holonomy}
    \lean{MPOTensor.CZX.bondExponent_add_bondExponent_chainFlip,
        MPOTensor.CZX.bondExponent_holonomy,
        MPOTensor.CZX.bondExponent_chainFlip_of_disjoint,
        MPOTensor.CZX.fiberPhase,
        MPOTensor.CZX.fiberPhase_mul_fiberPhase_flip,
        MPOTensor.CZX.fiberPhase_flip_of_disjoint,
        MPOTensor.CZX.fiberPhase_holonomy,
        MPOTensor.CZX.fiberPhase_ne_zero}
    \leanok
    \uses{def:mpug_czx_chain_labels}
    For every $N\geq3$, every bond $j$, and every configuration $s$,
    \begin{align}
        q_j(s)+q_j(\sigma_js)
         & =0,
        \label{eq:mpug_czx_involution}                \\
        q_j(s)+q_{j+1}(\sigma_js)
         & =q_{j+1}(s)+q_j(\sigma_{j+1}s)+2b_{j+1}(s)
        \label{eq:mpug_czx_holonomy}
    \end{align}
    modulo $4$, and $q_j(\sigma_ks)=q_j(s)$ whenever
    $k\notin\{j-1,j,j+1\}$. For the phases
    $\varphi_j(s)=i^{q_j(s)}$, (\ref{eq:mpug_czx_involution}) gives
    $\varphi_j(s)\varphi_j(\sigma_js)=1$, while
    (\ref{eq:mpug_czx_holonomy}) gives the neighboring square relation with
    factor $(-1)^{b_{j+1}(s)}$.
\end{lemma}

\begin{proof}\leanok
    \uses{def:mpug_czx_chain_labels}
    The exponent $q_j$ depends only on the six bits at the sites $j$ and
    $j+1$, and $\sigma_k$ changes only the sites $k$ and $k+1$; this gives
    the disjoint-bond statement. For (\ref{eq:mpug_czx_involution}) the
    flip $\sigma_j$ changes $x_0$, $x_1$, $a$, and $a'$, and the identity is
    verified on the $64$ values of the six bits. For
    (\ref{eq:mpug_czx_holonomy}) the sites $j$, $j+1$, $j+2$ are distinct
    because $N\geq3$; $\sigma_j$ changes the left gauge bit seen by the bond
    $(j+1,j+2)$, and $\sigma_{j+1}$ changes $x_2$, $x_3$, and $a'$ seen by
    the bond $(j,j+1)$, so the identity is verified on the $512$ values of
    the nine bits at the three sites.
\end{proof}

\begin{lemma}[Fixed spaces on the orbit fibers]
    \label{thm:mpug_czx_good_fibers}
    \lean{MPOTensor.CZX.isTrivialHolonomy_fiberPhase_iff}
    \leanok
    \uses{def:mpug_trivial_holonomy, thm:mpug_czx_orbit_classification}
    Let $N\geq3$. In the fiber coordinates
    (\ref{eq:mpug_czx_fiber_coordinates}), the phases $i^{q_j}$ have trivial
    holonomy on the fiber over $(p,b)$ if and only if $b=0$.
\end{lemma}

\begin{proof}\leanok
    \uses{thm:mpug_czx_holonomy, thm:mpug_czx_orbit_classification}
    If $b=0$, then for $k=j+1$ or $j=k+1$ the square identity is
    (\ref{eq:mpug_czx_holonomy}) with $b_{j+1}=0$, for $k=j$ it is trivial,
    and otherwise the bonds are disjoint and both sides agree. Conversely,
    trivial holonomy at the square spanned by the bonds $k-1$ and $k$ at
    $\gamma=0$ gives $(-1)^{b_k}=1$ by (\ref{eq:mpug_czx_holonomy}), since the
    phases are nonzero, hence $b_k=0$ for every $k$.
\end{proof}

\begin{lemma}[The placed CZX Gauss operators]
    \label{thm:mpug_czx_placed_operators}
    \lean{MPOTensor.CZX.windowCoordinates,
        MPOTensor.CZX.placedGaussOperator,
        MPOTensor.CZX.sum_multiplicative_zmod_two,
        MPOTensor.CZX.gaussProjector_circuitTuple,
        MPOTensor.CZX.placedGaussProjector_circuitTuple,
        MPOTensor.CZX.inv_two_smul_add_mulVec_eq_iff,
        MPOTensor.CZX.commonFixedSubmodule_placedGaussProjector_circuitTuple,
        MPOTensor.CZX.placedGaussOperator_eq_monomial,
        MPOTensor.CZX.siteBits_eq_localBits_zero,
        MPOTensor.CZX.siteBits_eq_localBits_one,
        MPOTensor.CZX.barFlip_apply_zero,
        MPOTensor.CZX.barFlip_apply_one,
        MPOTensor.CZX.barFlip_apply_two,
        MPOTensor.CZX.barFlip_apply_three,
        MPOTensor.CZX.toAdd_mul_gen_inv,
        MPOTensor.CZX.toAdd_gen_mul,
        MPOTensor.CZX.gaussLegAction_fst,
        MPOTensor.CZX.gaussLegAction_snd,
        MPOTensor.CZX.chainDecode_window,
        MPOTensor.CZX.chainDecode_windowPerm,
        MPOTensor.CZX.localPhase_window,
        MPOTensor.CZX.fiberCoordinates,
        MPOTensor.CZX.placedGaussOperator_eq_reindex}
    \leanok
    \uses{def:mpug_placed_gauss_projector, def:mpug_czx_circuit_tuple,
        def:mpug_czx_chain_labels, def:mpug_common_fixed_subspace}
    Let $N\geq2$ and let $G_j$ be the local Gauss operator
    (\ref{eq:mpug_czx_monomial_action}) placed on the bond $(j,j+1)$. For
    $\Z_2$ the placed Gauss projector of
    Definition~\ref{def:mpug_placed_gauss_projector} is
    \begin{align}
        \widetilde P_j & =\tfrac12\bigl(\Id+G_j\bigr),
        \label{eq:mpug_czx_placed_projector}
    \end{align}
    so the common fixed subspace $\mathcal V_N(R^{\mathrm{CZX}})$ of the
    $\widetilde P_j$ is the common fixed subspace of the $G_j$. In the
    coordinates (\ref{eq:mpug_czx_configuration}), $G_j$ is the monomial
    operator $G_j\ket{s}=i^{q_j(s)}\ket{\sigma_js}$, and in the fiber
    coordinates (\ref{eq:mpug_czx_fiber_coordinates}) it is the monomial
    operator with permutation $\gamma\mapsto\gamma+e_j$ and phase
    $i^{q_j}$.
\end{lemma}

\begin{proof}\leanok
    \uses{def:mpug_gauss_projector, thm:mpug_gauss_representation,
        thm:mpug_czx_local_monomial_action, thm:mpug_embed_monomial,
        thm:mpug_czx_orbit_classification, thm:mpug_mem_common_fixed_subspace}
    The average over $\Z_2=\{e,g\}$ is $\frac12(\mathcal G_e+\mathcal G_g)$
    with $\mathcal G_e=\Id$, and relabelling and placement are linear and
    unital, which gives (\ref{eq:mpug_czx_placed_projector}). A vector $v$
    satisfies $\frac12(v+G_jv)=v$ exactly when $G_jv=v$.
    Lemma~\ref{thm:mpug_czx_local_monomial_action} and
    Lemma~\ref{thm:mpug_embed_monomial} express $G_j$ as a monomial
    operator; decoding the two-site window at $j$ into the bits of the sites
    $j$ and $j+1$ identifies its permutation with $\sigma_j$ and its phase
    with $i^{q_j(s)}$. Finally
    Lemma~\ref{thm:mpug_czx_orbit_classification} reindexes both to the
    fiber coordinates.
\end{proof}

\begin{lemma}[Exact dimension for the displayed CZX circuit tuple]
    \label{thm:mpug_czx_dimension}
    \lean{MPOTensor.CZX.finrank_commonFixedSubmodule_placedGaussProjector_circuitTuple}
    \leanok
    \uses{def:mpug_placed_gauss_projector, def:mpug_czx_circuit_tuple,
        def:mpug_common_fixed_subspace}
    Answering the growth question of
    \cite[lines~5198--5204]{FrancoRubio2025MPUGauging} for the displayed
    circuit tuple, on a periodic chain of $N\geq3$ blocked sites,
    \begin{align}
        \dim\mathcal V_N(R^{\mathrm{CZX}}) & =2^N.
        \label{eq:mpug_czx_dimension}
    \end{align}
\end{lemma}

\begin{proof}\leanok
    \uses{thm:mpug_czx_placed_operators,
        thm:mpug_monomial_fixed_space_criterion,
        thm:mpug_czx_holonomy, thm:mpug_czx_good_fibers}
    By Lemma~\ref{thm:mpug_czx_placed_operators},
    $\mathcal V_N(R^{\mathrm{CZX}})$ is the common fixed subspace of monomial
    operators which, in fiber coordinates, have permutations
    $\gamma\mapsto\gamma+e_j$ on the fibers over
    $(p,b)\in\mathbb F_2^N\times\mathbb F_2^N$, and by
    (\ref{eq:mpug_czx_involution}) these operators are involutions. The
    monomial fixed-space criterion
    (\ref{eq:mpug_monomial_fixed_space_count}) counts the fibers with trivial
    holonomy, which by Lemma~\ref{thm:mpug_czx_good_fibers} are exactly
    the fibers with $b=0$. There is one such $b$ and there are $2^N$ values
    of $p$, so
    \begin{align}
        \dim\mathcal V_N(R^{\mathrm{CZX}})
         & =\sum_{p\in\mathbb F_2^N}1=2^N.
    \end{align}
\end{proof}

\begin{lemma}[The CZX dimension in the gauge-invariant subspace of the displayed tuple]
    \label{thm:mpug_czx_dimension_of_completion}
    \lean{MPOTensor.CZX.finrank_gaugeInvariantSubspace_circuitTuple}
    \leanok
    \uses{def:mpug_gauge_invariant_subspace, def:mpug_czx_circuit_tuple}
    For the gauge-invariant subspace~(\ref{eq:mpug_gauge_invariant_subspace}) at
    $R=R^{\mathrm{CZX}}$ and $N\geq3$,
    $\dim_{\C}\mathcal V_N(R^{\mathrm{CZX}})\allowbreak=2^N$. The value is thereby
    recorded for the subspace that
    Definition~\ref{def:mpug_gauge_invariant_subspace} attaches to an arbitrary
    tuple, so exhibiting $R^{\mathrm{CZX}}$ as a completion of prescribed
    defect maps with matter defect states satisfying local defect support,
    exact defect covariance, and nonvanishing of the gauged vector would put
    the exact value beside the lower bound $\dim_{\C}\mathcal V_N(R)\geq1$ of
    Lemma~\ref{thm:mpug_gauge_invariant_subspace_finrank}. None of those
    hypotheses is established here. The statement is about the displayed
    circuit tuple and asserts no membership in the full physical completion
    class of the CZX defect data.
\end{lemma}

\begin{proof}\leanok
    \uses{thm:mpug_czx_dimension, def:mpug_placed_gauss_projector}
    Both subspaces are the common fixed subspace of the placed Gauss
    projectors $\widetilde P_j(R^{\mathrm{CZX}})$ of
    Definition~\ref{def:mpug_placed_gauss_projector} acting on the
    configuration space of the chain, so
    Lemma~\ref{thm:mpug_czx_dimension} gives the value.
\end{proof}

\subsection{A second displayed circuit tuple at three and four sites}

Constraining the block $U_{1,1}$ of (\ref{eq:mpug_czx_tuple}) on the defect
subspace $\spn_{\C}\{\ket{0011},\ket{1100}\}$ leaves gauge freedom in
$\tilde\lambda$. Replacing $\tilde\lambda$ by another unitary with identical
action on this subspace yields an alternate circuit tuple whose common $+1$
eigenspace dimension differs at three and four blocked sites.

\begin{definition}[The second displayed circuit tuple]
    \label{def:mpug_czx_star_tuple}
    \lean{MPOTensor.CZX.rStarExponent,
        MPOTensor.CZX.rStar,
        MPOTensor.CZX.rStar_eq,
        MPOTensor.CZX.rStarExponent_zero_zero_one_one,
        MPOTensor.CZX.rStarExponent_one_one_zero_zero,
        MPOTensor.CZX.rStar_mulVec_eq_self,
        MPOTensor.CZX.tildeLambdaStar,
        MPOTensor.CZX.tildeLambdaStar_eq,
        MPOTensor.CZX.tildeLambdaStar_mem_unitaryGroup,
        MPOTensor.CZX.circuitTupleStar,
        MPOTensor.CZX.circuitTupleStar_one,
        MPOTensor.CZX.circuitTupleStar_gen_one,
        MPOTensor.CZX.circuitTupleStar_gen_gen}
    \leanok
    \uses{def:mpug_czx_circuit_tuple, thm:mpug_czx_phase_tables}
    In the qubit numbering of Definition~\ref{def:mpug_czx_circuit_tuple}, set
    \begin{align}
        r_\star(x) & =x_1+x_3+x_0x_1+x_0x_3+x_1x_2+x_2x_3,
        \label{eq:mpug_czx_star_phase}                                         \\
        R_\star    & =Z_1Z_3\,\mathrm{CZ}_{01}\mathrm{CZ}_{03}\mathrm{CZ}_{12}
        \mathrm{CZ}_{23}.
        \label{eq:mpug_czx_star_gate}
    \end{align}
    Then $R_\star\ket{x}=(-1)^{r_\star(x)}\ket{x}$, and
    $r_\star(0011)=r_\star(1100)=0$, so $R_\star$ fixes the prescribed defect
    subspace $\spn_{\C}\{\ket{0011},\ket{1100}\}$ pointwise. The second
    displayed circuit tuple is
    \begin{align}
        R^\star
         & =\bigl(U^\star_{0,0},U^\star_{0,1},U^\star_{1,0},U^\star_{1,1}\bigr)
        =\bigl(\Id,\Id,w,\tilde\lambda R_\star\bigr),
        \label{eq:mpug_czx_star_tuple}
    \end{align}
    with
    $\tilde\lambda R_\star\ket{x}=-i(-1)^{f(x)+r_\star(x)}\ket{\overline x}$;
    after reindexing to the two-site matter coordinates it is unitary. It has the
    same action as $R^{\mathrm{CZX}}$ on that defect subspace. The four defect
    domains and prescribed maps of
    Definition~\ref{def:mpug_czx_defect_maps} determine a four-sector
    completion class, and both tuples extend all four prescribed maps: the
    first by Lemma~\ref{thm:mpug_czx_tuple_completion} and the second by
    Lemma~\ref{thm:mpug_czx_star_tuple_completion}, because $R_\star$ fixes
    the defect domain of the pair $(g,g)$ pointwise and the other three
    operators are those of $R^{\mathrm{CZX}}$.
    % Source anchor: arXiv:2502.20257, lines 4503--5183 (CZX model, defect and
    % fusion data, modified Gauss law).
\end{definition}

\begin{lemma}[Local monomial action of the second tuple's Gauss operator]
    \label{thm:mpug_czx_star_local_monomial_action}
    \lean{MPOTensor.CZX.rStarExponent_barFlip,
        MPOTensor.CZX.starDefect,
        MPOTensor.CZX.phaseExponentStar,
        MPOTensor.CZX.phaseExponentStar_zero_zero,
        MPOTensor.CZX.phaseExponentStar_zero_one,
        MPOTensor.CZX.phaseExponentStar_one_zero,
        MPOTensor.CZX.phaseExponentStar_one_one,
        MPOTensor.CZX.localPhaseStar,
        MPOTensor.CZX.gaussOperator_circuitTupleStar_gen}
    \leanok
    \uses{def:mpug_czx_star_tuple, def:mpug_gauss_operator}
    The nontrivial local Gauss operator obtained from
    (\ref{eq:mpug_czx_star_tuple}) and (\ref{eq:mpug_gauss_operator}) is the
    monomial operator $G^\star\ket{s}=i^{q^\star(s)}\ket{\sigma(s)}$ with the
    same permutation part $\sigma$ as in
    (\ref{eq:mpug_czx_monomial_action}) and with
    \begin{align}
        q^\star(x_0,x_1,a,x_2,x_3,a') & =q(x_0,x_1,a,x_2,x_3,a')
        +2\delta_\star(x,a,a')
        \label{eq:mpug_czx_star_phase_table}
    \end{align}
    modulo $4$, where $\delta_\star(x,0,0)=r_\star(\overline x)$,
    $\delta_\star(x,1,1)=r_\star(x)$, and $\delta_\star$ vanishes in the two
    remaining gauge sectors. The phase polynomial satisfies
    \begin{align}
        r_\star(\overline x) & =r_\star(x)+x_0+x_1+x_2+x_3.
        \label{eq:mpug_czx_star_bar}
    \end{align}
\end{lemma}

\begin{proof}\leanok
    \uses{thm:mpug_gauss_operator_monomial, thm:mpug_czx_phase_tables,
        def:mpug_czx_star_tuple}
    By Lemma~\ref{thm:mpug_gauss_operator_monomial} the operator is monomial,
    and the two blocks $U^\star_{0,1}=\Id$ and $U^\star_{1,0}=w$ agree with
    those of (\ref{eq:mpug_czx_tuple}), so the $(0,1)$ and $(1,0)$ rows of the
    phase table are unchanged. In the $(1,1)$ sector the phase acquires the
    factor $(-1)^{r_\star(x)}$ from
    $\tilde\lambda R_\star\ket{x}=-i(-1)^{f(x)+r_\star(x)}\ket{\overline x}$; in
    the $(0,0)$ sector the adjoint formula
    (\ref{eq:mpug_monomial_adjoint}) evaluates that phase at $\overline x$ and
    conjugates it, which contributes $(-1)^{r_\star(\overline x)}$. Since
    $R_\star$ is diagonal, the permutation part is that of
    (\ref{eq:mpug_czx_monomial_action}). Identity
    (\ref{eq:mpug_czx_star_bar}) and the four rows of
    (\ref{eq:mpug_czx_star_phase_table}) are checked on the sixteen bit
    strings.
\end{proof}

\begin{lemma}[Involution and neighboring holonomy for the second tuple]
    \label{thm:mpug_czx_star_holonomy}
    \lean{MPOTensor.CZX.bondExponentStar,
        MPOTensor.CZX.starHolonomyLabel,
        MPOTensor.CZX.bondExponentStar_add_bondExponentStar_chainFlip,
        MPOTensor.CZX.bondExponentStar_holonomy,
        MPOTensor.CZX.bondExponentStar_chainFlip_of_disjoint,
        MPOTensor.CZX.starFiberLabel,
        MPOTensor.CZX.starHolonomyLabel_fiberEquiv,
        MPOTensor.CZX.fiberPhaseStar,
        MPOTensor.CZX.fiberPhaseStar_mul_fiberPhaseStar_flip,
        MPOTensor.CZX.fiberPhaseStar_flip_of_disjoint,
        MPOTensor.CZX.fiberPhaseStar_holonomy,
        MPOTensor.CZX.fiberPhaseStar_ne_zero}
    \leanok
    \uses{def:mpug_czx_chain_labels, def:mpug_czx_star_tuple}
    Let $q^\star_j(s)$ be the table (\ref{eq:mpug_czx_star_phase_table})
    evaluated on the bond $(j,j+1)$ as in
    Definition~\ref{def:mpug_czx_chain_labels}, and set
    \begin{align}
        c_j(s) & =b_{j+1}(s)\bigl(1+p_j(s)+p_{j+1}(s)\bigr)
        +b_{j+2}(s)\bigl(p_{j+1}(s)+p_{j+2}(s)\bigr).
        \label{eq:mpug_czx_star_holonomy_label}
    \end{align}
    For every $N\geq3$, every bond $j$, and every configuration $s$,
    \begin{align}
        q^\star_j(s)+q^\star_j(\sigma_js)
         & =0,
        \label{eq:mpug_czx_star_involution}                   \\
        q^\star_j(s)+q^\star_{j+1}(\sigma_js)
         & =q^\star_{j+1}(s)+q^\star_j(\sigma_{j+1}s)+2c_j(s)
        \label{eq:mpug_czx_star_holonomy}
    \end{align}
    modulo $4$, and $q^\star_j(\sigma_ks)=q^\star_j(s)$ whenever
    $k\notin\{j-1,j,j+1\}$. In the fiber coordinates
    (\ref{eq:mpug_czx_fiber_coordinates}) the label $c_j$ depends only on the
    orbit labels, $c_j(p,b)=b_{j+1}(1+p_j+p_{j+1})+b_{j+2}(p_{j+1}+p_{j+2})$,
    and the phases $\varphi^\star_j=i^{q^\star_j}$ are involutive and satisfy
    the neighboring square relation with factor $(-1)^{c_j(p,b)}$.
\end{lemma}

\begin{proof}\leanok
    \uses{thm:mpug_czx_star_local_monomial_action,
        thm:mpug_czx_orbit_classification}
    The exponent $q^\star_j$ depends only on the six bits at the sites $j$ and
    $j+1$, which gives the disjoint-bond statement. For
    (\ref{eq:mpug_czx_star_involution}) the flip $\sigma_j$ changes $x_0$,
    $x_1$, $a$, and $a'$, and the identity is verified on the $64$ values of
    those six bits. For (\ref{eq:mpug_czx_star_holonomy}) the sites $j$,
    $j+1$, $j+2$ are distinct because $N\geq3$, and the identity is verified on
    the $512$ values of the nine bits at those three sites. In fiber
    coordinates $p$ and $b$ are read off by
    Lemma~\ref{thm:mpug_czx_orbit_classification}, which turns $c_j$ into the
    displayed polynomial in the orbit labels alone.
\end{proof}

\begin{lemma}[Fixed spaces on the orbit fibers for the second tuple]
    \label{thm:mpug_czx_star_good_fibers}
    \lean{MPOTensor.CZX.isTrivialHolonomy_fiberPhaseStar_iff}
    \leanok
    \uses{def:mpug_trivial_holonomy, thm:mpug_czx_orbit_classification,
        thm:mpug_czx_star_holonomy}
    Let $N\geq3$. In the fiber coordinates
    (\ref{eq:mpug_czx_fiber_coordinates}), the phases $i^{q^\star_j}$ have
    trivial holonomy on the fiber over $(p,b)$ if and only if $c_j(p,b)=0$ for
    every $j\in\Z/N\Z$.
\end{lemma}

\begin{proof}\leanok
    \uses{thm:mpug_czx_star_holonomy, thm:mpug_czx_orbit_classification}
    If all $c_j(p,b)$ vanish, then for $k=j+1$ or $j=k+1$ the square identity is
    (\ref{eq:mpug_czx_star_holonomy}) with vanishing correction, for $k=j$ it is
    trivial, and otherwise the bonds are disjoint and both sides agree.
    Conversely, trivial holonomy at the square spanned by the bonds $j$ and
    $j+1$ at $\gamma=0$ gives $(-1)^{c_j(p,b)}=1$ by
    (\ref{eq:mpug_czx_star_holonomy}), since the phases are nonzero, hence
    $c_j(p,b)=0$.
\end{proof}

\begin{lemma}[The placed Gauss operators of the second tuple]
    \label{thm:mpug_czx_star_placed_operators}
    \lean{MPOTensor.CZX.placedGaussOperatorStar,
        MPOTensor.CZX.gaussProjector_circuitTupleStar,
        MPOTensor.CZX.placedGaussProjector_circuitTupleStar,
        MPOTensor.CZX.commonFixedSubmodule_placedGaussProjector_circuitTupleStar,
        MPOTensor.CZX.placedGaussOperatorStar_eq_monomial,
        MPOTensor.CZX.localPhaseStar_window,
        MPOTensor.CZX.placedGaussOperatorStar_eq_reindex}
    \leanok
    \uses{def:mpug_placed_gauss_projector, def:mpug_czx_star_tuple,
        def:mpug_czx_chain_labels, def:mpug_common_fixed_subspace}
    Let $N\geq2$ and let $G^\star_j$ be the local Gauss operator of
    Lemma~\ref{thm:mpug_czx_star_local_monomial_action} placed on the bond
    $(j,j+1)$. The placed Gauss projector of
    Definition~\ref{def:mpug_placed_gauss_projector} is
    $\widetilde P^\star_j=\tfrac12(\Id+G^\star_j)$, so the common fixed subspace
    $\mathcal V_N(R^\star)$ of the $\widetilde P^\star_j$ is the common fixed
    subspace of the $G^\star_j$. In the coordinates
    (\ref{eq:mpug_czx_configuration}), $G^\star_j$ is the monomial operator
    $G^\star_j\ket{s}=i^{q^\star_j(s)}\ket{\sigma_js}$, and in the fiber
    coordinates (\ref{eq:mpug_czx_fiber_coordinates}) it is the monomial
    operator with permutation $\gamma\mapsto\gamma+e_j$ and phase
    $i^{q^\star_j}$.
\end{lemma}

\begin{proof}\leanok
    \uses{def:mpug_gauss_projector, thm:mpug_czx_star_local_monomial_action,
        thm:mpug_embed_monomial, thm:mpug_czx_orbit_classification,
        thm:mpug_czx_placed_operators}
    The average over $\Z_2=\{e,g\}$ is $\frac12(\mathcal G_e+\mathcal G_g)$ with
    $\mathcal G_e=\Id$, and relabelling and placement are linear and unital. A
    vector $v$ satisfies $\frac12(v+G^\star_jv)=v$ exactly when
    $G^\star_jv=v$. Lemma~\ref{thm:mpug_czx_star_local_monomial_action} and
    Lemma~\ref{thm:mpug_embed_monomial} express $G^\star_j$ as a monomial
    operator; the permutation part is that of
    Lemma~\ref{thm:mpug_czx_placed_operators}, and decoding the two-site
    window at $j$ identifies the phase with $i^{q^\star_j(s)}$. Finally
    Lemma~\ref{thm:mpug_czx_orbit_classification} reindexes both to the
    fiber coordinates.
\end{proof}

\begin{lemma}[Reduction of the dimension to a finite count]
    \label{thm:mpug_czx_star_reduction}
    \lean{MPOTensor.CZX.finrank_commonFixedSubmodule_placedGaussProjector_circuitTupleStar}
    \leanok
    \uses{def:mpug_placed_gauss_projector, def:mpug_czx_star_tuple,
        def:mpug_common_fixed_subspace}
    On a periodic chain of $N\geq3$ blocked sites,
    \begin{align}
        \dim\mathcal V_N(R^\star)
         & =\#\bigl\{(p,b)\in\mathbb F_2^N\times\mathbb F_2^N:
        c_j(p,b)=0\text{ for all }j\in\Z/N\Z\bigr\}.
        \label{eq:mpug_czx_star_count}
    \end{align}
\end{lemma}

\begin{proof}\leanok
    \uses{thm:mpug_czx_star_placed_operators,
        thm:mpug_monomial_fixed_space_criterion,
        thm:mpug_czx_star_holonomy, thm:mpug_czx_star_good_fibers}
    By Lemma~\ref{thm:mpug_czx_star_placed_operators},
    $\mathcal V_N(R^\star)$ is the common fixed subspace of monomial operators
    which, in fiber coordinates, have permutations $\gamma\mapsto\gamma+e_j$ on
    the fibers over $(p,b)$, and by (\ref{eq:mpug_czx_star_involution}) these
    operators are involutions. The monomial fixed-space criterion
    (\ref{eq:mpug_monomial_fixed_space_count}) counts the fibers with trivial
    holonomy, which by Lemma~\ref{thm:mpug_czx_star_good_fibers} are exactly
    the fibers whose labels satisfy the $N$ equations $c_j(p,b)=0$.
\end{proof}

\begin{lemma}[Dimensions for the second displayed tuple at three and four sites]
    \label{thm:mpug_czx_star_dimension}
    \lean{MPOTensor.CZX.finrank_commonFixedSubmodule_placedGaussProjector_circuitTupleStar_three,
        MPOTensor.CZX.finrank_commonFixedSubmodule_placedGaussProjector_circuitTupleStar_four}
    \leanok
    \uses{def:mpug_placed_gauss_projector, def:mpug_czx_star_tuple,
        def:mpug_common_fixed_subspace}
    On periodic chains of three and four blocked sites,
    \begin{align}
        \dim\mathcal V_3(R^\star) & =14,
        \label{eq:mpug_czx_star_dim_three} \\
        \dim\mathcal V_4(R^\star) & =36.
        \label{eq:mpug_czx_star_dim_four}
    \end{align}
    These two values are certificates at two system sizes. They carry no
    asymptotic content, and no growth law for $\dim\mathcal V_N(R^\star)$
    follows from them. Since $R^\star$ extends the four prescribed maps of
    Definition~\ref{def:mpug_czx_defect_maps} by
    Lemma~\ref{thm:mpug_czx_star_tuple_completion}, these are two values of
    the dimension on that completion class, recorded as lower bounds on its
    supremum in Lemma~\ref{thm:mpug_czx_star_class_values}; they give no
    value for the minimum or the maximum of the dimension over completions.
    Compared with
    (\ref{eq:mpug_czx_dimension}), which gives $8$ and $16$ at the same two
    sizes, they exhibit dependence of the dimension on the displayed circuit
    data at those two sizes.
\end{lemma}

\begin{proof}\leanok
    \uses{thm:mpug_czx_star_reduction}
    Evaluating the right-hand side of (\ref{eq:mpug_czx_star_count}) by
    exhaustive enumeration gives the two dimensions: at $N=3$ the $64$
    pairs $(p,b)$ yield $14$ solutions of the three equations
    $c_j(p,b)=0$, and at $N=4$ the $256$ pairs yield $36$ solutions of
    the four equations.
\end{proof}

\section{The CZX defect domains and prescribed defect maps}
\label{sec:mpug_czx_defect_maps}

The modified-fusion lemma of
\cite[lines~4215--4254]{FrancoRubio2025MPUGauging} attaches to each pair
$(a,b)$ of gauge labels and each injective block $x$ a two-site defect vector
$v(a,b,x)$ mapped unitarily to $v(e,ab,x)$. For the CZX model, the defect
subspaces and prescribed maps are evaluated by finite contraction in
\cite{gap:fbc25_czx_defect_domains}.

\begin{definition}[Single-site CZX defects]
    \label{def:mpug_czx_site_defects}
    \lean{MPOTensor.CZX.identityDefect,
        MPOTensor.CZX.flipDefect,
        MPOTensor.CZX.siteDefect,
        MPOTensor.CZX.identityDefect_zero,
        MPOTensor.CZX.identityDefect_one,
        MPOTensor.CZX.flipDefect_zero,
        MPOTensor.CZX.flipDefect_one,
        MPOTensor.CZX.siteDefect_one,
        MPOTensor.CZX.siteDefect_gen}
    \leanok
    \discussion{7739}
    One blocked site of the CZX model carries two matter qubits, and the two
    injective blocks of the once-blocked GHZ state are the product vectors
    $\ket{\circ}=\ket{00}$ and $\ket{\bullet}=\ket{11}$
    \cite[lines~4660--4670]{FrancoRubio2025MPUGauging}, both of bond dimension
    one. Writing $e\leftrightarrow0$ and $g\leftrightarrow1$ for the labels of
    $\Z_2=\{e,g\}$, the identity defect and the domain-wall defect of
    \cite[lines~2831--2845]{FrancoRubio2025MPUGauging} at the block
    $x\in\mathbb F_2$ are
    \begin{align}
        e[x] & =\ket{xx},   \label{eq:mpug_czx_identity_defect}  \\
        g[0] & =\ket{00},   \label{eq:mpug_czx_flip_defect_zero} \\
        g[1] & =-i\ket{11}. \label{eq:mpug_czx_flip_defect_one}
    \end{align}
    As carried out in~\cite{gap:fbc25_czx_defect_domains} and displayed
    in \cite[lines~4800--4859]{FrancoRubio2025MPUGauging}, contracting the
    action tensors $A^{\Gamma}_{g,0}=-i\bra{1}$ and
    $A^{\Gamma}_{g,1}=\bra{0}$ of
    \cite[lines~4690--4694]{FrancoRubio2025MPUGauging} against the defect
    diagram gives (\ref{eq:mpug_czx_flip_defect_zero}) and
    (\ref{eq:mpug_czx_flip_defect_one}).
\end{definition}

\begin{definition}[Two-site CZX defect vectors]
    \label{def:mpug_czx_defect_vectors}
    \lean{MPOTensor.CZX.defectVector,
        MPOTensor.CZX.defectVector_one_one,
        MPOTensor.CZX.defectVector_one_gen,
        MPOTensor.CZX.defectVector_gen_one,
        MPOTensor.CZX.defectVector_gen_gen}
    \leanok
    \uses{def:mpug_czx_site_defects}
    The two-site defect vector of the label pair $(a,b)$ and the block $x$ is
    the left defect $a[bx]$ tensored with the right defect $b[x]$, in that
    order \cite[lines~4234--4245]{FrancoRubio2025MPUGauging},
    \begin{align}
        v(a,b,x) & =a[bx]\otimes b[x].
        \label{eq:mpug_czx_defect_vector}
    \end{align}
    The four qubits of two neighboring blocked sites are numbered $0,1$ on the
    left site and $2,3$ on the right, and the computational basis of the
    sixteen-dimensional two-site matter space is written $\ket{x_0x_1x_2x_3}$.
\end{definition}

\begin{lemma}[The eight CZX defect vectors]
    \label{thm:mpug_czx_defect_vector_values}
    \lean{MPOTensor.CZX.defectVector_one_one_zero,
        MPOTensor.CZX.defectVector_one_one_one,
        MPOTensor.CZX.defectVector_one_gen_zero,
        MPOTensor.CZX.defectVector_one_gen_one,
        MPOTensor.CZX.defectVector_gen_one_zero,
        MPOTensor.CZX.defectVector_gen_one_one,
        MPOTensor.CZX.defectVector_gen_gen_zero,
        MPOTensor.CZX.defectVector_gen_gen_one,
        MPOTensor.CZX.star_defectVector_dotProduct_defectVector}
    \leanok
    \uses{def:mpug_czx_defect_vectors}
    The eight two-site defect vectors of the CZX model are
    \begin{align}
        v(e,e,0) & =\ket{0000},   & v(e,e,1) & =\ket{1111},
        \label{eq:mpug_czx_v_ee}                              \\
        v(e,g,0) & =\ket{1100},   & v(e,g,1) & =-i\ket{0011},
        \label{eq:mpug_czx_v_eg}                              \\
        v(g,e,0) & =\ket{0000},   & v(g,e,1) & =-i\ket{1111},
        \label{eq:mpug_czx_v_ge}                              \\
        v(g,g,0) & =-i\ket{1100}, & v(g,g,1) & =-i\ket{0011}.
        \label{eq:mpug_czx_v_gg}
    \end{align}
    For every label pair the two vectors $v(a,b,0)$ and $v(a,b,1)$ are
    orthonormal.
\end{lemma}

\begin{proof}\leanok
    \uses{def:mpug_czx_site_defects, def:mpug_czx_defect_vectors}
    Substituting (\ref{eq:mpug_czx_identity_defect}),
    (\ref{eq:mpug_czx_flip_defect_zero}), and
    (\ref{eq:mpug_czx_flip_defect_one}) into
    (\ref{eq:mpug_czx_defect_vector}), with $bx=x+b$ in $\mathbb F_2$, gives
    the eight values. The two vectors of a label pair differ in the last two
    bits, hence are orthogonal, and each carries a scalar of modulus one,
    hence is a unit vector.
\end{proof}

\begin{lemma}[Reindexed monomial operators on the computational basis]
    \label{lem:mpug_czx_matter_monomial_action}
    \lean{MPOTensor.CZX.matterKet,
        MPOTensor.CZX.matterMatrix_monomial_mulVec_matterKet}
    \leanok
    \uses{def:mpug_monomial_matrix, def:mpug_czx_circuit_tuple,
        def:mpug_czx_defect_vectors}
    Write $\ket{x}$ for the computational basis vector of the two-site matter
    space attached to $x\in\mathbb F_2^4$ by the qubit numbering of
    Definition~\ref{def:mpug_czx_defect_vectors}. After reindexing to that
    space a monomial four-qubit operator still acts on the computational basis
    by its phase and its permutation, so that for every $x\in\mathbb F_2^4$
    \begin{align}
        M(\sigma,\varphi)\ket{x} & =\varphi(x)\ket{\sigma(x)}.
        \label{eq:mpug_czx_matter_monomial_action}
    \end{align}
\end{lemma}

\begin{proof}\leanok
    \uses{thm:mpug_monomial_matrix_laws, def:mpug_czx_circuit_tuple}
    Reindexing is the simultaneous relabelling of rows and columns along the
    qubit numbering. It carries $M(\sigma,\varphi)$ to the monomial matrix of
    the relabelled permutation and phase, and carries the basis vector of $x$
    to $\ket{x}$, so reading (\ref{eq:mpug_monomial_matrix}) in the new labels
    gives (\ref{eq:mpug_czx_matter_monomial_action}).
\end{proof}

\begin{lemma}[Unmodified CZX fusion contractions]
    \label{thm:mpug_czx_unmodified_fusion}
    \lean{MPOTensor.CZX.matterMatrix_lambda_mulVec_defectVector_zero,
        MPOTensor.CZX.matterMatrix_lambda_mulVec_defectVector_one}
    \leanok
    \discussion{7355}
    \uses{def:mpug_czx_circuit_tuple, def:mpug_czx_defect_vectors}
    Let $\Z_2=\{e,g\}$, and let $v(a,b,x)$ be the two-site defect vectors
    of Definition~\ref{def:mpug_czx_defect_vectors}, where $x\in\{0,1\}$.
    As in the two contractions displayed in
    \cite[lines~4989--5067]{FrancoRubio2025MPUGauging}, on the four matter
    qubits numbered $0,1,2,3$, the unmodified operator (as distinct from
    $\tilde\lambda=\lambda Z_0$ used in the displayed circuit tuple)
    \begin{align}
        \lambda & =\lambda^R_{g,g}
        =-i\,(X_0X_1X_2Z_3) \mathrm{CZ}_{01}\mathrm{CZ}_{12}\,X_2
    \end{align}
    satisfies
    \begin{align}
        \lambda\,v(g,g,0) & =-v(e,e,0), \\
        \lambda\,v(g,g,1) & =v(e,e,1).
    \end{align}
\end{lemma}

\begin{proof}\leanok
    \uses{thm:mpug_czx_phase_tables, thm:mpug_czx_defect_vector_values,
        lem:mpug_czx_matter_monomial_action}
    The phase formula~(\ref{eq:mpug_czx_lambda_action}), read on the two-site
    matter space through (\ref{eq:mpug_czx_matter_monomial_action}), gives
    \begin{align}
        \lambda\ket{1100} & =-i\ket{0000},
        \label{eq:mpug_czx_lambda_ket_even} \\
        \lambda\ket{0011} & =i\ket{1111},
        \label{eq:mpug_czx_lambda_ket_odd}
    \end{align}
    since $h(1100)=0$ and $h(0011)=1$.
    By Lemma~\ref{thm:mpug_czx_defect_vector_values}, the input defect
    vectors are $-i\ket{1100}$ and $-i\ket{0011}$, while the target
    vectors are $\ket{0000}$ and $\ket{1111}$. Multiplying
    (\ref{eq:mpug_czx_lambda_ket_even}) and
    (\ref{eq:mpug_czx_lambda_ket_odd}) by $-i$ gives the claimed signs.
\end{proof}

\begin{definition}[CZX defect domains]
    \label{def:mpug_czx_defect_domains}
    \lean{MPOTensor.CZX.defectDomain,
        MPOTensor.CZX.evenTarget,
        MPOTensor.CZX.oddTarget}
    \leanok
    \uses{def:mpug_czx_defect_vectors}
    The defect domain of the label pair $(a,b)$ and the two target planes are
    \begin{align}
        \mathcal K_{a,b} & =\spn_{\C}\{v(a,b,x):x\in\mathbb F_2\},
        \label{eq:mpug_czx_domain}                                 \\
        \mathcal L_0     & =\spn_{\C}\{\ket{0000},\ket{1111}\},
        \label{eq:mpug_czx_even_target}                            \\
        \mathcal L_1     & =\spn_{\C}\{\ket{1100},\ket{0011}\}.
        \label{eq:mpug_czx_odd_target}
    \end{align}
\end{definition}

\begin{lemma}[The four CZX defect domains]
    \label{thm:mpug_czx_defect_domains}
    \lean{MPOTensor.CZX.defectDomain_one_one,
        MPOTensor.CZX.defectDomain_gen_one,
        MPOTensor.CZX.defectDomain_one_gen,
        MPOTensor.CZX.defectDomain_gen_gen,
        MPOTensor.CZX.defectDomain_eq,
        MPOTensor.CZX.star_dotProduct_eq_zero_of_mem_evenTarget_of_mem_oddTarget,
        MPOTensor.CZX.finrank_evenTarget,
        MPOTensor.CZX.finrank_oddTarget,
        MPOTensor.CZX.finrank_defectDomain}
    \leanok
    \uses{def:mpug_czx_defect_domains}
    The four CZX defect domains are
    \begin{align}
        \mathcal K_{e,e}=\mathcal K_{g,e} & =\mathcal L_0,
        \label{eq:mpug_czx_even_domains}                   \\
        \mathcal K_{e,g}=\mathcal K_{g,g} & =\mathcal L_1,
        \label{eq:mpug_czx_odd_domains}
    \end{align}
    so the domain depends only on the right label. The two planes are
    orthogonal, $\mathcal L_0\perp\mathcal L_1$, and each has dimension two.
\end{lemma}

\begin{proof}\leanok
    \uses{thm:mpug_czx_defect_vector_values, def:mpug_czx_defect_domains}
    By (\ref{eq:mpug_czx_v_ee})--(\ref{eq:mpug_czx_v_gg}) the two spanning
    vectors of $\mathcal K_{a,b}$ are the two spanning vectors of
    $\mathcal L_b$ up to nonzero scalars, and rescaling a spanning vector by a
    nonzero scalar leaves the span unchanged, which gives
    (\ref{eq:mpug_czx_even_domains}) and (\ref{eq:mpug_czx_odd_domains}). The
    four basis vectors occurring in (\ref{eq:mpug_czx_even_target}) and
    (\ref{eq:mpug_czx_odd_target}) are pairwise distinct members of the
    computational basis, so expanding a vector of $\mathcal L_0$ and a vector
    of $\mathcal L_1$ in that basis makes every one of the four cross terms of
    their inner product vanish, and the two spanning vectors of each plane are
    linearly independent.
\end{proof}

\begin{definition}[The prescribed CZX defect maps]
    \label{def:mpug_czx_prescribed_maps}
    \lean{MPOTensor.CZX.prescribedMap}
    \leanok
    \uses{def:mpug_czx_defect_vectors, def:mpug_czx_defect_domains}
    The prescribed defect map of the label pair $(a,b)$ is the partial
    isometry
    \begin{align}
        D_{a,b} & =\sum_{x\in\mathbb F_2}\ketbra{v(e,ab,x)}{v(a,b,x)}
        \label{eq:mpug_czx_prescribed_map}
    \end{align}
    of the two-site matter space. It is formed from the defect vectors
    (\ref{eq:mpug_czx_v_ee})--(\ref{eq:mpug_czx_v_gg}) and not from the
    circuits $w$ and $\tilde\lambda$ of
    Definition~\ref{def:mpug_czx_circuit_tuple}.
\end{definition}

\begin{lemma}[Action of the prescribed CZX defect maps]
    \label{thm:mpug_czx_prescribed_action}
    \lean{MPOTensor.CZX.prescribedMap_mulVec_defectVector,
        MPOTensor.CZX.prescribedMap_mem_defectDomain,
        MPOTensor.CZX.isometry_on_defectDomain,
        MPOTensor.CZX.prescribedMap_one_one_mulVec,
        MPOTensor.CZX.prescribedMap_one_gen_mulVec,
        MPOTensor.CZX.prescribedMap_gen_one_mulVec_zero,
        MPOTensor.CZX.prescribedMap_gen_one_mulVec_one,
        MPOTensor.CZX.prescribedMap_gen_gen_mulVec_zero,
        MPOTensor.CZX.prescribedMap_gen_gen_mulVec_one}
    \leanok
    \uses{def:mpug_czx_prescribed_maps}
    For every label pair and every block,
    \begin{align}
        D_{a,b}v(a,b,x) & =v(e,ab,x),
        \label{eq:mpug_czx_prescription}
    \end{align}
    the map $D_{a,b}$ carries $\mathcal K_{a,b}$ into $\mathcal K_{e,ab}$, and
    it is an isometry on $\mathcal K_{a,b}$, that is,
    $D_{a,b}^{\dagger}D_{a,b}\ket{\xi}=\ket{\xi}$ for every
    $\ket{\xi}\in\mathcal K_{a,b}$. In the computational basis,
    \begin{align}
        D_{e,e} & =\Id_{\mathcal L_0},
        \label{eq:mpug_czx_map_ee}                     \\
        D_{e,g} & =\Id_{\mathcal L_1},
        \label{eq:mpug_czx_map_eg}                     \\
        D_{g,e}\ket{0000}=\ket{1100},
                & \quad D_{g,e}\ket{1111}=\ket{0011},
        \label{eq:mpug_czx_map_ge}                     \\
        D_{g,g}\ket{1100}=i\ket{0000},
                & \quad D_{g,g}\ket{0011}=i\ket{1111}.
        \label{eq:mpug_czx_map_gg}
    \end{align}
\end{lemma}

\begin{proof}\leanok
    \uses{thm:mpug_czx_defect_vector_values, def:mpug_czx_prescribed_maps}
    Applying (\ref{eq:mpug_czx_prescribed_map}) to $v(a,b,x)$ and using the
    orthonormality of $v(a,b,0)$ and $v(a,b,1)$ of
    Lemma~\ref{thm:mpug_czx_defect_vector_values} leaves the single term
    $v(e,ab,x)$, which is (\ref{eq:mpug_czx_prescription}). The same
    computation for the adjoint gives
    $D_{a,b}^{\dagger}v(e,ab,x)=v(a,b,x)$, and composing the two identities on
    a linear combination of $v(a,b,0)$ and $v(a,b,1)$ gives the isometry
    property and the inclusion of the images. Substituting
    (\ref{eq:mpug_czx_v_ee})--(\ref{eq:mpug_czx_v_gg}) into
    (\ref{eq:mpug_czx_prescription}) and cancelling the scalars of modulus one
    gives (\ref{eq:mpug_czx_map_ee})--(\ref{eq:mpug_czx_map_gg}).
\end{proof}

\begin{definition}[Prescribed defect maps of the CZX model]
    \label{def:mpug_czx_defect_maps}
    \lean{MPOTensor.CZX.defectMaps,
        MPOTensor.CZX.defectMaps_domain,
        MPOTensor.CZX.defectMaps_prescribed}
    \leanok
    \uses{def:mpug_defect_maps, def:mpug_czx_defect_domains,
        def:mpug_czx_prescribed_maps, thm:mpug_czx_prescribed_action}
    The four domains (\ref{eq:mpug_czx_even_domains}) and
    (\ref{eq:mpug_czx_odd_domains}) together with the four maps
    (\ref{eq:mpug_czx_prescribed_map}) are prescribed defect maps in the sense
    of Definition~\ref{def:mpug_defect_maps} for the group $\Z_2$ and the
    two-site matter space of the CZX model, the isometry condition
    (\ref{eq:mpug_defect_isometry}) being
    Lemma~\ref{thm:mpug_czx_prescribed_action}. That the circuit tuple
    $R^{\mathrm{CZX}}$ lies in the resulting completion class
    (\ref{eq:mpug_completion_class}) is proved unconditionally in
    Lemma~\ref{thm:mpug_czx_tuple_completion} below.
\end{definition}

\section{The displayed CZX tuple as a completion}
\label{sec:mpug_czx_completion}

The four defect domains and four prescribed maps of
Section~\ref{sec:mpug_czx_defect_maps} constrain each fusion operator on one
plane of the two-site matter space. The operators
$(\Id,\Id,w,\tilde\lambda)$ of Definition~\ref{def:mpug_czx_circuit_tuple}
satisfy these constraints, so the displayed tuple belongs to the completion
class of the CZX prescribed defect maps $D^{\mathrm{CZX}}$, and
(\ref{eq:mpug_czx_dimension}) lower-bounds the supremum of the gauge-invariant
dimension over the full class.

\begin{lemma}[The displayed CZX tuple is a completion]
    \label{thm:mpug_czx_tuple_completion}
    \lean{MPOTensor.CZX.circuitTuple_one_mulVec_defectVector,
        MPOTensor.CZX.circuitTuple_gen_one_mulVec_defectVector_zero,
        MPOTensor.CZX.circuitTuple_gen_one_mulVec_defectVector_one,
        MPOTensor.CZX.circuitTuple_gen_gen_mulVec_defectVector_zero,
        MPOTensor.CZX.circuitTuple_gen_gen_mulVec_defectVector_one,
        MPOTensor.CZX.circuitTuple_mulVec_defectVector,
        MPOTensor.CZX.isCompletion_circuitTuple,
        MPOTensor.CZX.circuitTuple_mem_completionClass}
    \leanok
    \uses{def:mpug_czx_defect_maps, def:mpug_czx_circuit_tuple,
        def:mpug_completion_class}
    For every ordered pair $(a,b)$ of labels of $\Z_2$ and every injective
    block $x$, the displayed operator of the pair carries the defect vector of
    the pair to the defect vector of the product label,
    \begin{align}
        U^{\mathrm{CZX}}_{a,b}v(a,b,x) & =v(e,ab,x),
        \label{eq:mpug_czx_tuple_fusion}
    \end{align}
    that is,
    \begin{align}
        \Id\,v(e,b,x)           & =v(e,b,x),
        \label{eq:mpug_czx_tuple_fusion_trivial}  \\
        w\,v(g,e,x)             & =v(e,g,x),
        \label{eq:mpug_czx_tuple_fusion_movement} \\
        \tilde\lambda\,v(g,g,x) & =v(e,e,x).
        \label{eq:mpug_czx_tuple_fusion_modified}
    \end{align}
    Consequently the displayed operator of every pair agrees on
    $\mathcal K_{a,b}$ with the prescribed map $D_{a,b}$, and the displayed
    tuple lies in the full completion class,
    \begin{align}
        R^{\mathrm{CZX}} & \in\mathfrak C(D^{\mathrm{CZX}}).
        \label{eq:mpug_czx_tuple_completion}
    \end{align}
\end{lemma}

\begin{proof}\leanok
    \uses{thm:mpug_czx_phase_tables, thm:mpug_czx_defect_vector_values,
        thm:mpug_czx_defect_domains, thm:mpug_czx_prescribed_action}
    The phase tables (\ref{eq:mpug_czx_w_action}) and
    (\ref{eq:mpug_czx_tilde_lambda_action}), evaluated at the four
    computational basis states carrying the defect vectors of
    Lemma~\ref{thm:mpug_czx_defect_vector_values}, give
    \begin{align}
        w\ket{0000}             & =\ket{1100},
                                & w\ket{1111}
                                & =\ket{0011},
        \label{eq:mpug_czx_w_basis_action}                \\
        \tilde\lambda\ket{1100} & =i\ket{0000},
                                & \tilde\lambda\ket{0011}
                                & =i\ket{1111},
        \label{eq:mpug_czx_tilde_lambda_basis_action}
    \end{align}
    because $e(0000)=e(1111)=0$ and $f(1100)=f(0011)=1$. Substituting
    (\ref{eq:mpug_czx_v_ee})--(\ref{eq:mpug_czx_v_gg}) and cancelling the
    factors $-i$ gives (\ref{eq:mpug_czx_tuple_fusion_movement}) and
    (\ref{eq:mpug_czx_tuple_fusion_modified}); the pairs with left label $e$
    carry the identity operator, so
    (\ref{eq:mpug_czx_tuple_fusion_trivial}) is immediate. The two defect
    vectors of a pair span $\mathcal K_{a,b}$ by (\ref{eq:mpug_czx_domain}),
    and the prescribed map obeys the same identity
    (\ref{eq:mpug_czx_prescription}), so the displayed operator and the
    prescribed map agree on the whole plane, which is
    (\ref{eq:mpug_czx_tuple_completion}).
\end{proof}

\begin{lemma}[Lower bound on the supremum over the CZX completion class]
    \label{thm:mpug_czx_class_maximum}
    \lean{MPOTensor.CZX.bddAbove_range_finrank_gaugeInvariantSubspace,
        MPOTensor.CZX.two_pow_le_iSup_finrank_gaugeInvariantSubspace_completionClass}
    \leanok
    \uses{def:mpug_czx_defect_maps, def:mpug_czx_circuit_tuple,
        def:mpug_completion_class, def:mpug_gauge_invariant_subspace,
        thm:mpug_czx_dimension_of_completion}
    On a periodic chain of $N\geq3$ blocked sites the gauge-invariant
    dimensions of the members of the CZX completion class are bounded above by
    the dimension of the chain space, so their supremum is finite, and
    \begin{align}
        \sup_{R\in\mathfrak C(D^{\mathrm{CZX}})}\dim\mathcal V_N(R) & \geq2^N.
        \label{eq:mpug_czx_class_maximum}
    \end{align}
    Nothing is asserted about the minimum of $\dim\mathcal V_N(R)$ over the
    class, about the exact value of the supremum, or about which completion
    attains it. With the domains and maps entering the class derived
    in~\cite{gap:fbc25_czx_defect_domains}, those extrema are questions about
    this completion class, not questions the source raises;
    \cite[lines~5198--5204]{FrancoRubio2025MPUGauging} asks only whether the
    gauge-invariant subspace of the state-level gauging stays bounded or grows
    with the system size, and nothing is asserted about that growth either.
\end{lemma}

\begin{proof}\leanok
    \uses{thm:mpug_czx_tuple_completion, thm:mpug_czx_dimension_of_completion}
    Every gauge-invariant subspace (\ref{eq:mpug_gauge_invariant_subspace}) is
    a subspace of the chain space, so the dimensions of the members of the
    class are bounded above by the dimension of that space and their supremum
    is finite. By Lemma~\ref{thm:mpug_czx_tuple_completion} the displayed
    tuple lies in the class, and by
    Lemma~\ref{thm:mpug_czx_dimension_of_completion} its gauge-invariant
    subspace has dimension $2^N$, which is therefore at most the supremum.
\end{proof}

\begin{lemma}[The second displayed CZX tuple is a completion]
    \label{thm:mpug_czx_star_tuple_completion}
    \lean{MPOTensor.CZX.matterMatrix_tildeLambdaStar_mulVec_matterKet_zero,
        MPOTensor.CZX.matterMatrix_tildeLambdaStar_mulVec_matterKet_one,
        MPOTensor.CZX.circuitTupleStar_gen_gen_mulVec_defectVector_zero,
        MPOTensor.CZX.circuitTupleStar_gen_gen_mulVec_defectVector_one,
        MPOTensor.CZX.circuitTupleStar_mulVec_defectVector,
        MPOTensor.CZX.isCompletion_circuitTupleStar,
        MPOTensor.CZX.circuitTupleStar_mem_completionClass}
    \leanok
    \uses{def:mpug_czx_defect_maps, def:mpug_czx_star_tuple,
        def:mpug_completion_class}
    The second displayed tuple obeys the same prescription
    (\ref{eq:mpug_czx_tuple_fusion}) as the first,
    \begin{align}
        U^\star_{a,b}v(a,b,x) & =v(e,ab,x),
        \label{eq:mpug_czx_star_tuple_fusion}
    \end{align}
    and therefore lies in the full completion class,
    \begin{align}
        R^\star & \in\mathfrak C(D^{\mathrm{CZX}}).
        \label{eq:mpug_czx_star_tuple_completion}
    \end{align}
\end{lemma}

\begin{proof}\leanok
    \uses{thm:mpug_czx_tuple_completion, thm:mpug_czx_defect_domains,
        def:mpug_czx_star_tuple}
    The pairs $(e,e)$, $(e,g)$, and $(g,e)$ carry the same operators in
    $R^\star$ as in $R^{\mathrm{CZX}}$, so
    Lemma~\ref{thm:mpug_czx_tuple_completion} gives
    (\ref{eq:mpug_czx_star_tuple_fusion}) there. For the pair $(g,g)$ the
    operator is $\tilde\lambda R_\star$, and the defect domain is
    $\mathcal K_{g,g}=\mathcal L_1=\spn_{\C}\{\ket{1100},\ket{0011}\}$ by
    (\ref{eq:mpug_czx_odd_domains}). Since
    $r_\star(1100)=r_\star(0011)=0$, the factor $R_\star$ fixes both spanning
    vectors, so $\tilde\lambda R_\star$ has the same values there as
    $\tilde\lambda$, namely $i\ket{0000}$ and $i\ket{1111}$. The two defect
    vectors of the pair span its domain, which gives
    (\ref{eq:mpug_czx_star_tuple_completion}).
\end{proof}

\begin{lemma}[Two further values of the dimension on the CZX completion
        class]
    \label{thm:mpug_czx_star_class_values}
    \lean{MPOTensor.CZX.finrank_gaugeInvariantSubspace_circuitTupleStar_three,
        MPOTensor.CZX.finrank_gaugeInvariantSubspace_circuitTupleStar_four,
        MPOTensor.CZX.fourteen_le_iSup_finrank_gaugeInvariantSubspace_completionClass,
        MPOTensor.CZX.thirtySix_le_iSup_finrank_gaugeInvariantSubspace_completionClass}
    \leanok
    \uses{def:mpug_czx_defect_maps, def:mpug_completion_class,
        def:mpug_gauge_invariant_subspace, thm:mpug_czx_star_tuple_completion,
        thm:mpug_czx_star_dimension}
    The dimensions $14$ at three sites and $36$ at four sites are values of the
    gauge-invariant dimension on the CZX completion class, so
    \begin{align}
        \sup_{R\in\mathfrak C(D^{\mathrm{CZX}})}\dim\mathcal V_3(R) & \geq14,
        \label{eq:mpug_czx_star_class_three}                                  \\
        \sup_{R\in\mathfrak C(D^{\mathrm{CZX}})}\dim\mathcal V_4(R) & \geq36.
        \label{eq:mpug_czx_star_class_four}
    \end{align}
    Beside the values $8$ and $16$ of the first displayed tuple at the same two
    chain lengths, these exhibit two different values of the dimension on the
    class at each of those lengths. Nothing follows about the minimum, the
    exact supremum, which completion attains it, or the behaviour of any of
    these as $N$ grows.
\end{lemma}

\begin{proof}\leanok
    \uses{thm:mpug_czx_star_tuple_completion, thm:mpug_czx_star_dimension}
    By Lemma~\ref{thm:mpug_czx_star_tuple_completion} the second displayed
    tuple lies in the class, and by
    Lemma~\ref{thm:mpug_czx_star_dimension} its gauge-invariant subspace has
    dimensions $14$ and $36$ at three and four sites. Each value is therefore
    at most the supremum over the class at that chain length, by the same
    boundedness argument as in
    Lemma~\ref{thm:mpug_czx_class_maximum}.
\end{proof}

\section{The gauged CZX state and the modified Gauss law}
\label{sec:mpug_czx_gauged_state}

The state-level gauging of \cite[lines~5121--5183]{FrancoRubio2025MPUGauging}
promotes the symmetry defects of the once-blocked GHZ state to gauge degrees of
freedom. Each site of the gauged chain carries one blocked matter index in
$\{0,1,2,3\}$ and one gauge label in $\Z_2=\{e,g\}$. The local Gauss operator
of the unmodified fusion operators does not fix the gauged state. After the
modification $\lambda^R_{g,g}\mapsto\lambda^R_{g,g}Z_1$ the gauged state is
fixed, at the cost of neighbouring Gauss projectors that do not commute. The
modified operators $\widetilde{\mathcal G}_h$ form a representation of $\Z_2$
by Lemma~\ref{thm:mpug_gauss_representation}, as every local Gauss operator
does. The source does not determine the common fixed subspace of the modified
projectors beyond the gauged state; its dimension for the displayed tuple is
Lemma~\ref{thm:mpug_czx_dimension}.

\begin{definition}[The gauged CZX tensor]
    \label{def:mpug_czx_gauged_tensor}
    \lean{MPOTensor.CZX.gaugeLabelMatrix,
        MPOTensor.CZX.gaugedSite,
        MPOTensor.CZX.gaugedTensor,
        MPOTensor.CZX.gaugedMPS,
        MPOTensor.CZX.gaugedDefectState}
    \leanok
    \discussion{7356}
    \uses{def:mpug_czx_blocked_ghz, def:mpv}
    Let $\widehat A$ be the once-blocked GHZ tensor of
    Definition~\ref{def:mpug_czx_blocked_ghz}, and let
    \begin{align}
        G_e & =\Id,                                   &
        G_g & =\begin{pmatrix}0&1\\-i&0\end{pmatrix}.
        \label{eq:mpug_czx_gauge_matrices}
    \end{align}
    The gauged tensor of \cite[lines~5129--5178]{FrancoRubio2025MPUGauging}
    has, at a site with blocked matter index $m$ and gauge label $a$, the
    letter $G_a\widehat A^m$: the gauge matrix sits on the bond to the left of
    the site. Its periodic matrix product vector on $N$ sites is the gauged
    CZX state $\ket{\Phi_N}$. For a configuration $\alpha$ of gauge labels, let
    $\ket{\Psi_\alpha}$ be the periodic GHZ state with $G_{\alpha_k}$ inserted
    on the bond to the left of the site $k$.
\end{definition}

\begin{lemma}[The gauged CZX state as a product of nearest-neighbour factors]
    \label{thm:mpug_czx_gauged_state_product}
    \lean{MPOTensor.CZX.gaugedMPS_eq_gaugedVector,
        MPOTensor.CZX.gaugedLeftBond,
        MPOTensor.CZX.gaugedBondWeightZ,
        MPOTensor.CZX.gaugedBondFactorZ,
        MPOTensor.CZX.gaugedTensor_apply,
        MPOTensor.CZX.gaugedMPS_eq_prod,
        MPOTensor.CZX.gaugedMPS_replaceWindow}
    \leanok
    \discussion{7356}
    \uses{def:mpug_czx_gauged_tensor, def:mpug_gauged_vector}
    The gauged CZX state is the gauged vector
    $\sum_\alpha\ket{\Psi_\alpha}\otimes\ket{\alpha}$ of
    Definition~\ref{def:mpug_gauged_vector}, the form of
    \cite[lines~3511--3530]{FrancoRubio2025MPUGauging}. For a site $u=(m,a)$
    with $m=2y+y'$, identify $e,g$ with $0,1$, let $\ell(u)=y+a$ modulo two,
    and for $r\in\{0,1\}$ let
    \begin{align}
        c_u(r) & =\delta_{m,3r}\cdot
        \begin{cases}
            -i & \text{if } a=g \text{ and } r=0, \\
            1  & \text{otherwise}.
        \end{cases}
    \end{align}
    Then every letter satisfies
    $(G_a\widehat A^m)_{lr}=\delta_{l,\ell(u)}c_u(r)$, and
    \begin{align}
        \braket{s|\Phi_N}
         & =\prod_{n=1}^Nc_{s_n}\bigl(\ell(s_{n+1})\bigr),
        \label{eq:mpug_czx_gauged_product}
    \end{align}
    with the indices read cyclically. Replacing the sites $j$ and $j+1$
    without changing $\ell$ at $j$ changes only the factors at $n=j$ and
    $n=j+1$.
\end{lemma}

\begin{proof}\leanok
    \uses{def:mpug_czx_gauged_tensor}
    The letter $\widehat A^m$ is the diagonal projector onto the GHZ sector of
    $m$, and $G_a$ shifts the bond by $a$ with the weights
    (\ref{eq:mpug_czx_gauge_matrices}); comparing the entries of the eight
    letters gives the entry formula. Since the incoming bond of every letter is forced by the
    letter, the trace of a product of letters is the product of the entries
    along the forced bonds. Replacing the window changes the bond forced at
    $j+1$, which enters only the factors at $j$ and $j+1$.
\end{proof}

\begin{definition}[The unmodified CZX fusion operators]
    \label{def:mpug_czx_unmodified_tuple}
    \lean{MPOTensor.CZX.lambda_mem_unitaryGroup,
        MPOTensor.CZX.unmodifiedCircuitTuple}
    \leanok
    \discussion{7356}
    \uses{def:mpug_czx_circuit_tuple, def:mpug_gauss_operator}
    The unmodified fusion operators are
    $(\lambda^R_{e,e},\lambda^R_{e,g},\lambda^R_{g,e},\lambda^R_{g,g})
        =(\Id,\Id,w,\lambda)$, with $w$ and $\lambda$ the unitaries
    (\ref{eq:mpug_czx_w}) and (\ref{eq:mpug_czx_lambda}). Their local Gauss
    operator of the generator is
    \begin{align}
        \mathcal G_g
         & =\lambda^\dagger\otimes\ketbra{g,g}{e,e}
        +w^\dagger\otimes\ketbra{g,e}{e,g}
        +w\otimes\ketbra{e,g}{g,e}
        +\lambda\otimes\ketbra{e,e}{g,g},
        \label{eq:mpug_czx_unmodified_gauss}
    \end{align}
    the operator of \cite[lines~5121--5128]{FrancoRubio2025MPUGauging}.
\end{definition}

\begin{lemma}[The unmodified Gauss law does not fix the gauged state]
    \label{thm:mpug_czx_unmodified_not_invariant}
    \lean{MPOTensor.CZX.unmodifiedPlacedGaussOperator_mulVec_gaugedMPS_ne}
    \leanok
    \discussion{7356}
    \uses{def:mpug_czx_unmodified_tuple, def:mpug_placed_gauss_operator,
        thm:mpug_czx_gauged_state_product}
    For every $N\geq2$ and every window $(j,j+1)$, the unmodified Gauss
    operator (\ref{eq:mpug_czx_unmodified_gauss}) placed on the window does not
    fix $\ket{\Phi_N}$, as stated in
    \cite[line~5179]{FrancoRubio2025MPUGauging}.
\end{lemma}

\begin{proof}\leanok
    \uses{thm:mpug_czx_gauged_state_product, thm:mpug_embed_monomial}
    The unmodified fusion operators are monomial, so the placed operator is
    monomial and fixes a vector $v$ exactly when $v(\sigma s)=\varphi(s)v(s)$
    for every configuration $s$. Take $s$ with matter $\ket{00}$ and label $e$
    at every site, so that $\braket{s|\Phi_N}=1$. The operator sends $s$ to the
    configuration with $(\ket{11},g)$ at $j$ and $(\ket{00},g)$ at $j+1$, with
    phase $i$. By (\ref{eq:mpug_czx_gauged_product}) the gauged state has
    coefficient $-i$ there, and $-i\neq i$.
\end{proof}

\begin{lemma}[The modified Gauss law fixes the gauged state]
    \label{thm:mpug_czx_modified_invariant}
    \lean{MPOTensor.CZX.iPowZ,
        MPOTensor.CZX.toComplex_iPowZ,
        MPOTensor.CZX.gaussWindowImage,
        MPOTensor.CZX.placedGaussOperator_mulVec_gaugedMPS,
        MPOTensor.CZX.gaugedMPS_mem_gaugeInvariantSubspace}
    \leanok
    \discussion{7356}
    \uses{def:mpug_czx_circuit_tuple, thm:mpug_czx_placed_operators,
        thm:mpug_czx_gauged_state_product, def:mpug_gauge_invariant_subspace}
    Replace $\lambda^R_{g,g}$ by $\tilde\lambda^R_{g,g}=\lambda^R_{g,g}Z_1$ and
    keep the other three fusion operators, as in
    \cite[lines~5179--5183]{FrancoRubio2025MPUGauging}. For every $N\geq2$,
    every placed modified Gauss operator $G_j$ fixes $\ket{\Phi_N}$, and hence
    so does every placed modified Gauss projector
    $\widetilde P_j=\frac12(\Id+G_j)$. Thus $\ket{\Phi_N}$ lies in the
    gauge-invariant subspace $\mathcal V_N(R^{\mathrm{CZX}})$.
\end{lemma}

\begin{proof}\leanok
    \uses{thm:mpug_czx_local_monomial_action, thm:mpug_czx_placed_operators,
        thm:mpug_czx_gauged_state_product}
    By Lemma~\ref{thm:mpug_czx_placed_operators} the operator $G_j$ is
    monomial: it complements the matter bits at $j$, moves the two labels of
    the window by the Gauss leg action, and multiplies by $i^{q(s)}$. It does
    not change $\ell$ at the site $j$. By
    Lemma~\ref{thm:mpug_czx_gauged_state_product} the claim reduces to
    \begin{align}
        c_{u_0'}\bigl(\ell(u_1')\bigr)\,c_{u_1'}(b)
         & =i^{q(u_0,u_1)}\,c_{u_0}\bigl(\ell(u_1)\bigr)\,c_{u_1}(b)
    \end{align}
    for every window $(u_0,u_1)$ with image $(u_0',u_1')$ and every bond
    $b\in\{0,1\}$. All values are Gaussian integers, and the identity is
    checked over the $128$ cases.
\end{proof}

\begin{lemma}[Neighbouring modified projectors do not commute]
    \label{thm:mpug_czx_projectors_noncommuting}
    \lean{MPOTensor.CZX.not_commute_placedGaussOperator,
        MPOTensor.CZX.not_commute_placedGaussProjector}
    \leanok
    \discussion{7356}
    \uses{thm:mpug_czx_placed_operators, thm:mpug_czx_modified_invariant}
    For every $N\geq3$ and every $j$, the modified Gauss operators $G_j$ and
    $G_{j+1}$ do not commute, and neither do the projectors $\widetilde P_j$
    and $\widetilde P_{j+1}$. This is the cost stated in
    \cite[line~5183]{FrancoRubio2025MPUGauging}.
\end{lemma}

\begin{proof}\leanok
    \uses{thm:mpug_czx_placed_operators}
    Both products $G_jG_{j+1}$ and $G_{j+1}G_j$ are monomial. Let $s$ be the
    configuration with matter $\ket{00}$ and label $g$ at every site. Since
    $N\geq3$, the windows at $j$ and $j+1$ share only the site $j+1$, and the
    phases of the two products at $s$ are read off the phase table of
    Lemma~\ref{thm:mpug_czx_local_monomial_action}: they differ by a sign.
    The column of $s$ therefore differs in the two products, whether or not
    their permutations agree at $s$. Since $G_k=2\widetilde P_k-\Id$,
    commuting projectors would give commuting operators.
\end{proof}
